% Este fichero es parte del N�mero 4 NE de la Revista Online Occam's Razor
% Revista On-line Occam's Razor N�mero 4 Nueva �poca
%
% (c)  2019, The Occam's Razor Team
%
% Esta obra est� bajo una licencia Reconocimiento 3.0 Espa�a de
% Creative Commons. Para ver una copia de esta licencia, visite
% http://creativecommons.org/licenses/by/3.0/es/ o envie una carta a
% Creative Commons, 171 Second Street, Suite 300, San Francisco,
% California 94105, USA. 


\documentclass[10pt,a4paper,twoside]{article}


% Paquetes... probablemente alguno no sea necesario
% 

\usepackage[latin1]{inputenc}
%\usepackage[utf8]{inputenc}
\usepackage[greek,spanish]{babel}  % S�mbolo del euro
\usepackage{graphicx}
\usepackage{a4,fancyhdr, multicol}
\usepackage{float}
\usepackage{pdftricks}
%\usepackage[pdf]{pstricks}
\usepackage{color}
%\usepackage{pst-plot}
\usepackage{pst-eps}
\usepackage{wrapfig}
\usepackage{eso-pic}
\usepackage{listings}
\usepackage{textpos}
\usepackage{epsf}
\usepackage{setspace}
\usepackage{hyperref}
\usepackage{colortbl}
\usepackage{multirow}
\usepackage{mathtools,amssymb,amsthm} 

% Configuraci�n de Hyperenlaces.
% Uso del paquete hyperref sugerido por Cruz Enrique Borges Hern�ndez
\hypersetup{
    bookmarks=true,        % show bookmarks bar?
    colorlinks=true,       % false: boxed links; true: colored links
    linkcolor=red,         % color of internal links
    citecolor=green,       % color of links to bibliography
    filecolor=magenta,     % color of file links
    urlcolor=blue          % color of external links
}



%\usepackage[sfdefault,condensed]{roboto}
%\usepackage[sfdefault]{roboto}
\usepackage[default]{lato}
\renewcommand*\sfdefault{ugq}
\usepackage[T1]{fontenc} 




% Configuraci�n de tama�o de p�gina
\setlength{\parindent}{0in}
\setlength{\parskip}{0.1cm}
\setlength{\oddsidemargin}{0.05mm}
\setlength{\evensidemargin}{0.05mm}

\setlength{\headsep}{-0.6cm}
%\setlength{\footsep}{-1.0cm}
\setlength\headheight{2.1cm}

\addtolength{\textwidth}{4cm}
\addtolength{\topmargin}{-2.5cm}
\addtolength{\textheight}{3.4cm}

%\addtolength{\headwidth}{-1.2\marginparsep}
%\addtolength{\headwidth}{-1.2\marginparwidth}

\newcommand{\hsection}[4]{
  {\begin{flushright}{
        {\psset{linecolor=#1,linestyle=dotted}
          \psline(-14,0.05)(1,0.05)}
        \colorbox{#1}{
          \begin{minipage}{#3\linewidth}
            \center
            \textcolor{#2}{
              \textsf{\textbf{#4}}}            
          \end{minipage}
        }
        }
  \end{flushright}}
  \vspace{4mm}
}


\newcommand{\hhsection}[4]{
  {\begin{flushleft}{
        {\psset{linecolor=#1,linestyle=dotted}
          \psline(17,0.05)(1,0.05)}
        \colorbox{#1}{
          \begin{minipage}{#3\linewidth}
            \center
            \textcolor{#2}{
              \textsf{\textbf{#4}}}
          \end{minipage}
        }
        }
  \end{flushleft}}
  \vspace{4mm}
}





\fancypagestyle{divulgacion}{
    %\fancyhead{}
  \fancyhead{}
  \fancyfoot{} % clear all footer fields

  \fancyhead[RE]{\hsection{purple}{white}{0.14}{DIVULGACI�N}}
  \fancyhead[LO]{\hhsection{purple}{white}{0.14}{DIVULGACI�N}}

  \fancyfoot[LE]{\textbf{\textsf{OCCAM's Razor | \thepage}}}
  \fancyfoot[RO]{\textbf{\textsf{\thepage | OCCAM's Razor}}}
  \renewcommand{\footrulewidth}{0.4pt}
  
  \renewcommand{\headrulewidth}{0pt}
  }


\fancypagestyle{trucos}{
    %\fancyhead{}
  \fancyhead{}
  \fancyfoot{} % clear all footer fields

  \fancyhead[RE]{\hsection{lime}{black}{0.14}{TRUCOS}}
  \fancyhead[LO]{\hhsection{lime}{black}{0.14}{TRUCOS}}

  \fancyfoot[LE]{\textbf{\textsf{OCCAM's Razor | \thepage}}}
  \fancyfoot[RO]{\textbf{\textsf{\thepage | OCCAM's Razor}}}
  \renewcommand{\footrulewidth}{0.4pt}
  
  \renewcommand{\headrulewidth}{0pt}
  }






\fancypagestyle{ratas}{
    %\fancyhead{}
  \fancyhead{}
  \fancyfoot{} % clear all footer fields

  %\fancyhead[LE,RO]{\textbf{\textsf{RATAS}}}
  \fancyhead[RE]{\hsection{introcolor}{black}{0.2}{RATAS DE BIBLIOTECA}}
  \fancyhead[LO]{\hhsection{introcolor}{black}{0.2}{RATAS DE BIBLIOTECA}}
  \fancyfoot[LE]{\textbf{\textsf{OCCAM's Razor | \thepage}}}
  \fancyfoot[RO]{\textbf{\textsf{\thepage | OCCAM's Razor}}}
  \renewcommand{\footrulewidth}{0.4pt}
  
  \renewcommand{\headrulewidth}{0pt}
  }



\fancypagestyle{geek}{
    %\fancyhead{}
  \fancyhead{}
  \fancyfoot{} % clear all footer fields

  %\fancyhead[LE,RO]{\textbf{\textsf{RATAS}}}
  \fancyhead[RE]{\hsection{violet}{white}{0.2}{ZONA GEEK}}
  \fancyhead[LO]{\hhsection{violet}{white}{0.2}{ZONA GEEK}}
  \fancyfoot[LE]{\textbf{\textsf{OCCAM's Razor | \thepage}}}
  \fancyfoot[RO]{\textbf{\textsf{\thepage | OCCAM's Razor}}}
  \renewcommand{\footrulewidth}{0.4pt}
  
  \renewcommand{\headrulewidth}{0pt}
  }



\fancypagestyle{seguridad}{
    %\fancyhead{}
  \fancyhead{}
  \fancyfoot{} % clear all footer fields

  \fancyhead[RE]{\hsection{black}{green}{0.14}{SEGURIDAD}}
  \fancyhead[LO]{\hhsection{black}{green}{0.14}{SEGURIDAD}}

  \fancyfoot[LE]{\textbf{\textsf{OCCAM's Razor | \thepage}}}
  \fancyfoot[RO]{\textbf{\textsf{\thepage | OCCAM's Razor}}}
  \renewcommand{\footrulewidth}{0.4pt}
  
  \renewcommand{\headrulewidth}{0pt}
  }


\fancypagestyle{pvc}{
    %\fancyhead{}
  \fancyhead{}
  \fancyfoot{} % clear all footer fields

  %\fancyhead[LE,RO]{\textbf{\textsf{RATAS}}}
  \fancyhead[RE]{\hsection{orange}{black}{0.25}{VITAMINA C}}
  \fancyhead[LO]{\hhsection{orange}{black}{0.25}{VITAMINA C}}
  \fancyfoot[LE]{\textbf{\textsf{OCCAM's Razor | \thepage}}}
  \fancyfoot[RO]{\textbf{\textsf{\thepage | OCCAM's Razor}}}
  \renewcommand{\footrulewidth}{0.4pt}
  
  \renewcommand{\headrulewidth}{0pt}
  }


\fancypagestyle{como}{
    %\fancyhead{}
  \fancyhead{}
  \fancyfoot{} % clear all footer fields

  %\fancyhead[LE,RO]{\textbf{\textsf{RATAS}}}
  \fancyhead[RE]{\hsection{violet}{white}{0.2}{COMO FUNCIONA}}
  \fancyhead[LO]{\hhsection{violet}{white}{0.2}{COMO FUNCIONA}}

  \fancyfoot[LE]{\textbf{\textsf{OCCAM's Razor | \thepage}}}
  \fancyfoot[RO]{\textbf{\textsf{\thepage | OCCAM's Razor}}}
  \renewcommand{\footrulewidth}{0.4pt}
  
  \renewcommand{\headrulewidth}{0pt}
  }



\fancypagestyle{mala-bestia}{
      %\fancyhead{}
  \fancyhead{}
  \fancyfoot{} % clear all footer fields

  %\fancyhead[LE,RO]{\textbf{\textsf{RATAS}}}
    \fancyhead[RE]{\hsection{cyan}{black}{0.15}{MALA BESTIA}}
  \fancyhead[LO]{\hhsection{cyan}{black}{0.15}{MALA BESTIA}}

  \fancyfoot[LE]{\textbf{\textsf{OCCAM's Razor | \thepage}}}
  \fancyfoot[RO]{\textbf{\textsf{\thepage | OCCAM's Razor}}}
  \renewcommand{\footrulewidth}{0.4pt}
  
  \renewcommand{\headrulewidth}{0pt}
  }

\fancypagestyle{linux}{
    %\fancyhead{}
  \fancyhead{}
  \fancyfoot{} % clear all footer fields

  %\fancyhead[LE,RO]{\textbf{\textsf{RATAS}}}
  \fancyhead[RE]{\hsection{teal}{white}{0.25}{GNU/LINUX}}
  \fancyhead[LO]{\hhsection{teal}{white}{0.25}{GNU/LINUX}}
  \fancyfoot[LE]{\textbf{\textsf{OCCAM's Razor | \thepage}}}
  \fancyfoot[RO]{\textbf{\textsf{\thepage | OCCAM's Razor}}}
  \renewcommand{\footrulewidth}{0.4pt}
  
  \renewcommand{\headrulewidth}{0pt}
  }

\fancypagestyle{fpn}{
    %\fancyhead{}
  \fancyhead{}
  \fancyfoot{} % clear all footer fields

  \fancyhead[RE]{\hsection{pink}{black}{0.14}{FEN�MENOS EXTRA�OS}}
  \fancyhead[LO]{\hhsection{pink}{black}{0.14}{FEN�MENOS EXTRA�OS}}

  \fancyfoot[LE]{\textbf{\textsf{OCCAM's Razor | \thepage}}}
  \fancyfoot[RO]{\textbf{\textsf{\thepage | OCCAM's Razor}}}
  \renewcommand{\footrulewidth}{0.4pt}
  
  \renewcommand{\headrulewidth}{0pt}
  }


%\pagestyle{fancy}
% Configuraci�n de Cabeceras Fancy Header
%\fancyhead{}
%\fancyfoot{} % clear all footer fields
%\fancyfoot[LE]{\textbf{\textsf{OCCAM's Razor | \thepage}}}
%\fancyfoot[RO]{\textbf{\textsf{\thepage | OCCAM's Razor}}}
%\renewcommand{\footrulewidth}{0.4pt}

% Elimina l�neas de cabecera
% 
%\renewcommand{\headrulewidth}{0pt}

% Colores
\definecolor{introcolor}{rgb}{0.8,1.0,0.8}
\definecolor{listcolor}{rgb}{0.2,0.4,0.2}
\definecolor{titlecolor}{rgb}{0.4,0.5,0.1}
\definecolor{excolor}{rgb}{0.8,0.8,0.8}
\definecolor{notecolor}{rgb}{0.4,0.4,1.0}

% ********************************************************
% Definici�n de Comandos y Entornos
% ********************************************************

% Comandos de uso general
% ---------------------------------------------------------
% Secciones t�tulos y subt�tulos de cada p�gina
\newcommand{\msection}[4]{
{\begin{flushright}{
{\psset{linecolor=black,linestyle=dotted}\psline(-17,0)}
\colorbox{#1}{
\begin{minipage}{#3\linewidth}
\center
  \textcolor{#2}{
    \textsf{\textbf{#4}}}
\end{minipage}
}}\end{flushright}}

\vspace{4mm}
}

\newcommand{\mtitle}[2]{
  {\resizebox{#1}{0.7cm}{\textbf{\textsf{#2}}}}
  \vspace{1mm}
}


\newcommand{\msubtitle}[2]{
  {\resizebox{#1}{0.5cm}{{\gray{\textbf{\textsf{#2}}}}}}
  \vspace{1mm}
}

% Principio de P�gina. Pone el cuadro superior con la secci�n
\newcommand{\bOpage}[3]{
  \msection{#1}{black}{#2}{#3}
  \begin{multicols}{2}
}

% Fin de p�gina. Termina el entorno multicols
\newcommand{\eOpage}{
\pagebreak
\end{multicols}
}

% Fin e Inicio de P�gina. Sino utilizamos figuras fuera de las
% columnas del cuerpo principal, esta es la forma adecuada de marcar
% cada p�gina
\newcommand{\ebOpage}[3]{
\eOpage
\bOpage{#1}{#2}{#3}
}

% Crea el cuadro de introducci�n al principio de cada art�culo
\newcommand{\intro}[3]{
\colorbox{#1}{\large
  \begin{minipage}{.98\linewidth}
    \vspace{2mm}
    {\colorbox{black}{\color{white}{\resizebox{!}{1.0cm}{#2}}}{ #3}}
  \vspace{1mm}
  \end{minipage}
}
\vspace{4mm}
}



\newcommand{\introA}[3]{
\colorbox{#1}{
  \begin{minipage}{\linewidth}
    \vspace{2mm}
    {\colorbox{black}{\color{white}{\resizebox{!}{1.0cm}{#2}}}{#3}}
  \vspace{1mm}
  \end{minipage}
}
\vspace{4mm}
}



% Comando para introducir figura en entorno multicol
\newcommand{\myfig}[3]{
\begin{center}
  \includegraphics[width=#3\hsize,angle=#1]{#2}
  \nobreak
\end{center}}

% Caption para figuras en entorno multicol
\setcounter{figure}{1}
\newcommand{\mycaption}[1]{
  \begin{quote}
    {\small
    {{\sc Figura} \arabic{figure}: #1}
    }
  \end{quote}
  \stepcounter{figure}
}


% Comandos para utilizar tablas y figuras en entorno multicols
\makeatletter
\newenvironment{tablehere}
  {\def\@captype{table}}
  {}

\newenvironment{figurehere}
  {\def\@captype{figure}}
  {}
\makeatother

% Caption para figuras en entorno multicol sin contador
\newcommand{\nncaption}[1]{
  %\begin{quote}
    {\footnotesize{\textbf{
    {#1}
    }}}
  %\end{quote}
}

\newcommand{\sectiontext}[3]{\vspace{4mm}{{\textcolor{titlecolor}{\large{\textbf{\textsf{#3}}}}}}
\vspace{1mm}
}

%\newcommand{\sbullet}{\,\begin{picture}(1,1)(0,-3)\circle*{3}\end{picture}\ }
\newcommand{\EOP}{\psframe[fillstyle=solid, fillcolor=titlecolor,
linecolor=titlecolor](0,0)(4pt,4pt)}
%\newcommand{\EOP}{\psframe[fillstyle=solid, fillcolor=black, linecolor=black](0,0)(4pt,6pt)}

% Entornos (begin... end)
% ----------------------------------------------------
% Entorno para introducir ejemplos
\newenvironment{mexample}{
  \vspace{2mm}
  \bgroup
  \tiny
}{
  \egroup
  \vspace{4mm}
}

% Entorno para introducir entradillas en el texto
\newenvironment{entradilla}{
  \vspace{5mm}
  \color{introcolor}{\hrule }
  \color{black}
  \vspace{2mm}
  \bgroup
  \LARGE
\begin{spacing}{0.6}
}{
\end{spacing}
  \egroup
  \vspace{5mm}
  \color{introcolor}{\hrule height 3pt}
  \vspace{5mm}
}



% **************************************************************************
% Comienza el documento
\begin{document}

% Portada, no utiliza Fancy Header e introduce imagen de portada con PStricks
\pagestyle{empty}

\rput(8,-13){\epsfbox{images/portada.eps}}

\clearpage
\pagebreak

% Imagen con el sumario en la siguiente p�gina
%\rput(8,-14.0){\resizebox{!}{30cm}{{\epsfbox{images/sumario.eps}}}}



%\pagebreak

% Activa Fancy Headers stilo e incluye los distintos art�culos
\pagestyle{fancy}

\definecolor{introcolor}{rgb}{0.6,0.7,0.3}

\pagestyle{empty}

{\color{white}{.}}
%\vspace{10mm}

\rput(7.9,-25.4){\resizebox{21.0cm}{!}{{\epsfbox{images/general/tools.eps}}}}

\begin{flushright}

{\footnotesize{N�MERO 6 - MARZO 2024}}\\
{\textsf{Segunda �poca}} 


{\large{REVISTA ONLINE OCCAM'S RAZOR}}\\
{\color{introcolor}{\bf Porque lo m�s sencillo es lo m�s probable}}

\vspace{0.5cm}


{\resizebox{!}{1.8cm}{\textsf{SUMARIO}}}
\end{flushright}


\begin{tabular}{ll}
  \hline
  \multirow{4}{*}{
%    \colorbox{lightgray}{\resizebox{!}{1.3cm}{{0\pageref{ratas:aes}}}}
    \colorbox{lightgray}{\resizebox{!}{1.3cm}{\hyperlink{ratas:aes1}{\pageref*{ratas:aes}}}} 
  }
  &  \\
  %& \colorbox{introcolor}{{\large\textsf{RATAS DE BIBLIOTECA}}} \\
  & {\textsf{RATAS DE BIBLIOTECA}}\\
  %& \textsf{Calculando Hashes en tus Programas con OpenSSL} \\
  & \large\bf{Cifrado Militar en tus programas} \\
  %& \textbf{por Nelly Kerningham}   \vspace {1.0cm}\\
  & {\em por Don Bit0}   \vspace {0.5cm}\\
  



  \hline
  \multirow{4}{*}{
%    \colorbox{lightgray}{\resizebox{!}{1.3cm}{{\pageref{linux:kernel}}}}
    \colorbox{lightgray}{\resizebox{!}{1.3cm}{\hyperlink{linux:kernel1}{\pageref*{linux:kernel}}}}
    }
  &  \\
  %& {\large\textsf{GNU/LINUX}} \\
  & {\textsf{GNU/LINUX}} \\
  %& \colorbox{teal}{\color{white}{\large\textsf{GNU/LINUX}}} \\   
  %& \textsf{Manejando Volcados de Disco Como un Pro} \\
  & \large\bf{Multitarea Apropiativa en Ordenadores Dom�sticos} \\
  & {\em por Nelly Kerningham}   \vspace {0.5cm}\\


  \hline
  \multirow{4}{*}{
    %\colorbox{lightgray}{\resizebox{!}{1.3cm}{{\pageref{pildoras:alloca}}}}
    \colorbox{lightgray}{\resizebox{!}{1.3cm}{\hyperlink{pildoras:alloca1}{\pageref*{pildoras:alloca}}}}                    
  }
  &  \\
  %& {\large\textsf{SEGURIDAD}} \\
  & {\textsf{PROGRAMACI�N}} \\
  %& \colorbox{black}{\color{green}{\large\textsf{SEGURIDAD}}} \\
  %& \textsf{Seguridad para meros mortales. Parte II. VIRUS} \\
  & \large\bf{P�ldoras de Vitamina C: alloca} \\
  %& \textbf{por Don Bit0}   \vspace {1.0cm}\\
  & {\em por Richio C. Poweri}   \vspace {0.5cm}\\




  \hline
  \multirow{4}{*}{
    %\colorbox{lightgray}{\resizebox{!}{1.3cm}{{\pageref{sec:criptografia}}}}
    \colorbox{lightgray}{\resizebox{!}{1.3cm}{\hyperlink{sec:criptografia1}{\pageref*{sec:criptografia}}}}                
  }
  &  \\
  %& {\large\textsf{SEGURIDAD}} \\
  & {\textsf{SEGURIDAD}} \\
  %& \colorbox{black}{\color{green}{\large\textsf{SEGURIDAD}}} \\
  %& \textsf{Seguridad para meros mortales. Parte II. VIRUS} \\
  & \large\bf{Seguridad para meros mortales. Criptograf�a} \\
  %& \textbf{por Don Bit0}   \vspace {1.0cm}\\
  & {\em por Er Enigma}   \vspace {0.5cm}\\

  \hline
  \multirow{4}{*}{
    %\colorbox{lightgray}{\resizebox{!}{1.3cm}{{\pageref{sec:criptografia}}}}
    \colorbox{lightgray}{\resizebox{!}{1.3cm}{\hyperlink{fpn:bucle}{\pageref*{fpn:bucle}}}}                
  }
  &  \\
  %& {\large\textsf{SEGURIDAD}} \\
  & {\textsf{FEN�MENOS EXTRA�OS}} \\
  %& \colorbox{black}{\color{green}{\large\textsf{SEGURIDAD}}} \\
  %& \textsf{Seguridad para meros mortales. Parte II. VIRUS} \\
  & \large\bf{Bucle Infinito en C. Modo Esot�rico} \\
  %& \textbf{por Don Bit0}   \vspace {1.0cm}\\
  & {\em por Carolyn Lightrun}   \vspace {0.5cm}\\



  \hline
  \multirow{4}{*}{
%    \colorbox{lightgray}{\resizebox{!}{1.3cm}{{\pageref{divulgacion:compresion}}}}
    \colorbox{lightgray}{\resizebox{!}{1.3cm}{\hyperlink{divulgacion:compresion1}{\pageref*{divulgacion:compresion}}}}            
  }
  &  \\
  %& {\large\textsf{DIVULGACI�N}} \\
  & {\textsf{DIVULGACI�N}} \\
  %& \colorbox{purple}{\color{white}{\large\textsf{DIVULGACI�N}}} \\ 
  %& \textsf{Realidad Virtual. Realidad Aumentada. Realidad Mixta} \\
  & \large\bf{Compresi�n de datos sin p�rdidas} \\
  %& \textbf{por Signul Floid}   \vspace {1.0cm}\\
  & {\em por Signul Floid}   \vspace {0.5cm}\\


  \hline
  \multirow{4}{*}{
    %\colorbox{lightgray}{\resizebox{!}{1.5cm}{{\pageref{geek:bf}}}}
    \colorbox{lightgray}{\resizebox{!}{1.3cm}{\hyperlink{geek:bf1}{\pageref*{geek:bf}}}}      
  }
  &  \\
  %& {\large\textsf{DIVULGACI�N}} \\
  & {\textsf{ZONA GEEK}} \\
  %& \colorbox{purple}{\color{white}{\large\textsf{DIVULGACI�N}}} \\ 
  %& \textsf{Realidad Virtual. Realidad Aumentada. Realidad Mixta} \\
  & \large\bf{Interprete de BrainFuck} \\
  %& \textbf{por Signul Floid}   \vspace {1.0cm}\\
  & {\em por Friki Piki}   \vspace {0.5cm}\\



  \hline
  \multirow{4}{*}{
    %\colorbox{lightgray}{\resizebox{!}{1.5cm}{{\pageref*{trucos}}}}
    \colorbox{lightgray}{\resizebox{!}{1.3cm}{\hyperlink{trucos1}{\pageref*{trucos}}}}

  }
  &  \\
  %& {\large\textsf{TRUCOS}} \\
  & {\textsf{TRUCOS}} \\
  %& \colorbox{lime}{\color{black}{\large\textsf{TRUCOS}}} \\   
  %& \textsf{Brujer�as en la l�nea de comandos} \\
  & \large\bf{Con un por... de l�neas} \\
  & {\em por Tamariz el de la Perd�z}   \vspace {0.5cm}\\

\end{tabular}

\vspace{1.8cm}

\textsf{Esta revista ha sido realizada con:}

\clearpage

\pagebreak


% Este fichero es parte del N�mero 4 de la Revista Online Occam's Razor
% Revista Online Occam's Razor N�mero 4
%
% (c)  20019, The Occam's Razor Team
%
% Esta obra est� bajo una licencia Reconocimiento 3.0 Espa�a de
% Creative Commons. Para ver una copia de esta licencia, visite
% http://creativecommons.org/licenses/by/3.0/es/ o envie una carta a
% Creative Commons, 171 Second Street, Suite 300, San Francisco,
% California 94105, USA. 

\rput(6,-1.3){\resizebox{!}{6cm}{{\epsfbox{images/general/editorial-pluma.eps}}}}
\rput(0.0,-13){\resizebox{7cm}{35.0cm}{{\epsfbox{images/general/bar.eps}}}}
\rput(0.4,-4.3){\resizebox{!}{4.8cm}{{\epsfbox{images/portada-thumb.eps}}}}
\rput(0.4,-18.6){\resizebox{!}{1.7cm}{{\epsfbox{images/general/ibolcode.eps}}}}
\rput(0.7,-26){\resizebox{!}{0.9cm}{{\epsfbox{images/general/licencia.eps}}}}

\begin{textblock}{9.2}(3,0)
\begin{flushright}
{\resizebox{!}{1.3cm}{\textsf{Editorial}}}

\vspace{2mm}

{\LARGE Estamos de vuelta}

by The Occam Team
\end{flushright}
\end{textblock}

\vspace{4mm}


\definecolor{barcolor}{rgb}{0.9,0.9,0.9}

\begin{textblock}{30}(-1.5, -1)
\begin{minipage}{0.12\linewidth}
\sf\color{barcolor}
\begin{center}

\vspace{1cm}

\colorbox{black}{
{\resizebox{3cm}{0.7cm}{\textcolor{white}{\bf\sf\large Occam's}}}
}

{\resizebox{2.5cm}{0.4cm}{\bf\sf\large Razor}}

\vspace{4mm}

{\bf N�mero 6, A�o 2024}

\vspace{5cm}

\hrule

\vspace{4mm}

Direcci�n:

\vspace{1mm}

{\bf David Mart�nez Oliveira}

\vspace{2mm}



{Colaboradores:}

\vspace{1mm}
{\bf Nelly Kerningham, Er Enigma, Richi C. Poweri, Don Bit0, Signul Floid, Carolyn LightRun, Tamariz el de la Perdiz}



\vspace{4mm}

%\vspace{8mm}

\hrule

\vspace{4mm}

{Maquetaci�n y Grafismo}

{\bf DeMO}


\vspace{2mm}

\hrule

\vspace{4mm}

{Contacto Colaboradores}

{\tt colabora@ibolcode.net}

\vspace{2mm}

\hrule

\vspace{4mm}

{ Publicidad}

\vspace{1mm}

{\bf Occam's Razor Direct}

{\tt roor-ad@ibolcode.net}

\vspace{2mm}

\hrule

\vspace{4mm}

{Impresi�n}

{\bf Por ahora tu mismo\ldots Si te apetece}

\vspace{2mm}

\hrule

\hrule

\vspace{4mm}

{Editada por}


\vspace{20mm}

{\bf \copyright 2024 IbolCode/Occam's Razor

Esta obra est� bajo una licencia Reconocimiento 3.0 Espa�a de Creative
Commons. Para ver una copia de esta licencia, visite
http://creativecommons.org/\\licenses/by/3.0/es/ o envie una carta a
Creative Commons, 171 Second Street, Suite 300, San Francisco,
California 94105, USA.
}
\end{center}
\end{minipage}

\end{textblock}

\begin{textblock}{20}(3,2.0)

\begin{minipage}{.45\linewidth}
\colorbox{introcolor}{
\begin{minipage}{1\linewidth}

{{\resizebox{!}{1.0cm}{U}}{na vez m�s, despu�s de un largo letargo, volvemos a la carga con un nuevo y excitante n�mero de Occam's Razor y algunas novedades que os iremos comentando en las pr�ximas semanas.
}
}

\end{minipage}
}

\vspace{6mm}
La principal novedad y de la que os queremos informar antes de nada es de que cambiamos nuestra web. La revista Online Occam's Razor pasa a formar parte de un proyecto m�s grande y ambicioso que hemos bautizado como {\em Ibol Code}. Pod�is encontrar todos los detalles en nuestra nueva web:

\medskip

\begin{center}
\Large{\href{https://ibolcode.net}{https://ibolcode.net}}
\end{center}

\medskip

Esto tambi�n significa que hemos cambiado nuestras redes sociales. En la nueva p�gina web  encontrar las nuevas cuentas para estar a la �ltima sobre nuestras actividades, entre las que se encuentran publicar la revista con una periodicidad de 3 meses y emocionantes proyectos de los que todav�a no os vamos a desvelar ning�n detalle.

\bigskip

Respecto a este n�mero, os traemos varios art�culos que esperamos os resulten interesantes. {\em Don Bit0} y {\em Er Enigma} forman equipo para ofrecernos un mont�n de informaci�n sobre criptograf�a. Mientras {\em Er Enigma} nos explica en una nueva entrega de {\em Seguridad Para Meros Mortales} los conceptos b�sicos sobre criptograf�a que todos deber�amos conocer, {\em Don Bit0} se va al m�s bajo nivel y nos cuenta como utilizar OpenSSL para a�adir cifrado de nivel militar en nuestros programas.

\bigskip

{\em Signul Floid} nos trae un interesante art�culo sobre que es la compresi�n de datos y como funcionan los principales algoritmos que utilizamos a diario. Y nuestra querida {\em Richi} nos cuanta todos los secretos sobre la funci�n \texttt{alloca} y nos ayuda a mejorar en nuestro conocimiento del lenguaje de programaci�n C. {\em Nelly} nos cuenta como ha sido la historia de la incorporaci�n de la multitarea real en los sistemas operativos dom�sticos. Quiz�s m�s de uno se lleve una sorpresa... Os adelantamos que los Linuxeros no.

\bigskip

Incluimos dos nuevas secciones. Zona Geek, en la que comenzamos a programar un int�rprete de BrainFuck en ensamblador RISC-V :). Y Fen�menos extra�os en la que os contamos... bueno... fen�menos extra�os.

\bigskip

Esperamos que os guste, y como siempre esperamos vuestros comentarios, ideas, cr�ticas (siempre constructivas) para mejorar esta vuestra revista. Solo deciros que ya estamos trabajando en el siguiente n�mero y que como siempre, si tienes alg�n tema del que quieras escribir o quieras que escribamos, solo tienes que ponerte en contacto con nosotros.


\begin{flushright}
{\Large\sc{The Occam's Razor \\Team}}
\end{flushright}

\end{minipage}

\bigskip

\colorbox{introcolor}{
\begin{minipage}{.45\linewidth}

\bigskip

{\footnotesize\sf{\color{white}
Las opiniones expresadas en los art�culos, as� como los contenidos de
los mismos, son responsabilidad de los autores de �stos.

Puede obtener la versi�n electr�nica de esta publicaci�n, as� como el
{\em c�digo fuente} de la misma y los distintos ficheros de datos
asociados a cada art�culo en el sitio web:

{\tt http://ibolcode.net/roor}

}}

\bigskip

\end{minipage}
}


\end{textblock}

\pagebreak


\hypertarget{ratas:aes1}{}\label{ratas:aes}
\pagestyle{ratas}

\rput(7.9,-0.5){\resizebox{!}{14cm}{{\epsfbox{images/ratas/header-aes.eps}}}}

\psset{fillstyle=solid}
\psframe[fillcolor=black,opacity=0.7](2,-4.5)(17,0)



% -------------------------------------------------
% Cabecera
\begin{flushright}


{\color{introcolor}\mtitle{14cm}{Cifrado Militar en tus programas}}

\msubtitle{8cm}{Cifrando datos con OpenSSL}

{\sf\color{white}{ por Don Bit0}}

%{{\psset{linecolor=black,linestyle=dotted}\psline(-12,0)}}
\end{flushright}


\vspace{2mm}
% -------------------------------------------------


%\lstset{language=C,frame=tb,framesep=5pt,basicstyle=\footnotesize}
\lstset{language=C,frame=tb,framesep=5pt,basicstyle=\scriptsize}



\intro{introcolor}{E}{n el n�mero anterior os contamos como utilizar \emph{OpenSSL} paragenerar hashes utilizando una amplia variedad de algoritmos. Pero \emph{OpenSSL} nos permite mucho m�s. En este art�culo vamos a contaros como utilizar esta poderosa librer�a para encriptar datos utilizando distintos cifrados.}

\begin{multicols}{2}

Si, si, el t�tulo del art�culo se podr�a decir que es un poco \emph{Click Bait}, sin embargo si que os vamos a contar como cifrar datos utilizando varios algoritmos entre los que se encuentra AES-256. AES o \emph{Est�ndar de Encriptado Avanzado} (de sus siglas en ingl�s \emph{Advanced Encryption Standar}) es el algoritmo seleccionado oficialmente por el gobierno federal de los Estados Unidos. La versi�n con clave de 256 bits, conocida como AES-256 (eso no lo visteis venir eh?) se considera actualmente (en 2024, para los que le�is esto en el futuro), como el algoritmo de cifrado sim�trico m�s seguro.

El c�digo de este art�culo va a ser muy f�cil de seguir si hab�is le�do el
\href{http://www.papermint-designs.com/roor/2019-02-calculando-hashes-en-tus-programas-con-openssl}{art�culo anterior}. Sino, os recomiendo que lo hag�is, ya que la forma de cifrar y descrifrar datos es muy parecida a la forma de generar hashes y no vamos a volver a repetir aqu� los conceptos b�sicos. As� que, sino has le�do ese articulo, este es un buen momento pare hacer y as�, entender mejor lo que sigue.

\hypertarget{cuxf3digos-de-bloque-y-cuxf3digos-de-secuencia}{%
\sectiontext{white}{black}{C�DIGOS BLOQUE Y C�DIGOS DE SECUENCIA}\label{cuxf3digos-de-bloque-y-cuxf3digos-de-secuencia}}

Seg�n la forma en la que el algoritmo de cifrado consume los datos de
entrada, podemos clasificarlos en dos tipos: Bloque y Secuencia.

Los primeros consumen datos de entra en bloques de un tama�o definido,
mientras que los segundos los consumen bit a bit y suelen ser m�s
comunes en implementaciones HW. En este art�culo vamos a hablar de
cifrados de bloque, en concreto de AES como ya os adelantamos. Pero
primero tenemos que introducir un par de conceptos para poder entender
lo que sigue.

Un cifrado de bloque define un determinado tama�o de bloque con el que
va a trabajar. Y eso tiene dos consecuencias fundamentales. La primera
es que, normalmente la clave necesaria por el algoritmo suele ser del
mismo tama�o que el bloque (m�s sobre esto enseguida) y la segunda
es\ldots{} que pasa cuando no tenemos suficientes datos para llenar el
�ltimo bloque de nuestro mensaje cifrado?

Imaginemos que nuestro tama�o de bloque es de 128 bits o si lo prefer�s
16 bytes. Este es de hecho el caso para el algoritmo AES,
independientemente del tama�o de la clave que utilicemos.

Ahora imaginad que tenemos que cifrar un mensaje de 17 bytes. El
algoritmo toma los primeros 16 bytes y hace sus cosas para generar un
nuevo bloque de 16 bytes cifrado, pero que hacemos con el �ltimo byte
que nos queda por cifrar?. La respuesta es rellenar o hacer
\emph{Padding} como dicen en ingl�s, de forma que tengamos un nuevo
bloque que cifrar con nuestro algoritmo.

Sin embargo, no podemos coger simplemente el primer byte del bloque
cifrado. Tenemos que almacenar todo el bloque y disponer de un m�todo
para saber que solo un byte del �ltimo bloque forma parte del mensaje
original. As� que, en general, cuando ciframos un mensaje con un
algoritmo de bloque, el mensaje resultado siempre tendr� un tama�o mayor
o como mucho igual al mensaje original.

\hypertarget{padding}{%
\sectiontext{white}{black}{PADDING}\label{padding}}

Hay distintas estrategias para rellenar el bloque final, pero solo os
vamos a comentar dos de ellas. Para m�s informaci�n hay un excelente
art�culo en la
\href{https://en.wikipedia.org/wiki/Padding_(cryptography)}{wikipedia
sobre el tema}.

La forma m�s sencilla de rellenar el �ltimo bloque es usar ceros. Si
estamos cifrando cadenas de texto, el cero ya es un indicador de fin de
cadena, pero en el caso general, necesitar�amos alg�n tipo de cabecera
que nos informara del tama�o original. Super sencillo.

El otro tipo de relleno del que vamos a hablar es el utilizado por
OpenSSL y por eso nos interesa. Se conoce como PKCS\#7 y est� descrito
en el \href{https://tools.ietf.org/html/rfc5652\#section-6.3}{RFC5652}.
Este m�todo consiste en rellenar los espacios pendientes con el n�mero
de espacios pendientes en el �ltimo bloque. Imaginemos que nuestro
tama�o de bloque es de 8 bytes, en este caso tendr�amos las siguientes
posibilidades:

\begin{lstlisting}
|XX XX XX XX XX XX XX XX|XX 07 07 07 07 07 07 07|
|XX XX XX XX XX XX XX XX|XX XX 06 06 06 06 06 06|   
|XX XX XX XX XX XX XX XX|XX XX XX 05 05 05 05 05|       
|XX XX XX XX XX XX XX XX|XX XX XX XX 04 04 04 04|           
|XX XX XX XX XX XX XX XX|XX XX XX XX XX 03 03 03|               
|XX XX XX XX XX XX XX XX|XX XX XX XX XX XX 02 02|                   
|XX XX XX XX XX XX XX XX|XX XX XX XX XX XX XX 01|                       
\end{lstlisting}

Esos son todos los posibles patrones para el bloque final de un c�digo
bloque con tama�o 8 bytes. Si os fij�is, el �ltimo bloque es
problem�tico. Que pasar�a si en lugar de tener un byte que rellenar,
tuvi�ramos un bloque entero cuyo ultimo byte es
\verb!01!?\ldots{} No podr�amos diferenciarlo de otro
bloque al que le faltase el �ltimo byte. Lo que PKCS\#7 hace en este
caso es a�adir un bloque extra cuyo contenido es el tama�o del bloque.
As�, en el caso de que el tama�o del mensaje es multiplo del tama�o de
bloque (es decir, el �ltimo bloque est� relleno), el mensaje que
generar�a PKCS\#7 ser�a:

\begin{lstlisting}
|XX XX XX XX XX XX XX XX|08 08 08 08 08 08 08 08|
\end{lstlisting}

Como veremos en breve, OpenSSL genera siempre un bloque de datos extra
puesto que utiliza este tipo de rellenado.

\hypertarget{encriptando-datos}{%
\sectiontext{white}{black}{ENCRIPTANDO DATOS}\label{encriptando-datos}}

Ahora que estos peque�os detalles de implementaci�n han sido aclarados,
veremos que el proceso de encriptar y desencriptar datos es muy
sencillo. Comenzaremos viendo el c�digo para encriptar un mensaje con
m�s detalle ya que el proceso de descriptado es pr�cticamente id�ntico.

Vamos a utilizar un sencillo programa que encripta y desencripta un
mensaje que recibe de la entrada est�ndar. Como ver�is el API que ofrece
OpenSSL es muy f�cil de usar y podr�is modificarlo muy f�cilmente para
escribir vuestras propias utilidades.

\begin{entradilla}
{\em El {\color{introcolor}API \verb!EVP_! que ofrece OpenSSL} hace que resulte muy sencillo el cifrado/descifrado de datos }
\end{entradilla}

Dividiremos el c�digo en dos partes. La primera cifra el mensaje, y la
pod�is ver a continuaci�n. La segunda lo descifra y la ten�is que a�adir
tras esta parte, ya que usa las mismas variables y claves que la primera
parte.

{\scriptsize
\begin{lstlisting}[language=C]
#include <stdio.h>
#include <string.h>
#include <stdlib.h>
#include <openssl/evp.h>

int main (int argc, char *argv[]) {
 unsigned char key[] = {
   0xC7,0x75,0xE7,0xB7,0x57,0xED,0xE6,0x30,
   0xCD,0x0A,0xA1,0x11,0x3B,0xD1,0x02,0x66,
   0x1A,0xB3,0x88,0x29,0xCA,0x52,0xA6,0x42,
   0x2A,0xB7,0x82,0x86,0x2F,0x26,0x86,0x46
 };
 unsigned char iv[] = {
   0x15,0x25,0x35,0x45,0x55,0x65,0x75,0x85,
   0x16,0x26,0x36,0x46,0x56,0x66,0x76,0x86
 };
 char               txt[1024];
 unsigned char      txt_enc[1024];
 long unsigned long txt_len = 0;
 long unsigned long txt_enc_len = 0;
 int                len, r;

 const EVP_CIPHER *c =
                  EVP_get_cipherbyname(argv[1]);
 size_t bsize = EVP_CIPHER_get_block_size(c);

 // Initializamos los buffers
 memset (txt, 0, 1024); 
 memset (txt_enc, 0, 1024);
 
 // Leemos un mensaje de la entrada est�ndard
 printf ("%s", "Mensaje ? \n");
 fgets (txt, 1024, stdin);
 
 txt_len = strlen(txt);
 txt_enc_len = 0;
  
 // Creamos e inicializamos un contexto de cifrado 
 EVP_CIPHER_CTX *ctx;
 ctx = EVP_CIPHER_CTX_new();
 EVP_EncryptInit (ctx, c, key, iv);
 // Mientras tengamos bloques que cifrar.... Ciframos
 while (txt_len > bsize) {
  r = EVP_EncryptUpdate (ctx,
                          txt_enc + txt_enc_len, &len, 
                          txt + txt_enc_len, txt_len);
  txt_enc_len += len;
  txt_len -= len;
 }
 len = 0;
 // Procesa cualquier dato que quede pendiente
 EVP_EncryptFinal (ctx, txt_enc + txt_enc_len, &len);
 txt_enc_len += len;
 
 // Liberamos el contexto
 EVP_CIPHER_CTX_free (ctx);
 
 printf ("Total mensaje cifrado : %d\n", txt_enc_len);
 BIO_dump_fp (stdout, txt_enc, txt_enc_len);
 
 // Descifrado
 // (...)
\end{lstlisting}
}
\hypertarget{seleccionando-el-cifrado}{%
\sectiontext{white}{black}{SELECCIONANDO EL CIFRADO}\label{seleccionando-el-cifrado}}

Lo primero que hace el programa es seleccionar el tipo de cifrado que
vamos a utilizar, el cual pasaremos como primer par�metro al programa.
Una vez seleccionado determina el tama�o de bloque del algoritmo en
cuesti�n. Y ya sabemos por que eso es importante.

\begin{lstlisting}[language=C]
  const EVP_CIPHER *c =
                 EVP_get_cipherbyname (argv[1]);
  size_t bsize = EVP_CIPHER_get_block_size(c);
\end{lstlisting}

Para obtener una lista de los posibles algoritmos, podemos utilizar la
utilidad \verb!openssl!. En nuestro sistema obtenemos
la siguiente salida

\end{multicols}

\begin{center}
{\normalsize
\begin{verbatim}
$ openssl
(...)
Cipher commands (see the `enc' command for more details)
aes-128-cbc       aes-128-ecb       aes-192-cbc       aes-192-ecb
aes-256-cbc       aes-256-ecb       aria-128-cbc      aria-128-cfb
aria-128-cfb1     aria-128-cfb8     aria-128-ctr      aria-128-ecb
aria-128-ofb      aria-192-cbc      aria-192-cfb      aria-192-cfb1
aria-192-cfb8     aria-192-ctr      aria-192-ecb      aria-192-ofb
aria-256-cbc      aria-256-cfb      aria-256-cfb1     aria-256-cfb8
aria-256-ctr      aria-256-ecb      aria-256-ofb      base64
bf                bf-cbc            bf-cfb            bf-ecb
bf-ofb            camellia-128-cbc  camellia-128-ecb  camellia-192-cbc
camellia-192-ecb  camellia-256-cbc  camellia-256-ecb  cast
cast-cbc          cast5-cbc         cast5-cfb         cast5-ecb
cast5-ofb         des               des-cbc           des-cfb
des-ecb           des-ede           des-ede-cbc       des-ede-cfb
des-ede-ofb       des-ede3          des-ede3-cbc      des-ede3-cfb
des-ede3-ofb      des-ofb           des3              desx
rc2               rc2-40-cbc        rc2-64-cbc        rc2-cbc
rc2-cfb           rc2-ecb           rc2-ofb           rc4
rc4-40            seed              seed-cbc          seed-cfb
seed-ecb          seed-ofb          sm4-cbc           sm4-cfb
sm4-ctr           sm4-ecb           sm4-ofb
\end{verbatim}
}
\textbf{Lista de Algoritmos de cifrado de OpenSSL}

\end{center}
\begin{multicols}{2}

Nuestro programa podr� utilizar cualquiera de estos algoritmos para
cifrar nuestros datos. Como pod�is ver, para cada algoritmo hay varias
opciones. Algunas indican el tama�o de clave, como sucede con
\verb!aes!, \verb!aria! o \verb!camellia!, mientras que otras definen modos de
operaci�n como los terminados en {\verb!-cbc!} o
\verb!-ecb! por ejemplo. Hablaremos de esto un poco
m�s tarde.

\hypertarget{cifrando-datos}{%
\sectiontext{white}{black}{CIFRANDO DATOS}\label{cifrando-datos}}

El proceso de cifrado con el API \verb!EVP_! de
\emph{OpenSSL} es igual al proceso de generaci�n de hashes, pero con
c�digos de cifrado en lugar de con funciones hash. As� que los pasos son
los mismos:

\begin{itemize}
%\tightlist
\item
  Genera un objeto contexto para mantener el estado de nuestro cifrado
\item
  Inicializa el cifrado con los valores necesarios
\item
  Pasa la informaci�n a cifrar al contexto y almacena el resultado
\item
  Procesa lo que este pendiente y terminar
\end{itemize}

Si comparamos este c�digo con el que usamos para generar hashes podemos
ver dos diferencias principales (dejando a un lado el nombre de las
funciones)

{\scriptsize
\begin{lstlisting}[language=C]
EVP_CIPHER_CTX *ctx = EVP_CIPHER_CTX_new();
EVP_EncryptInit (ctx, c, key, iv);

while (txt_len > bsize) {
  r = EVP_EncryptUpdate (ctx, txt_enc + txt_enc_len, &len, 
                              txt + txt_enc_len, txt_len);
  txt_enc_len += len;
  txt_len -= len;
}
EVP_EncryptFinal (ctx, txt_enc + txt_enc_len, &len);
text_enc_len += len;
EVP_CIPHER_CTX_free (ctx);
  
BIO_dump_fp (stdout, txt_enc, txt_enc_len);
\end{lstlisting}
}

La primera es que en la funci�n de inicializaci�n es en la que pasamos
la clave y el vector inicial. Olvidaros por ahora del vector de
inicializaci�n \verb!IV!, volveremos sobre esto un
poco m�s tarde.

La segunda es que cuando llamamos a
\verb!EVP_EncryptUpdate! tenemos que incrementar
ambos punteros el de los datos originales y el del resultado. Recordad
que con el hash, los resultado siempre se almacenaban en el mismo sitio,
pero aqu� tenemos que cifrar los datos originales y cada bloque que
cifremos tenemos que almacenarlo en su correspondiente bloque de
memoria.

La funci�n \verb!EVP_EncryptUpdate! permite utilizar
el mismo buffer como entrada y salida siempre que almacenemos los datos
en el mismo sitio de donde los sacamos. En este caso hemos decidido
utilizar dos buffers separados por claridad. Pero que lo sep�is\ldots{}

\begin{entradilla}
{\em Durante el proceso de {\color{introcolor}cifrado y descifrado} debemos tener en cuanta el \textbf{ Padding} }
\end{entradilla}


La otra es que, como comentamos cuando hablamos del \emph{Padding},
debemos procesar los datos hasta que nos queden menos bytes que los que
conforman un bloque (condici�n en el \verb!while!), y
finalmente llamar a \verb!EVP_EncryptFinal! para
procesar ese �ltimo bloque, ya sea el que hemos rellenado o hemos
a�adido seg�n el algoritmo PKCS\#7.

A parte de eso, como pod�is ver, la forma de cifrar datos es muy,muy
sencilla. Una vez que los datos est�n cifrados, en este programa de
ejemplo, nosotros los volcamos en la consola, para echarles un ojo.
Podr�amos escribir nuestra propia funci�n, pero en este caso utilizamos
la funci�n \verb!BIO_dump_fp! que propociona
\emph{OpenSSL}.

\hypertarget{desencriptando-datos}{%
\sectiontext{white}{black}{DESCIFRANDO DATOS}\label{desencriptando-datos}}

Lo prometido es deuda, y aqu� ten�is la segunda parte del programa que
desencripta los datos que acabamos de encriptar. Simplemente a�adid este
c�digo a continuaci�n del primero para tener el programa completo.

{\scriptsize
\begin{lstlisting}[language=C]
  // Destruir el buffer de texto plano
  memset (txt, 0, 1024); 
  txt_len = 0;

  // Creammos e inicializamos un nuevo contexto de cifrado
  ctx = EVP_CIPHER_CTX_new();
  EVP_DecryptInit (ctx, c, key, iv);

  // Mientras tengamos datos pendientes, desencriptamos
  while (txt_enc_len > bsize) {
    EVP_DecryptUpdate (ctx, txt+txt_len, &len,
                       txt_enc+txt_len, txt_enc_len);
    txt_len += len;
    txt_enc_len -= len;
  }
  len = 0;
      BIO_dump_fp (stdout, txt, txt_len);
  if ((r = EVP_DecryptFinal (ctx, txt + txt_len, &len)) == 0) 
    fprintf (stderr, "Datos de entrada corruptos\n");
  txt_len += len;
  BIO_dump_fp (stdout, txt, txt_len);
  EVP_CIPHER_CTX_free (ctx);
  }
\end{lstlisting}
}

Como pod�is ver el c�digo es casi id�ntico. Procesamos todos los bloques
excepto el �ltimo. Puesto que no hemos desativado el padding, el �ltimo
bloque es siempre un bloque de padding y debe ser procesado con
\verb!EVP_Decrypt_Final!. Esta funci�n comprobar�
cuantos bytes de padding fueron usados y nos devolver� en
\verb!len! el n�mero de bytes final que debemos
a�adir a nuestro buffer.

Al final del programa volcamos el buffer con el tama�o real, tras la
correci�n con el valor devuelto por
\verb!EVP_DecryptFinal! y con el tama�o de bloque,
para que pod�is ver el padding.

\hypertarget{probando-nuestro-ejemplo}{%
\sectiontext{white}{black}{PROBANDO NUESTRO EJEMPLO}\label{probando-nuestro-ejemplo}}

Es hora de compilar y probar nuestro ejemplo. Para compilar este
programa tenemos que utilizar la librer�a
\verb!libcrypto! tal que as�:

\begin{lstlisting}
$ gcc -g -o acrypt acrypt.c -lcrypto
\end{lstlisting}

Y ahora ya podemos ejecutarlo con distintos algoritmos. Estos son
algunos ejemplos. Primero veamos que sucede cuando la entrada es un
m�ltiplo de bloque.

\begin{entradilla}
{\em {\color{introcolor}OpenSSL a�ade un bloque extra} durante el cifrado si el tama�o del mensaje es un m�ltiplo exacto del tama�o de bloque}
\end{entradilla}


{\scriptsize
\begin{lstlisting}
# (perl -e 'print "A"x64;'| ./k aes-256-ecb)
Mensaje ?
Total mensaje cifrado : 80
0000 - 2e bf 50 07 e2 c0 d6 60-d8 f3 dc 96 71 c1 e8 72   ..P....`....q..r
0010 - 2e bf 50 07 e2 c0 d6 60-d8 f3 dc 96 71 c1 e8 72   ..P....`....q..r
0020 - 2e bf 50 07 e2 c0 d6 60-d8 f3 dc 96 71 c1 e8 72   ..P....`....q..r
0030 - 2e bf 50 07 e2 c0 d6 60-d8 f3 dc 96 71 c1 e8 72   ..P....`....q..r
0040 - 82 34 3f 32 92 bb 77 28-82 70 35 7e d8 40 e8 5f   .4?2..w(.p5~.@._
------------------------------
0000 - 41 41 41 41 41 41 41 41-41 41 41 41 41 41 41 41   AAAAAAAAAAAAAAAA
0010 - 41 41 41 41 41 41 41 41-41 41 41 41 41 41 41 41   AAAAAAAAAAAAAAAA
0020 - 41 41 41 41 41 41 41 41-41 41 41 41 41 41 41 41   AAAAAAAAAAAAAAAA
0030 - 41 41 41 41 41 41 41 41-41 41 41 41 41 41 41 41   AAAAAAAAAAAAAAAA
\end{lstlisting}
}




Como pod�is ver, al cifrar 64 bytes que es un m�ltiplo del tama�o de
bloque de AES-256, durante el cifrado a�adimos un nuevo bloque con el
valor \verb!0x10! que representa el padding PKCS.
Tambi�n podemos ver como, tras utilizar el valor devuelto por
\verb!EVP_DecryptFinal! obtenemos el tama�o correcto
del mensaje original.



Veamos que pasa si nos pasamos por dos bytes.

{\scriptsize
\begin{lstlisting}
# (perl -e 'print "A"x66;'| ./k aes-256-ecb)
Mensaje ?
Total mensaje cifrado : 80
0000 - 2e bf 50 07 e2 c0 d6 60-d8 f3 dc 96 71 c1 e8 72   ..P....`....q..r
0010 - 2e bf 50 07 e2 c0 d6 60-d8 f3 dc 96 71 c1 e8 72   ..P....`....q..r
0020 - 2e bf 50 07 e2 c0 d6 60-d8 f3 dc 96 71 c1 e8 72   ..P....`....q..r
0030 - 2e bf 50 07 e2 c0 d6 60-d8 f3 dc 96 71 c1 e8 72   ..P....`....q..r
0040 - bb f1 ef 09 1b 99 d0 4c-d1 30 bd b4 e9 5a 15 70   .......L.0...Z.p
------------------------------
0000 - 41 41 41 41 41 41 41 41-41 41 41 41 41 41 41 41   AAAAAAAAAAAAAAAA
0010 - 41 41 41 41 41 41 41 41-41 41 41 41 41 41 41 41   AAAAAAAAAAAAAAAA
0020 - 41 41 41 41 41 41 41 41-41 41 41 41 41 41 41 41   AAAAAAAAAAAAAAAA
0030 - 41 41 41 41 41 41 41 41-41 41 41 41 41 41 41 41   AAAAAAAAAAAAAAAA
0040 - 41 41                                             AA
\end{lstlisting}
}

Mismo resultado, solo que ahora podemos ver el algoritmo de padding en
vivo. Al usar solo dos bytes del �ltimo bloque, debemos a�adir 14 bytes
para completar el bloque que, en hexadecimal es
\verb!0x0e!.

Bueno, parece que funciona bastante bien. Pod�is obtener una lista de
los algoritmos disponibles en vuestro sistema utilizando el comando:

\begin{lstlisting}
openssl list -cipher-algorithms
\end{lstlisting}

\hypertarget{porque-no-me-puedo-cifrar-con-rc4}{%
\sectiontext{white}{black}{�PORQU� NO PUEDO CIFRAR CON RC4?}\label{porque-no-me-puedo-cifrar-con-rc4}}

Has probado a pasar como par�metro a nuestro programa
\verb!rc4! o \verb!bf! y no ha
sucedido nada?. No te preocupes, ahora mismo lo solucionamos. Lo primero
que debemos hacer es modificar nuestro programa ligeramente para que nos
muestre mensajes de error, para ello debemos a�adir el fichero de
cabecera \verb!<openssl/err.h>!. Ahora podemos
realizar la siguiente modificaci�n:

{\scriptsize
\begin{lstlisting}[language=C]
#include <openssl/provider.h>
(...)
main (...)

  while (txt_len >= bsize) {
    r = EVP_EncryptUpdate (ctx, txt_enc + txt_enc_len, &len, 
                        txt     + txt_enc_len, txt_len);
    if (!r) {ERR_print_errors_fp (stderr);exit (1);}
    txt_enc_len += len;
    txt_len -= len;
    }
(...)
\end{lstlisting}
}
La funci�n \verb!EVP_EncryptUpdate! retorna
\verb!0! si se produce un error. Si ahora
recompilamos y ejecutamos esta nueva versi�n de nuestro programa
obtendr�amos:

\end{multicols}


{\scriptsize
\begin{lstlisting}
$ (perl -e 'print "A"x65;'| ./k bf)
BLock Size : 8 (0x7f7530432b80)
Mensaje ?
40E7DE2F757F0000:error:0308010C:digital envelope routines: inner_evp_generic_fetch:unsupported:../crypto/evp/evp_fetch.c:
373:Global default library context, Algorithm (BF-CBC : 11), Properties ()
40E7DE2F757F0000:error:03000083:digital envelope routines: EVP_EncryptUpdate:no cipher set:../crypto/evp/evp_enc.c:665:
\end{lstlisting}
}

\begin{multicols}{2}
Si bien el mensaje parece que est� cifrado con el propio algoritmo
\verb!bf! (que es como OpenSSL conoce al algoritmo
\emph{BlowFish}), lo que realmente nos est� diciendo es que no puede
encontrar ese algoritmo para utilizarlo.

Las versiones m�s modernas de OpenSSL, est�n configuradas por defecto
para no utilizar algoritmos que se no se consideran seguros. OpenSSL
utiliza lo que ellos llaman proveedores (\emph{Providers} en ingl�s).
Existen varios tipos pero los que nos interesan son los proveedores por
defecto y los llamados \verb!legacy! que incluyen
algoritmos digamos\ldots{} viejos u obsoletos. Podemos activar eso
algoritmos \verb!legacy! modificando el fichero de
configuraci�n de OpenSSL, o a�adiendo unas pocas l�neas de c�digo a
nuestro programa.

En cualquier caso, todos los detalles los pod�is encontrar en el fichero
\verb!README-PROVIDERS.md! del
\href{https://github.com/openssl/openssl/blob/master/README-PROVIDERS.md}{github
official de OpenSSL}.

\hypertarget{tamauxf1os-de-bloque-y-claves}{%
\sectiontext{white}{black}{TAMA�OS DE BLOQUE Y CLAVE}\label{tamauxf1os-de-bloque-y-claves}}

En los primeros algoritmos criptogr�ficos, como DES (\emph{Data
Encription Standard}), el tama�o de la clave coincid�a con el tama�o del
bloque. En los algoritmos actuales ese ya no es el caso. Por ejemplo,
AES-256 utiliza una clave de 256 bits, pero su tama�o de bloque es de
128 bits. AES-128 utiliza una clave de 128 bits, aunque su tama�o de
bloque es tambi�n 128 bits.

Lo que estos nuevos tipos de cifrado utilizan es un algoritmo de
transformaci�n de la clave original. Estos algoritmos se suelen conocer
como \emph{Key Schedule}, algo as� como planificaci�n de clave. A veces
se utilizan para crear claves m�s peque�as a partir de claves mayores, o
si lo prefer�s muchas claves de un determinado tama�o (como
\href{https://en.wikipedia.org/wiki/AES_key_schedule}{AES}), otras para
generar una clave m�s grande a partir de una m�s corta (como en
\href{https://en.wikipedia.org/wiki/RC4\#Key-scheduling_algorithm_(KSA)}{RC4}).

\begin{entradilla}
{\em Muchos algoritmos criptogr�ficos {\color{introcolor}utilizan un algoritmo Jey Schedule} para extender, comprimir o generar nuevas claves}
\end{entradilla}


Lo que es importante tener claro es que si bien el tama�o del bloque y
el tama�o de clave no tienen porque coincidir, de alguna forma,
internamente es necesario utilizar una clave del tama�o del bloque usado
por el algoritmo y cuando ese no es el caso, existen algoritmos para
conseguir claves del tama�o adecuado, a partir de la clave original.
Para entender lo que significa este �ltimo p�rrafo y saber para que es
esa variable \verb!iv! de la que no hemos dicho nada
a prop�sito, tenemos que introducir los modos de cifrado.

\hypertarget{modos-de-cifrado}{%
\sectiontext{white}{black}{MODOS DE CIFRADO}\label{modos-de-cifrado}}

Record�is el \emph{Vector de Inicializacion} (\emph{IV} en ingl�s) del
que hablamos al principio?. Bueno, es hora de volver sobre este tema.
Pero primero quiero mostraros algo:

{\scriptsize
\begin{lstlisting}
# (perl -e 'print "A"x64;'| ./k aes-256-ecb)
BLock Size : 16 (0x7f6eaa38b740)
Mensaje ?
Total mensaje cifrado : 80
0000 - 9e 9d 05 b2 b8 8e 96 1c-1d e0 0c 90 27 9e a3 42   ............'..B
0010 - 9e 9d 05 b2 b8 8e 96 1c-1d e0 0c 90 27 9e a3 42
............'..B
0020 - 9e 9d 05 b2 b8 8e 96 1c-1d e0 0c 90 27 9e a3 42   ............'..B
0030 - 9e 9d 05 b2 b8 8e 96 1c-1d e0 0c 90 27 9e a3 42
............'..B
0040 - 31 c7 81 8d 54 9d c2 6c-c1 8b 04 71 38 42 17 8d   1...T..l...q8B..
(...)
\end{lstlisting}
}
Como pod�is ver, si utilizamos el modo \verb!ECB!
(\emph{Electronic Code Book} o \emph{Libro de c�digos electr�nico}),
vemos como los primeros 16 bytes (el tama�o de bloque de AES) se repiten
en bloques sucesivos, excepto para el �ltimo bloque, el cual, si
record�is, es un bloque de relleno que contiene el valor
\verb!0x10! repetido. A primera vista esto no parece
un grave problema, pero, puesto que, bloques id�nticos en el fichero
original se codificaran de la misma forma, ciertos patrones se podr�n
identificar incluso si el fichero est� cifrado. Al final del art�culo
veremos un ejemplo pr�ctico.

Para evitar este problema se introdujeron formas de manejar los datos
para que una misma secuencia de bytes no produzca siempre el mismo
resultado. B�sicamente, estos modos de cifrado simplemente toman alg�n
dato del bloque anteriormente cifrado y lo usan para modificar los datos
de entrada en el siguiente bloque. Existen 4 modos de cifrado cl�sicos:

\begin{itemize}
%\tightlist
\item
  \verb!ECB! \emph{Electronic Code Book}. Este es el
  modo b�sico en el que los datos de entrada se dividen en bloques y
  cada uno de ellos se cifra tal cual utilizando la clave proporcionada
  por el usuario.
\item
  \verb!CBC! \emph{Cipher Block Chaining}. En este
  modo, cada uno de los bloques de datos a cifrar se modifican aplicando
  una operaci�n \verb!xor! con el resultado del
  cifrado del bloque anterior.
\item
  \verb!CFB! \emph{Cipher FeedBack}. En este modo,
  cada bloque de datos cifrados se calcula aplicando una operaci�n
  \verb!xor! entre el bloque sin cifrar y el
  resultado del cifrado del bloque anterior.
\item
  \verb!OFB! \emph{Ouput Feedback}. Este modo
  funciona igual que el anterior, pero en lugar de pasar como entrada al
  siguiente bloque el resultado del \verb!xor! con el
  texto plano, env�a la salida del bloque actual.
\end{itemize}

En la actualidad se han definido algunos modos m�s, sin embargo,
conceptualmente, todos funcionan igual, as� que, sin p�rdida de
generalidad, como dicen los matem�ticos, vamos a hablar solo de el modo
\verb!CBC!.


\hypertarget{el-infame-iv}{%
\sectiontext{white}{black}{EL INFAME IV}\label{el-infame-iv}}

Como acabamos de decir, vamos a centrarnos, a modo de ejemplo, en el
modo \verb!CBC!. Veamos el diagrama de bloques de
este modo:



{\scriptsize
\begin{lstlisting}
          BLOQUE1              BLOQUE2
             |                    |
??? ------> XOR     +----------> XOR
             |      |             | 
             v      |             v
         +------+   |          +------+
CLAVE -->| ALGO |   | CLAVE -->| ALGO |
         +------+   |          +------+ 
             |      |              |
             +------+              +-----> Siguiente
             |                     |       Bloque
             v                     v
         CIFRADO1               CIFRADO2
         
\end{lstlisting}
}

Para el caso de AES-256, en el diagrama anterior, los bloques de entrada
(\verb!BLOQUE1! y \verb!BLOQUE2!)
contienen los datos a encriptar en bloques de 128 bits. Recordad el
tama�o de bloque para AES es de 128 bits independientemente de la clave.
Ahora ya sab�is lo que esto significa\ldots{} cada bloque de 16 bytes es
lo que se pasa por la caja que hemos llamado
\verb!ALGO! en el diagrama.

Cada uno de los bloques \verb!ALGO! implementan el
algoritmo de cifrado seleccionado, el cual necesita la clave del usuario
para funcionar. Para AES, esa clave debe ser tambi�n de 128 bits y como
pod�is ver es un par�metro de entrada para cada uno de los bloques
\verb!ALGO!. Cada una de esas
\verb!CLAVE!s son calculadas a partir de la clave
original utilizando un algoritmo de varios pasos. Dependiendo del tama�o
de la clave (128, 192 o 256), el algoritmo aplica un n�mero determinado
de iteraciones para generar la clave de 128bits que necesitamos para
cada bloque.

\begin{entradilla}
{\em {\color{introcolor} AES} puede utilizar claves de 128, 192 y 256 bits, si bien, el tama�o de bloque es siempre 128 bits}
\end{entradilla}


La salida del \verb!ALGO! es el resultado del cifrado
del bloque inicial. Es en este momento en el que el modelo de cifrado
\verb!CBC! entra en juego. Ese mismo resultado del
primer cifrado se propaga al siguiente bloque en el que realizamos un
\verb!XOR! entre este resultado y el siguiente bloque
a cifrar. De esta forma, conseguimos que aunque
{\verb!BLOQUE1!} y {\verb!BLOQUE2!}
sean exactamente iguales, el resultado del cifrado de estos dos bloques
ser� totalmente diferente, puesto que estamos modificando la entrada al
algoritmo. La operaci�n XOR es muy guay por que es muy f�cil de
invertir.



Bien, y aqu� es donde entra en juego el infame
{\verb!IV!}. Si os fij�is, utilizamos la salida de cada
etapa de cifrado para modificar los datos de entrada a la siguiente.
Esto funciona genial, para todas las etapas menos para la primera. Y ah�
lo ten�is, ese es el {\verb!IV!} el valor que
utilizamos en la primera etapa (por eso lo de inicializaci�n) para
modificar los datos del primer bloque a cifrar.

Cuando trabajamos en modo {\verb!ECB!} cada uno de los
bloques del diagrama anterior es independiente (no hay datos que pasan
de una etapa a otra) y por lo tanto no necesitamos un vector de
inicializaci�n. Como pod�is observar, nuestro programa siempre pasa un
valor de {\verb!IV!} al
{\verb!CIPHER!} sin embargo, para los modos
{\verb!ECB!}, ese valor simplemente se ignora.

En general, el valor de IV no tiene por que ser secreto pero si
aleatorio o, dependiendo del tipo de algoritmo con el que lo usemos,
solo es necesario que sea impredecible o �nico. Este valor normalmente
se almacena junto con los datos cifrados ya que es necesario para poder
decodificar el mensaje si bien, conocer su valor no ofrece ninguna
informaci�n �til para descifrar el mensaje.

\hypertarget{iv-en-aes-256}{%
\sectiontext{white}{black}{IV EN AES-256}\label{iv-en-aes-256}}

Con todo esto nuevo que sabemos, vamos a echar un ojo, una vez m�s, a la
salida de nuestro programa cuando utilizamos el modo
{\verb!ECB!} con un patr�n claramente definido:

{\scriptsize
\begin{lstlisting}
$ (perl -e 'print "A"x64;'| ./k aes-256-ecb)
BLock Size : 16 (0x7f6eaa38b740)
Mensaje ?
Total mensaje cifrado : 80
0000 - 9e 9d 05 b2 b8 8e 96 1c-1d e0 0c 90 27 9e a3 42   ............'..B
0010 - 9e 9d 05 b2 b8 8e 96 1c-1d e0 0c 90 27 9e a3 42
............'..B
0020 - 9e 9d 05 b2 b8 8e 96 1c-1d e0 0c 90 27 9e a3 42   ............'..B
0030 - 9e 9d 05 b2 b8 8e 96 1c-1d e0 0c 90 27 9e a3 42
............'..B
0040 - 31 c7 81 8d 54 9d c2 6c-c1 8b 04 71 38 42 17 8d   1...T..l...q8B..
(...)
\end{lstlisting}
}

Los primeros cuatro bloques id�nticos se deben a que la implementaci�n
de {\verb!ECB!} de \emph{OpenSSL} utiliza bloques de 16
bytes como ya hemos visto. Como cada etapa de cifrado es independiente,
obtenemos 4 bloques id�nticos. El �ltimo bloque es el de relleno y por
lo tanto es diferente a los anteriores.

Si utiliz�ramos el modo {\verb!CBC!}, el resultado es
bastante diferente:

{\scriptsize
\begin{lstlisting}
$ (perl -e 'print "A"x64;'| ./k aes-256-cbc)
BLock Size : 16 (0x7f5382f33940)
Mensaje ?
Total mensaje cifrado : 80
0000 - 51 c8 0e f0 60 34 5b 67-a1 60 45 4a 9d 88 2d ac   Q...`4[g.`EJ..-.
0010 - c0 a1 c0 6b db 0e 34 11-4e 06 04 15 f2 f3 0c 36   ...k..4.N......6
0020 - b1 ea a6 8f f1 10 1c 97-d7 ea 25 79 cc 09 01 88   ..........%y....
0030 - 50 51 b8 9a 76 31 3c 23-48 2c 3a 1f f4 1a c7 29   PQ..v1<#H,:....)
0040 - ea c5 9b 05 ee 46 4f 69-da cb 7a 37 38 6f 7e 72   .....FOi..z78o~r
\end{lstlisting}
}

Como podemos ver, ahora cada bloque es diferente y no se puede deducir
nada sobre la secuencia de datos original. En este sencillo ejemplo,
esto no parece un problema muy grave, pero bajo ciertas circunstancias y
dependiendo de nuestro caso de uso puede ser un problema. Veamos un
ejemplo.

Seleccionad una imagen que contenga un dibujo, es decir, una imagen que
tenga grandes �reas con un mismo color. Yo he elegido a nuestro querido
Tux.

\myfig{0}{images/ratas/tux.eps}{0.6}


Ahora convertiremos nuestra imagen al formato PGM. Este formato no tiene
nada de especial, simplemente, como veremos en un segundo, har� m�s
f�cil manipular la imagen de la forma que necesitamos.

\begin{lstlisting}
$ convert tux.jpg tux.pgm
\end{lstlisting}

\begin{quote}
Necesitar�is el paquete {\verb!imagemagick!} para poder
disponer de la utilidad {\verb!convert!}
\end{quote}

Lo interesante de el formato {\verb!PGM!} es que la
cabecera son solo 3 l�neas de texto. En nuestro caso est� (con la imagen
de Tux de la wikipedia) es la cabecera.

{\scriptsize
\begin{lstlisting}
P5
1727 2047
255
\end{lstlisting}
}





Ahora debemos eliminar estas 3 l�neas. Puedes abrir el fichero en tu
editor de textos favorito y borrarlas o:

{\scriptsize
\begin{lstlisting}
$ sed -i '1,3d' tux.pgm
\end{lstlisting}
}

Ahora encriptamos el fichero con la utilidad
{\verb!openssl!} o quiz�s utilizando tu propio program
;).



{\scriptsize
\begin{lstlisting}
openssl enc -aes-256-cbc -p -in tux.pgm -out tux-cbc.enc
openssl enc -aes-256-ecb -p -in tux.pgm -out tux-ecb.enc
\end{lstlisting}
}

\begin{entradilla}
{\em {\color{introcolor} EL modo de cifrado ECB} no es seguro cuando los datos a cifrar contienen patrones repetitivos}
\end{entradilla}



Ahora solo tenemos que poner la cabecera otra vez a los ficheros:

{\scriptsize
\begin{lstlisting}
$ cat << EOM > file1.pgm
> P5
> 1727 2047
> 255
> EOM
$ cp file1.pgm file2.pgm
$ cat tux-cbc.enc >> file1.pgm
$ cat tux-ecb.enc >> file2.pgm
\end{lstlisting}
}

Ahora abre los ficheros con tu visor de im�genes preferido y ver�s cual
es el problema con el modo ECB. Os reproducimos las dos im�genes al final del art�culo.

\hypertarget{conclusiones}{%
\sectiontext{white}{black}{CONCLUSIONES}}\label{conclusiones}

En este peque�o art�culo hemos visto como utilizar \emph{OpenSSL} para
encriptar y desencriptar datos utilizando cualquiera de los algoritmos
que nos proporciona. Hemos prestado algo m�s de atenci�n a AES ya que es
el algoritmo sim�trico m�s seguro hasta la fecha y hemos visto los
principales modos de cifrado (hay algunos m�s) con los que utilizar
algoritmos de cifrado por bloque como AES, lo cual nos ha permitido
entender la funci�n del vector de inicializaci�n (normalmente referido
como {\verb!IV!}) que se ve de vez en cuando por ah�.


%\raggedcolumns
%\clearpage
\end{multicols}

\begin{center}
\includegraphics[width=0.4\hsize]{images/ratas/tux1.eps}
\includegraphics[width=0.4\hsize]{images/ratas/tux2.eps}

\textbf{Imagen de Tux cifrada en mode ECB y en modo CBC}
\end{center}
%\clearpage
%\pagebreak

\hypertarget{linux:kernel1}{}\label{linux:kernel}
\pagestyle{linux}

\rput(8.0,-2.0){\resizebox{!}{12cm}{{\epsfbox{images/header-linux.eps}}}}

\psset{fillstyle=solid}
\psframe[fillcolor=black,opacity=0.6](2,-4.5)(17,0)



% -------------------------------------------------
% Cabecera
\begin{flushright}


{\color{teal}\mtitle{14cm}{Multitarea Apropiativa en ordenadores Dom�sticos}}

\msubtitle{10cm}{Una breve retrospectiva sobre innovaci�n}

{\sf\color{white}{ por Nelly Kerningham}}

{\psset{linecolor=black,linestyle=dotted}\psline(-12,0)}
\end{flushright}

\vspace{2mm}
% -------------------------------------------------



\lstset{language=C,frame=tb,framesep=5pt,basicstyle=\scriptsize}


\intro{introcolor}{E}{s lamentable como el tiempo y la propaganda tergiversa la realidad y
verdades a medias o mentiras completas se acaban convirtiendo en la
realidad de muchas personas. En este articulo vamos a explorar algunas
de esas \emph{inocentes} falacias que pulula por las mentes de la
sociedad.}

\begin{multicols}{2}


Esta es una cuesti�n que quiz�s resulte sorprendente a muchos de
vosotros. Lo fu� para mi, si bien, al tener la suerte de haber vivido
aquellos a�os hab�a cosas que no encajaban. Vamos a hacer un breve
repaso de la evoluci�n de los sistemas operativos m�s populares.

En esta secci�n vamos a ver las capacidades de los kernel de los OSes de
microsoft y Apple comparados con Linux. Pero primero, una breve introduci�n para que sep�is de que vamos a hablar en el resto del art�culo

\hypertarget{kernel}{%
\sectiontext{white}{black}{MULTITAREA APROPIATIVA}\label{kernel}}

Antes del advenimiento de los sistemas multi-core, la multitarea solo se pod�a conseguir repartiendo la CPU entre los procesos, es decir, ejecutando los programas a los pocos. Existen dos formas fundamentales de hacer esto que se conocen como: Multitarea colaborativa y Multitarea apropiativa (del ingl�s {\em preemptive}).

En la primera, el sistema operativo nunca inicia lo que se conoce como cambio de contexto (parar un proceso e iniciar otro). Cada proceso sigue ejecut�ndose hasta que decida liberar la CPU voluntariamente, o en ciertas circunstancias como si debe esperar por una operaci�n de entrada/salida u otro evento. En estos sistema, un error de programaci�n en cualquier proceso puede provocar que todo el sistema se bloquee.

\begin{entradilla}
{\em Los sistemas multiproceso se pueden implementar de forma {\color{introcolor}Colaborativa o Apropiativa.}}
\end{entradilla}

En la segunda, el sistema operativo asigna un tiempo fijo a cada tarea. Pasado ese tiempo (o en caso de que el proceso se bloquee esperando por alg�n evento), el sistema operativo asigna la CPU a otro proceso. En estos sistemas, aunque un proceso se bloquee, el sistema seguir� funcionando sin problemas.

Bien, como os pod�is imaginar, la primera es much�simo m�s sencilla de implementar y tambi�n mucho peor puesto que cualquier proceso puede bloquear la ejecuci�n del sistema operativo y por tanto de todo el ordenador.

�Qu� sistema operativo dom�stico cre�is que fue el primero en ofrecer este tipo de sistema operativo?. Seguid leyendo si no conoc�is la respuesta.




\hypertarget{microsoft}{%
\sectiontext{white}{black}{MICROSOFT}\label{microsoft}}

Dejando a un lado el DOS (\emph{Disk Operating System}) y centr�ndonos
en los otros sistemas operativos de esta compa��a, los conocidos como
\emph{Windows}, los cuales curiosamente no fueron, ni mucho menos los
primeros entornos de ventanas, ni siquiera para los PCs de la �poca.

\begin{entradilla}
{\em {\color{introcolor}Windows NT} fu� el primer sistema operativo de Micrsoft con multitarea pre-emptiva all� por 1993. }
\end{entradilla}


Hist�ricamente, la versi�n 3.1 es considerada la primera versi�n oficial
de Windwos. Lanzada en 1992 segu�a siendo una \emph{Shell} sobre DOS,
fue la primera en ofrece un soporte m�s que anecd�tico al modo protegido
introducido por los procesadores 80386. Versiones anteriores de Windows
hab�an explorado este campo, pero los avances no se pod�an considerar
suficientemente representativos. De hecho, fue la version 3.11 liberada
en 1993, la primera que empezaba a parecer algo m�s que un programa DOS
que pintaba ventanas.



En ese mismos a�o, 1993, la versi�n 3.1 de Windows NT tambi�n se
lanzaba. Windows NT era la versi�n profesional del sistema operativo de
Microsoft. El nombre al parecer viene de incremental las letras de
\textbf{VMS} (el sistema operativo de las m�quinas VAX, uno de los m�s
potentes de la �poca), es decir, \textbf{WNT} y de ah� windows NT.


La primera versi�n de Windows NT fue la 3.1, sin embargo fue la version
3.5 liberada en 1994 la mas popular. Windows NT utilizaba un modelo de
memoria plano y un kernel con multitarea pre=emptiva en modo protegido.
Nada que ver con el modelo de memoria segmentado de la otra familia de
windows y la multitarea cooperativa\ldots{} que tanta inestabilidad
causaba a los PCs de la �poca.

Tras esto, se lanz� windows 95 en\ldots{} 1995 con la intenci�n de sacar
m�s partido a las capacidades de los procesadores de la �poca (los
procesadores Pentium ya estaba disponibles). Windows 95, 98 y ME fueron
realmente un desastre, muy inestables, hasta el punto de corromper el
sistema de tal forma que los usuarios de esos sistemas consideraban
natural, reinstalar de cero todo el sistema cada pocos meses.

Finalmente, en 2000, microsoft consigui� aunar los avances en el GUI de
windows 95 y posteriores (aspecto en el que si hubo un gran desarrollo)
y el kernel de NT. Windows 2000 fue liberado, y esa es quiz�s la primera
versi�n realmente estable de los sistemas Windows. Para el tema que nos
ocupa, no tiene sentido avanzar m�s en el tiempo.

\hypertarget{apple}{%
\sectiontext{white}{black}{APPLE}\label{apple}}

Esta secci�n va a ser curiosa y levantar� algunas ampollas. Vamos a
ello. El otro jugador en el campo de juegos de los ordenadores
personales junto con el tandem IBM/Microsoft era Apple. En el albor de
los 80 y 90 los ordenadores Apple no eran, ni mucho menos, tan populares
como los PCs, fundamentalmente por su precio, varios ordenes de magnitud
por encima.

Los ordenadores Apple, a diferencia de los PCs de IBM, dispon�an de HW
especial para imagen y sonido, lo que los hicieron muy populares en ese
�mbito. Sin embargo ahora estamos hablando de el sistema operativo, y en
ese �mbito , dejando a un lado el interfaz de usuario, no era nada del
otro mundo.

Hasta el 91, los OSes de los populares Machintosh se conocen como System
Software N, siendo el \emph{System Software 1} el primero liberado en
1984 y el \emph{System Software 7}. De este �ltimo OS hubo varias
versiones liberadas desde 1991 hasta 1997 en la que la version 7.6 se
libero. Esta �ltima versi�n ya se conoce como \emph{MAC OS 7.6} y es a
partir de ella con la que se cambia el nombre de \emph{System Software}
a \emph{MAC OS}, si bien el n�mero de versi�n se mantiene.

\begin{entradilla}
{\em {\color{introcolor}Apple no utilizo multitare pre-emptiva} en sus sistemas Operativos hasta 1997}
\end{entradilla}



Pues lo cre�is o no, hasta este momento, los sistemas operativos de
Apple todav�a utilizaban multitarea cooperativa\ldots{} en 1997!!!. La
estrategia hay que admitir que era ingeniosa. Hago un GUI muy bonito y
f�cil de utilizar, pero dise�ado para ejecutar una sola aplicaci�n a la
vez\ldots{} de forma que no se note mucho que no era multitarea real.



Hasta que se liber� el proyecto Rhapsody (lo que se acabar�a
convirtiendo en Ma OS X) los sistemas de Apple no dispon�an de
multitarea real. De hecho hasta el lanzamiento de Mac OS X en 2001, no
se produjo la transici�n al kernel Mach que tomaron prestado de NeXT.

\hypertarget{otros-kernels-dignos-de-menciuxf3n}{%
\sectiontext{white}{black}{OTROS KERNELS DIGNOS DE MENCI�N}\label{otros-kernels-dignos-de-menciuxf3n}}

Si bien el mundo de los ordenadores personales estaba totalmente
dominado por Microsoft y Apple en los 80 y 90, y dejando a un lado otros
ordenadores como las m�quina Atari o Amiga que tuvieron una vida m�s
corta, hay un par de sistemas que son dignos de menci�n:

\begin{itemize}
\item
  OS/2. Este sistema operativo desarrollado por IBM (quien tambi�n
  colabor� con Microsoft en el desarrollo de NT) e inicialmente por
  Microsoft. EN 1992 OS/2 ofrec�a un sistema operativo nativo de 32bits,
  siendo la versi�n OS/2 Wrap la m�s popular por su mejorado interaz y
  el buen soporte de applicaciones DOS y WIndows. OS/2 Wrap se lanz� en
  1994
\item
  BeOS. BeOS fue lanzado oficialmente en 1995, como un moderno OS con
  todas las principales caracter�sticas de un sistema operativo moderno,
  todav�a no disponible en ninguno de los sistemas operativos de los que
  hemos hablado. Lamentablemente, este sistema operativo nunca lleg� a
  ser popular, si bien, todav�a hoy en d�a hay versiones disponibles con
  las que probarlo. Especial menci�n para el proyecto libre Haiku que
  todav�a est� en activo.
\item
  NextStep: Si bien, el sistema operativo NextStep as� como las m�quinas
  Next caen un poco fuera de la categor�a de ordenador personal (igual
  que las grandes m�quinas Unix o VAX que ofrec�an multitarea real desde
  hacia varios a�os), merece la pena mencionarla ya que fue el OS
  seleccionado por Apple para su transici�n final a la liga de los
  mayores. NextStep desarrollo su propio kernel basado en Mach, un
  microkernel desarrollado en la universidad Carnegie Mellon
\item
  Minix: Es un sistema operativo con arquitectura de microkernel. Creado
  inicialmente para fines educativos por el Profesor Andrew S. Tanenbaum
  una versi�n reducida de su c�digo fuente fue publicada en un libro. El
  sistema inicialmente desarrollado para plataformas intel de 16 bits se
  desarroll� hasta su versi�n 3 y ha permanecido activo hasta 2018. Al
  contrario que Linux que utiliza un dise�o monol�tico, Minix sigue un
  modelo de nanokernel como Mach. Esta discrepancia fue el origen de un
  conocido debate entre
  \href{https://en.wikipedia.org/wiki/Tanenbaum\%E2\%80\%93Torvalds_debate}{Tanenbaum
  y Torvalds}.
\end{itemize}

\hypertarget{linux}{%
\sectiontext{white}{black}{LINUX}\label{linux}}

El kernel Linux se liber� en el verano de 1991. Desde su primera versi�n
fue concebido como un sistema operativo de 32 bits ejecut�ndose en modo
protegido, con multitarea pre-emptiva, espacio de direccionamiento plano
y soporte nativo para redes TCP/IP.

En 1992 las primeras distribuciones de Linux ya estaban disponibles, si
bien es cierto que eran un poco complicadas de utilizar (por decirlo de
una forma suave), pero a partir de aquel momento nuevas distrbuciones de
Linux no dejaron de aparecer.

Debian en 1993 aunque la primera versi�n estable es de 1996, si bien
todos sabemos lo que significa estable para Debian. Slackware tambi�n
vio la luz en 1993, pero no fue hasta la versi�n 4.2 de Red-Hat que
GNU/Linux empez� a ser una opci�n viable para usuarios avanzados.

Dejando a un lado la complejidad de la instalaci�n y la deficiencia de
soporte HW que era uno de los principales motivos de complejidad en la
instalaci�n, el kernel linux fue de los primeros en ofrecer soporte real
para sistemas de 32bits. Mucho antes que cualquiera de los otros sistema
operativos desarrollados por compa��as privadas.

Cabe rese�ar que aunque GNU/Linux no estaba orientado a usuarios
dom�sticos inicialmente, siempre dispuso de entornos gr�ficos con
algunas capacidades que muchos sistemas comerciales han empezado a
soportar muy recientemente\ldots{} Pero como dec�a Conan, esa es otra
historia.


\hypertarget{conclusiones}{%
\sectiontext{white}{black}{CONCLUSIONES}\label{conclusiones}}

Si bien, algunas compa��as se han elegido como abanderadas de la
innovaci�n, la realidad es bien distinta y como hemos podido ver, cuanto
m�s triunfadora es una impresa menos innovadora es. Es algo
tremendamente curioso. IBM, Next o Be Inc realmente innovaron y
contribuyeron al progreso de los sistemas operativos. Dejando a IBM a un
lado, que nunca hizo negocio de sus sistemas operativos, las otras
compa��as han terminando desapareciendo.

NeXt consigui� transcender a trav�s de OpenStep y su implementaci�n
libre GNU/Step, y de los sistemas operativos de Apple\ldots{} Si bien
llev� casi una d�cada en que los avances de esa compa��a se reflejaran
en Mac OS. SI bien, a nivel de sistema operativo NeXt se aprovech� de
los desarrollos del mundo UNIX de a�os anteriores y de el microkernel
March.

BeOS por su parte tambi�n a transcendido gracias al mundo del software
libre con proyectos como Haiku.

Finalmente, resulta curioso como la mayor innovaci�n en esta �rea ha
estado dominada for el kernel Linux desde su creaci�n a principios de
los 90. Linuxha sido hist�ricamente uno de los kernel adoptando nuevas
tecnolog�as y conceptos m�s r�pidamente. Esto es gracias sobre todo a su
c�digo abierto que abre una puerta a investigadores y empresas para
implementar nuevas ideas sin necesidad de partir de cero.

As� que, independientemente de lo que algunas empresas pretenden dar a
entender, de la furiosa defensa de diversos fanboys y de la propaganda
barata, esas empresas que se presentan como avanzadilla del progreso, en
muy raras ocasiones lo han sido y cada vez menos lo volver�n a ser.

%\columnbreak

\textbf{Cronolog�a de los principales kernels multitarea desde lo 80s}
{\footnotesize
\begin{verbatim}
+- 1984 : System Software 1 
|  1985 : Windows 1.0
+- 1987 : System Software 5
|         Windows 2.0
|         Minix 1.0 (8088)
+- 1988 : System Software 6
|         NextStep 0.9 (m68k)
+- 1989 : NextStep 1.0 (m68k)
+- 1990 : WIndows 3.0
|         NextStep 2.0 (m68k) 
+- 1991 : Kernel Linux
|         System Software 7
|         Minix 1.5 (m68K, SPARC, 386)
|         NextStep 2.1 (m68k)
+- 1992 : Windows 3.1
|         OS/2 2.0
|         NextStep 3.0 (m68k)
+- 1993 : Windows 3.11
|         Windows NT 3.1
|         Slackware 1.0 (Kernel 0.99.11)
|         NextStep 3.1-3.2 (m68k, i386)
+- 1994 : Debian 0.91 (Linux 0.99)
|         Windows NT 3.5
|         OS/2 Wrap
|         Slackware 1.2-2.1 (Kernel 1.0.8-1.1.59)
+- 1995 : Windows 95
|         RedHat 1/2.0 (Kernel 1.2.8-1.2.13)
|         Debian 0.93 (Linux 1.2.13)
|         Slackware 2.2-3.0 (Kernel 1.2.1-1.2.13)
|         BeOS DR1
|         NextStep 3.3 (m68k, i386, SPARC, PA-RISC)
+- 1996 : OS/2 Wrap 4
|         Windows NT 40
|         Debian 1.1 (Buz) 1.2 (Rex) (Kernel 2.0)
|         RedHat 4 (Kernel 2.0.18)
|         Slackware 3.1 (Kernel 2.0.0)
|         NextStep 4.0 (m68k, i386, SPARC)
+- 1997 : Systems Software 7,6/Mac Os 7 
|         Mac Os 8 (Multitarea cooperativa)
|         Debian 1.3 (Bo) (Kernel 2.0.33)
|         Red Hat 4.2-5.0 (Kernel 2.0.27-2.0.32)
|         Slackware 3.1-3.4 (Kernel 2.0.29-2.0.30)
|         Minix 2.0
|         NextStep 4.1 (m68k, i386, SPARC)
+- 1998 : Windows 98
|         Mac OS 8.1-8.5.1
|         BeOs R3
|         Debian 2.0 (Hamm) (Kernel 2.0.34)
|         Red Hat 5.2 (Kernel 2.0.36)
|         Slackware 3.5-3.6 (Kernel (2.0.34-2.0.35)
+- 1999 : Mac Os 8.6
|         Mac Os 9
|         RedHat 6 (Kernel 2.2.5)
|         Debian 2.1 (Slink) (Kernel 2.0.34-38)
|         Slackware 3.9-7.0 (Kernel 2.0.37-2.2.13)
+- 2000 : Windows 2000 (NT 5.0)
|         Windows 98SE
|         Windows ME
|         Mac Os X
|         Debian 2.2 (Potato) (Kernel 2.2.19)
|         Red Hat 6.2 - 7.0 (Kernel 2.2.12-2.2.16)
|         Slackware 7.1 (Kernel 2.2.16)
\end{verbatim}
}

\end{multicols}
\clearpage
\pagebreak

\hypertarget{pildoras:alloca1}{}\label{pildoras:alloca}
\pagestyle{pvc}

\rput(7.9,0.0){\resizebox{!}{16cm}{{\epsfbox{images/header-alloca.eps}}}}

\psset{fillstyle=solid}
\psframe[fillcolor=black,opacity=0.6](2,-4.5)(17,0)




% -------------------------------------------------
% Cabecera
\begin{flushright}


{\color{orange}\mtitle{14cm}{Pildoras de vitamina C: alloca}}

\msubtitle{10cm}{Reserva memoria en la pila super r�pido}

{\sf\color{white}{ por Richi C. Poweri}}

%{\psset{linecolor=white,linestyle=dotted}\psline(-12,0)}
\end{flushright}

\vspace{2mm}
% -------------------------------------------------



%\lstset{language=C,frame=tb,framesep=5pt,basicstyle=\footnotesize}
\lstset{language=C,frame=tb,framesep=5pt,basicstyle=\scriptsize}




\intro{introcolor}{L}{a funci�n \texttt{alloca} no es muy conocida\ldots{}
de hecho su uso no est� recomendado puesto que su implementaci�n depende
de la m�quina y compilador que usemos. Sin embargo, es una funci�n muy
potente y merece la pena conocerla ya que, cuando tenga sentido
utilizarla puede ser un arma poderosa.}

\begin{multicols}{2}

La funci�n \verb!alloca! simplemente reserva espacio
en la pila, o en otras palabras, modifica el puntero de pila. Esa es la
raz�n por la que es tan eficiente, y es la misma raz�n por la que
depende de la plataforma. La otra gran ventaja de
\verb!alloca! es que, puesto que solo modifica el
puntero de pila, la memoria se libera autom�ticamente al salir de la
funci�n.

\hypertarget{que-significa-reservar-memoria}{%
\sectiontext{white}{black}{�QU� SIGNIFICA RESERVAR MEMORIA?}\label{que-significa-reservar-memoria}}

A no ser que ya sepas de que va esto de
\verb!alloca!, y en ese caso seguro que ya no est�s
leyendo esto, quiz�s no tengas muy claro como funciona el proceso de
reservar memoria. Aqu� tienes un curso intensivo :)

Cuando se lanza un proceso, el sistema operativo reserva varios bloques
de memoria para su ejecuci�n. Por una parte un bloque de memoria
contendr� el c�digo, otro los datos de solo lectura, otro las variables
globales y otro la pila. Los dos primeros bloques son, en principio
fijos, mientras que el tama�o de los dos �ltimos pueden variar durante
la ejecuci�n del programa.

El bloque de datos se puede modificar con las llamadas al sistema
\verb!brk!, mientras que la pila la modifica el
programa a voluntad, b�sicamente modificando el puntero de pila.

\begin{entradilla}
{\em El {\color{introcolor}tama�o de la pila puede variar din�micamente} durante la ejecuci�n de un programa.}
\end{entradilla}

En general, la pila crece hacia las direcciones bajas (hay algunos
procesadores en los que va al rev�s) y el segmento de datos crece hacia
las direcciones altas. Por esa raz�n, la pila se inicia en la zona alta
de la memoria, y el segmento de datos en la zona baja. De esta forma
tenemos el mayor hueco posible entre ambas, para reservar memoria, ya
sea en la pila o en el segmento de datos.

\emph{NOTA:En realidad las librer�as din�micas se mapean en medio y
medio de esa zona libre, pero eso no tiene relevancia para el tema que
nos ocupa.}

El mapa de memoria gen�rico de un proceso ser�a algo como esto:

\begin{lstlisting}
  +--------------------+
  | Pila               | | Crece hacia abajo
  +--------------------+ v
  |                    |
  ~ Memoria Libre      |
  |                    |
  +--------------------+ ^
  | Segmento de Datos  | | Crece hacia arriba
  +--------------------+
  | Codigo + rodata    | rodata -> Datos solo-lectura
  +--------------------+
\end{lstlisting}

\hypertarget{el-stack-frame}{%
\sectiontext{white}{black}{EL STACK FRAME}\label{el-stack-frame}}

Si eres uno de nuestros asiduos lectores seguro que ya sabes de que va
esto del \emph{Stack Frame}, pero para los que reci�n os incorpor�is,
aqu� ten�is una r�pida introducci�n.

Cada vez que ejecutamos una funci�n, el compilador crea para nosotros lo
que se conoce como un marco de pila (\emph{stack frame} en ingl�s). Este
\emph{stack frame} (voy a usar el t�rmino ingl�s, por que el espa�ol la
verdad que da miedito), comienza con la direcci�n de retorno, que no es
otra que la direcci�n de la siguiente instrucci�n a la que contiene el
salto a la funci�n. Dependiendo del procesador esto sucede de distintas
formas, por ejemplo, los x86 utilizan la instrucci�n
\verb!call! que autom�ticamente a�ade la direcci�n de
retorno a la pila. Otros procesadores almacenan la direcci�n en un
registro especial y depende del programador hacer lo correcto con �l.

El compilar a�ade a todas las funciones (a no ser que le digamos lo
contrario), lo que se conoce como el pr�logo de la funci�n, que no es
m�s que un par de l�neas de c�digo para crear el stack frame (si, yo
considero la direcci�n de retorno parte del \emph{stack frame}, pero eso
es discutible). Este pr�logo, normalmente almacena el valor del registro
de pila (que apunta a la siguiente posici�n disponible en la pila) en un
registro auxiliar que se suele conocer como registro base y que para los
procesadores intel ha sido tradicionalmente el
\verb!BP! (\emph{Base Pointer} ). A continuaci�n
modifica el registro de pila para reservar espacio para las variables
locales.

El registro base \verb!RBP! marca el comienzo de la
zona de la pila asociada a esta funci�n. Por esa raz�n debemos
almacenarlo en alg�n sitio antes de modificarlo\ldots{} bueno, la raz�n
es que si nuestra funci�n ha sido llamada por otra funci�n, al entrar en
ella, \verb!RBP! apunta al \emph{stack frame} de la
funci�n que llama, y cuando terminemos de ejecutar nuestra funci�n
queremos restaurarlo, para que al retornar la funci�n llamante pueda
continuar su ejecuci�n\ldots{} es decir, encuentre sus variables locales
y su direcci�n de retorno.

El c�digo del pr�logo es normalmente algo como esto:

\begin{lstlisting}
push RBP
mov  RBP, RSP
sub  RSP, ESPACIO_VARIABLES_LOCALES
\end{lstlisting}

Lo que nos deja la pila de esta forma

\begin{lstlisting}
 | RET  | 
 | RBP  |
 |      | <--- RBP
 ~ ...  ~
 |      | <--- RSP = RBP - ESPACIO_VARIABLES
\end{lstlisting}

Como pod�is ver, las variables locales se reservan en la pila,
simplemente haciendo hueco (cambiando el valor del puntero de pila). De
esta forma, podemos acceder a las variables locales de la funci�n, con
simples deplazamientos en el registro base (en realidad utilizando
direccionamiento indexado sobre �l).

Este c�digo va acompa�ado de un epilogo, al final del la funci�n, que
simplemente hace lo siguiente:

\begin{lstlisting}
mov RSP, EBP
pop RBP
ret
\end{lstlisting}

En otras palabras, restauramos nuestro puntero de pila al valor que
ten�a justo al entrar en la funci�n, y restauramos el registro base con
el \emph{stack frame} anterior que hab�amos almacenado en la
pila\ldots{} Tras eso, el valor al que apunta
\verb!RSP! es la direcci�n de retorno, con lo que
\verb!ret! retornar� al lugar adecuado.

\begin{entradilla}
{\em El {\color{introcolor}c�digo que genera el Stack Frame} de una funci�n se conoce como prologo}
\end{entradilla}

Los procesadores intel incorporan instrucciones para hacer estas dos
cosas de forma autom�tica. Las instrucciones
\verb!enter! y \verb!leave!
respectivamente ejecutan el pr�logo y epilogo de la funci�n. En general,
\verb!enter! no se utiliza ya que es menos eficiente
que el c�digo que os mostramos m�s arriba, as� que raramente lo ver�is
en c�digo real.

Esta sencilla estructura nos permite anidar f�cilmente funciones y es la
base sobre la que se pueden construir funciones recursivas.

\hypertarget{alloca-again}{%
\sectiontext{white}{black}{\texorpdfstring{\texttt{alloca}
AGAIN}{alloca again}}\label{alloca-again}}

Pues, despu�s de todo este rollo, ya os pod�is imaginar que es lo que
hace \verb!alloca!\ldots{} simplemente, modifica el
puntero de pila, para hacer m�s hueco. Si tomamos como ejemplo el
diagrama de pila anterior y suponemos que ejecutamos la funci�n
\verb!alloca (128)!, se nos quedar�a tal que as�:

\begin{lstlisting}
 | RET  | 
 | RBP  |
 |      | <--- RBP
 ~ ...  ~
 |      | <--- Termina ESPACIO_VARIABLES
 ~ 128  ~
 | bytes|
 |      | <-- RSP
\end{lstlisting}

Como pod�is ver, cuando se ejecuta el ep�logo de la funci�n, el puntero
de pila, simplemente se resetea (se sobre escribe con el valor de
\verb!RBP!) y como resultado, toda la memoria est�
ahora disponible de nuevo para utilizarla con otro
\verb!alloca! u otra \emph{Stack Frame}.

Esta es la raz�n por la que \verb!alloca! es tan
eficiente, b�sicamente hace una suma o una resta\ldots. Para
asegurarnos, veamos que genera \verb!gcc! para la
siguiente sencilla funci�n:

\begin{lstlisting}[language=C]
void* func () {
  char *ptr;
  ptr = alloca (1024);
  memset (ptr, 0, 1024);
  return ptr;
}
int main () { func (); return 0;}
\end{lstlisting}

{\em{Nota: El {\verb!memset!} es necesario para que al
optimizar el c�digo el compilador se de cuenta que hacemos algo con
{\verb!ptr!} y no lo elimine}.}

Si compilamos nuestro programa con la optimizaci�n `-O1':

\begin{lstlisting}
gcc -O1 -o alloca-test alloca-test.c
\end{lstlisting}

La funci�n \verb!func! se quedar�a m�s o menos as�:

\begin{lstlisting}
<func>:
    push   rbp
    xor    eax,eax
    mov    ecx,0x80
    mov    rbp,rsp
    sub    rsp,0x410
    lea    rdx,[rsp+0xf]
    and    rdx,0xfffffffffffffff0
    mov    rdi,rdx
    rep stos QWORD PTR es:[rdi],rax
    mov    rax,rdx
    leave
    ret
\end{lstlisting}

\hypertarget{analizando-el-cuxf3digo-generado}{%
\sectiontext{white}{black}{ANALIZANDO EL C�DIGO GENERADO}\label{analizando-el-cuxf3digo-generado}}

Si reordenamos un poco el c�digo\ldots{}

\begin{lstlisting}
    push   rbp
    mov    rbp,rsp
    sub    rsp,0x410
    
    lea    rdx,[rsp+0xf]
    and    rdx,0xfffffffffffffff0
        
    xor    eax,eax
    mov    ecx,0x80     
    mov    rdi,rdx
    rep stos QWORD PTR es:[rdi],rax
    mov    rax,rdx
    
    leave
    ret
\end{lstlisting}

Ahora podemos ver claramente el pr�logo de la funci�n y como reservamos
espacio para los 1024 bytes reservados con
\verb!alloca! y 16 bytes m�s para almacenar el
puntero\ldots{} que para una arquitectura de 64 bits debe estar alineado
a 16 bytes\ldots{} aunque solo usemos 8.

Las dos instrucciones siguientes se encargan de alinear el puntero en la
zona de la pila reservada. Veamos como est� la pila justo despu�s del
pr�logo, para entender que est� pasando:

\begin{lstlisting}
  | RET     |
  | RBP     | <- RBP
  |         |
  ~ ....    ~ <= 0x400 bytes de alloca
  |         | <= 0x10 bytes para el puntero
  |         | <- RSP = RBP - 0x410
\end{lstlisting}

Lo que el compilador ha hecho es reservar \verb!0x10!
bytes para almacenar nuestro puntero que realmente solo necesita 8 bytes
en lugar de 16, para a continuaci�n calcular la direcci�n de entre esas
16 que est� alineada a 16 bytes. Para ello lo que hace es lo siguiente:
sumar 15 y luego borrar los �ltimos 4 bytes (o dividirlo por 16 si
prefer�s)\ldots. que efectivamente significa alinear a 16 bytes (4 bits
-\textgreater{} 16 valores).

\begin{entradilla}
{\em Para {\color{introcolor}m�quinas {\verb!x86_64!} utilizan 16 bytes} para el alineamiento de memoria en la pila}
\end{entradilla}


Imaginad que el valor de \verb!RSP! despu�s de
reservar la memoria es por ejemplo \verb!0x01412!, lo
que significa que nuestro buffer est� en
\verb!0x01432!\ldots{} por lo que tenemos que elegir
una direcci�n entre esas dos que est� alineada a 16 bytes\ldots{} o lo
que es lo mismo que los �ltimos 4 bits sean cero.

Podr�amos pensar en simplemente borrarlos\ldots{} eso nos dar�a como
resultado\verb!0x01410! que est� fuera de nuestro
rango (fuera de la pila), as� que lo que hacemos es sumar 15\ldots{} de
forma que nos aseguramos que vamos a quedarnos en el rango
disponible\ldots{} pero sin tocar el buffer:

\begin{lstlisting}
  0x001412 + 0xf = 0x001421
\end{lstlisting}

Si ahora borramos los �ltimos 4 bits, obtenemos la direcci�n
\verb!0x001420!, que se encuentra en el rango
adecuado de nuestra pila, nos da espacio suficiente para almacenar
nuestro puntero de 8 bytes y adem�s est� alineado con la pila a 16
bytes.


El resto del programa no es m�s que la versi�n \emph{inline} de
\verb!memset!\ldots{} El �nico comentario es que
\verb!stos! en el c�digo generado, usa registros de
64 bytes, por lo que el contador en \verb!rcx! es
solo 0x80 (0x80 * 8 = 0x400 = 1024)\ldots{} vamos que copia de 8 en 8
bytes.

\hypertarget{stack-vs-heap.-fight}{%
\sectiontext{white}{black}{STACK vs HEAP\ldots. FIGHT!}\label{stack-vs-heap.-fight}}

Probablemente siempre has escuchado que la forma de reservar memoria en
C es utilizando las funciones \verb!malloc! y
\verb!free!. Si bien eso es lo normal, como hemos
visto el proceso de reservar memoria es un poco m�s complicado que eso.
As� que vamos a describir r�pidamente como funcionan esas dos funciones.

Ya hemos hablado brevemente del segmento de datos y como modificarlo.
Esa es la forma m�s sencilla de implementar un sistema de gesti�n de
memoria din�mico. De hecho es un buen ejercicio. Si lo hac�is, ver�is
que hay muchos problemas que tener en cuenta y que no es nada f�cil
hacer uno de estos sistemas que funcione decentemente en cualquier
situaci�n.

Esa es la raz�n por la que la librer�a C de GNU no implementa una �nica
estrategia para reservar memoria aunque nosotros, como programadores, no
lo vemos. Esto lo podemos dejar para otro art�culo en el futuro, ya que
el funcionamiento interno del sistema de gesti�n de memoria es
fundamental para entender ciertos ataques como los \emph{heap
overflows}, \emph{out by one},\ldots{} relacionados con el \emph{heap}.

Para terminar, deciros que existe otra forma m�s de reservar memoria (de
hecho la implementaci�n de \verb!malloc! de la
librer�a est�ndar tambi�n la utiliza), y no es otra que la llamada al
sistema \verb!mmap!. \verb!mmap!
nos permite reservar bloques de memoria, con permisos espec�ficos e
incluso, mapearlos en direcciones espec�ficas\ldots{} todo esto con
algunas limitaciones\ldots. Parece una aut�ntica pasada no? Os
preguntar�is porque no usamos \verb!mmap! y ya
est�\ldots{} bueno, hay distintas razones, pero lo fundamental es que
\verb!mmap! trabaja sobre p�ginas de memoria, no
sobre direcciones. Lo que significa que si necesitamos 8 bytes,
tendremos que reservar 4K (o el tama�o de p�gina que sea que utilice el
OS).

Como pod�is ver, no existen las panaceas. Cada t�cnica tiene sus pros y
sus contras y por ello debemos conocerlas todas para poder elegir la m�s
adecuada.

\hypertarget{problemas-con-alloca}{%
\sectiontext{white}{black}{\texorpdfstring{PROBLEMAS CON \texttt{alloca}}{Problemas con alloca}}\label{problemas-con-alloca}}

Para terminar este art�culo, vamos a discutir brevemente los problemas
que os pod�is encontrar con \verb!alloca!. La funci�n
no tiene ning�n problema, y usada con cuidado puede hacer que nuestro
programa sea mucho m�s eficiente, pero si nos despistamos, puede hacer
que el programa falle de formas poco obvias.

Lo primero que tenemos que entender es que, en contra de lo que podr�a
parecer, \verb!alloca! no es una versi�n optimizada
de \verb!malloc!, aunque se use con ese objetivo. La
principal consecuencia de esto es que objetos creados con
\verb!alloca! no pueden ser utilizados fuera de la
funci�n como ocurre con la memoria reservada con
\verb!malloc!.

\begin{entradilla}
{\em La memoria reservada con \texttt{alloca} se encuentra en la pila y por tanto se {\color{introcolor}libera autom�ticamente} al abandonar la funci�n.}
\end{entradilla}

Lo segundo es que la memoria reservada con
\verb!alloca! no hay que liberarla como ocurre con la
memoria reservada con \verb!malloc!. De hecho,
ejecutar \verb!free! sobre un bloque de memoria
reservado por \verb!alloca! puede ser fatal. La
memoria reservada con \verb!alloca! se libera de
forma autom�tica al dejar la funci�n (como hemos visto antes).

\verb!alloca! puede producir desbordamientos de pila,
ya que la implementaci�n muchas veces no hace ning�n tipo de
comprobaci�n sobre los par�metros. De esta forma, si pasamos un valor
negativo, por ejemplo, podr�amos sobre escribir el \emph{stack frame} de
la funci�n y hacer que el programa se estrelle\ldots{} Probad a escribir
\verb!alloca(-1024);! en nuestro programa de ejemplo
y observad que c�digo se genera.

De la misma forma, utilizar valores muy grandes nos puede llevar fuera
del bloque de memoria asociado a la pila y provocar un fallo de
segmentaci�n\ldots. S�, la pila no puede crecer indeterminadamente en el
hueco libre en medio del espacio de direccionamiento del proceso,
especialmente en un sistema multi-proceso. El segmento de pila tiene un
tama�o y si nos pasamos de �l se producir� una violaci�n de segmento.
Para ver donde se encuentra la pila para un determinado proceso pod�is
mirar el mapa de memoria del proceso en
\verb!/proc/PID/maps! donde
\verb!PID! es el identificador del proceso.

Otra cosa a tener en cuenta es que utilizar
\verb!alloca! solo tiene sentido cuando no conocemos
el tama�o del bloque a reservar, ya que de lo contrario simplemente
deber�amos declarar una variable local del tama�o adecuado. La principal
ventaja de declarar la variable es que el compilador podr� hacer un
mejor trabajo detectando posible bugs.

\hypertarget{siendo-portables}{%
\sectiontext{white}{black}{SIENDO PORTABLES}\label{siendo-portables}}

A partir de la versi�n 99 de C, lo que se conoce como C99, el est�ndar
incluye las llamadas matrices de longitud variable o VLA, de sus siglas
en ingl�s \emph{Variable Length Array}. B�sicamente funcionan como
\verb!alloca!, pero de forma est�ndar y a�adiendo
algo m�s de seguridad. Una vez m�s, al igual que sucede con
\verb!alloca!, esto solo tiene sentido cuando debemos
reservar un bloque de memoria de longitud variable (de ah� su
nombre)\ldots{} de lo contrario, simplemente declaramos la variable como
global y ya acabar� en la pila directamente.

\begin{entradilla}
{\em La {\color{introcolor}funci�n \texttt{alloca} no es portable}, pero C99 introdujo los VLAs.}
\end{entradilla}

El siguiente c�digo muestra como utilizar esta construcci�n introducida
en C99 aunque gcc ya ofrec�a esta funcionalidad como una extensi�n antes
de que fuera incorporada al est�ndar:

\begin{lstlisting}[language=C]
int func (int n) {
  char ptr[n];
  memset (ptr, 0, n);
  return 0;
}
\end{lstlisting}

Como pod�is ver, simplemente declaramos el array que necesitamos, pero
ahora podemos utilizar una variable para indicar su longitud. Esto no
era posible en versiones anteriores de C y este tipo de c�digo requer�a
del uso de \verb!malloc!.

Si le ech�is un ojo al ensamblador generado, es bastante m�s complicado
que simplemente que el generado para \verb!alloca!,
pero ese suele ser el precio a pagar por la portabilidad.

Como nota final, comentaros que los \verb!VLA!s no se
utilizan en el kernel Linux desde la version 4.20 debido a la baja
calidad del c�digo generado para tama�os peque�os del VLA.

\hypertarget{conclusiones}{%
\sectiontext{white}{black}{CONCLUSIONES}\label{conclusiones}}

Bueno, hasta aqu� esta peque�a introducci�n a la funci�n
\verb!alloca! y una no tan peque�a digresi�n sobre la
gesti�n de memoria en procesos. Conocer los detalles de como funciona la
gesti�n de memoria nos hace m�s f�cil entender lo que hacemos en cada
momento y reducir los bugs en nuestros programas.

\verb!alloca!, si bien , no deber�a ser tu primera
opci�n para manejar la memoria de tu programa, puede ser una
optimizaci�n muy interesante en algunos casos cuando es usada con
cuidado.

\end{multicols}

%\includegraphics[width=.9\linewidth]{images/promo/promo01.eps}
\rput(7.9,-3.8){\resizebox{!}{9.2cm}{{\epsfbox{images/promo/promo01.eps}}}}

\hypertarget{sec:criptografia1}{}\label{sec:criptografia}
\pagestyle{seguridad}

\rput(7.9,0.0){\resizebox{!}{16cm}{{\epsfbox{images/header-seguridad.eps}}}}

\psset{fillstyle=solid}
\psframe[fillcolor=black,opacity=0.6](2,-4.5)(17,0)




% -------------------------------------------------
% Cabecera
\begin{flushright}


{\color{green}\mtitle{10cm}{Introducci�n a la Criptograf�a}}

\msubtitle{8cm}{Seguridad para meros mortales}

{\sf\color{white}{ por Andr�s ``Andy'' Pajaker}}

%{\psset{linecolor=black,linestyle=dotted}\psline(-12,0)}
\end{flushright}

\vspace{2mm}
% -------------------------------------------------



%\lstset{language=C,frame=tb,framesep=5pt,basicstyle=\footnotesize}
\lstset{language=C,frame=tb,framesep=5pt,basicstyle=\scriptsize}



\intro{introcolor}{L}{o sepas o no, la criptograf�a forma parte de nuestro d�a a d�a desde
hace ya algunos a�os, si bien, la mayor�a de nosotros no sabemos muy
bien como funciona o que es lo que se puede o no se puede hacer. Muchas
compa��as tratan de convencernos de lo seguras que son, simplemente
afirmando que utilizan \emph{Criptograf�a fuerte} o \emph{Extremo a
Extremo}\ldots{} Pero realmente sabes de lo que est�n hablando?\ldots{}}

\begin{multicols}{2}

La criptograf�a es un elemento esencial de las comunicaciones actuales.
Pr�cticamente cualquier acceso a un servicio online est� cifrado. Los
datos almacenados en nuestros dispositivos tambi�n est�n cifrados. Los
datos almacenados en la nube\ldots{} Pero eso no significa que estemos
seguros.

En este art�culo os vamos a explicar como funciona la criptograf�a y
como se utiliza en las comunicaciones, para que pod�is entender mejor
los riesgos a los que os enfrent�is cada d�a. Prometemos que el art�culo
va a ser muy sencillo para que cualquiera pueda entenderlo y al terminar
de leerlo, tenga una idea muy clara de como funciona todo esto.

\hypertarget{servicios-criptogruxe1ficos}{%
\sectiontext{white}{black}{SERVICIOS CRIPTOGR�FICOS}\label{servicios-criptogruxe1ficos}}

Antes de entender como utilizamos criptograf�a en los servicios que
utilizamos cada d�a, debemos entender los fundamentos b�sicos. Es decir,
como funciona, no como crear nuestros propios sistemas
criptogr�ficos\ldots{} de eso ya se ocupar�n otros colegas de Occam's.

Cuando pensamos en criptograf�a lo primero que se nos viene a la mente
son mensajes secretos. Y si, esa es, quiz�s, la primera aplicaci�n de la
criptograf�a y es, adem�s, la operaci�n fundamental en la que se basan
las otras aplicaciones. En general, la criptograf�a busca ofrecernos las
siguientes aplicaciones:

\begin{itemize}

\item
  Privacidad. La privacidad es la propiedad de permite que cierta
  informaci�n no pueda ser accedida por personas que no queramos. Esta
  es la aplicaci�n m�s com�n de la criptograf�a\ldots. cifrar datos de
  forma que si no sabes la clave no podr�s acceder a ellos.
\item
  Autenticaci�n. Esta aplicaci�n busca formas para asegurar que una
  determinada persona (o sistema) es realmente quien dice ser. No
  queremos que nadie se haga pasar por nosotros. Utilizamos los mismos
  algoritmos, pero de una forma un tanto diferente.
\item
  Integridad. Esta aplicaci�n busca asegurar que ciertos datos no han
  sido modificados, al menos no por alguien no autorizado. Imaginad que
  os mandan una factura por un servicio que hab�is contratado, pero
  alguien la intercepta, cambia el n�mero de cuenta en el que ten�is que
  pagar y luego os la env�a. Los servicios de integridad nos permiten
  saber si esa factura que recibimos ha sido modificada por un criminal
  o no.
\item
  No repudiaci�n de origen. Esta aplicaci�n con un nombre tan curioso
  busca asegurar que las acciones que realiza una cierta persona o
  sistema las ha realizado esa cierta persona o sistema. Conceptualmente
  es muy parecido a la autenticaci�n , pero la no repudiaci�n de origen
  es un concepto un poco m�s amplio, aunque probablemente tendremos m�s
  que suficiente con el resto de aplicaciones que hemos mencionado.
\end{itemize}

Ahora es un buen momento para explicaros los dos fundamentos
criptogr�ficos en los que se fundamentan todos estos servicios. Vamos
all�.

\hypertarget{fundamentos-criptogruxe1ficos}{%
\sectiontext{white}{black}{FUNDAMENTOS CRIPTOGR�FICOS}\label{fundamentos-criptogruxe1ficos}}

La mayor�a de las aplicaciones que hemos introducido en la secci�n
anterior se basan en dos conceptos bastante sencillos: cifrado y
hashing.

El cifrado como todos sab�is consiste en transformar un determinado
mensaje en otro utilizando alguna informaci�n secreta de forma que, si
conocemos esa informaci�n secreta podemos obtener el mensaje inicial muy
f�cilmente, pero si no la conocemos (la informaci�n secreta o clave), es
muy, muy dif�cil obtener el mensaje original.

\begin{entradilla}
{\em Los algoritmos criptogr�fico utilizan {\color{introcolor}funciones trampa} que resultan muy dif�ciles de invertir.}
\end{entradilla}

Matem�ticamente esto se conoce como funci�n trampa. Se trata de
funciones que son muy f�ciles de evaluar en una direcci�n, pero que
calcular su inversa es muy dif�cil. Como dijimos que no �bamos a
meternos en detalles escabrosos, aqu� lo vamos a dejar.

Un ejemplo de cifrado muy b�sico es el conocido como cifrado de C�sar,
atribuido a este general Romano. Este c�digo, simplemente consiste en
asignar un valor num�rico a cada letra y generar el mensaje cifrado
sumando un determinado n�mero a la letra original. El n�mero que sumamos
es la clave. Imaginad que asignamos a la letra A el valor 1 y a la letra
Z el 26. El 0 representar� el espacio.

\end{multicols}

% TODO: Center this (Change to image)

{\color{introcolor}{\hrule}}
\begin{verbatim}
        A B C D E F G H I J  K  L  M  N  �  O  P  Q  R  S  T  U  V  W  X  Y  Z
        1 2 3 4 5 6 7 8 9 10 11 12 13 14 15 16 17 18 19 20 21 22 23 24 25 26 27
\end{verbatim}
{\color{introcolor}{\hrule height 1.5pt}}
\textbf{Equivalencia caracteres n�meros para cifrado C�sar}
{\color{introcolor}{\hrule height 1.5pt}}


\begin{multicols}{2}
Ahora podemos convertir cualquier mensaje en una secuencia de n�meros.
Por ejemplo:

\begin{verbatim}
H  O  L  A   M  U  N  D  O
8  16 12 1 0 13 22 14 4  16
\end{verbatim}

Si ahora le sumamos dos a cada uno de esos n�mero obtenemos:

\begin{verbatim}
H  O  L  A   M  U  N  D  O
8  16 12 1 0 13 22 14 4  16
10 18 14 3 2 15 24 16 6  18
\end{verbatim}

Y utilizando nuestra tabla que asocia letras y n�meros obtenemos el
mensaje:

\begin{verbatim}
H  O  L  A   M  U  N  D  O
8  16 12 1 0 13 22 14 4  16
10 18 14 3 2 15 24 16 6  18
J  Q  N  C B �  W  O  F  Q
\end{verbatim}

As� que \emph{HOLA MUNDO} se ha convertido en \emph{JQNCB�WOFQ}\ldots{}
Hemos cifrado nuestro mensaje y si no sabemos que tenemos que restarle
dos a cada uno de los valores num�ricos del texto cifrado no podemos
recuperar el mensaje original. El c�digo de C�sar es muy sencillo y
tambi�n muy malo, tiene muchos problemas como que es susceptible de
ataques de fuerza bruta o puede ser cryptoanalizado\ldots{} Pero ya
volveremos sobre esto m�s adelante.

\begin{entradilla}
{\em El sistema de cifrado sim�trico m�s popular es {\color{introcolor}AES}}
\end{entradilla}


El algoritmo de cifrado sim�trico m�s popular (y m�s seguro en 2024) es
AES (\emph{Advanced Encryption Standar} o \emph{Est�ndar de Cifrado
Avanzado}).

\hypertarget{fundamentos-criptogruxe1ficos-2.-cifrado-asimuxe9trico}{%
\sectiontext{white}{black}{CIFRADO SIM�TRICO Y ASIMETRICO}\label{fundamentos-criptogruxe1ficos-2.-cifrado-asimuxe9trico}}

Ahora que tenemos claro lo que significa cifrar y descifrar mensajes,
tenemos que saber que hay dos formas fundamentales de hacerlo: cifrado
sim�trico y asim�trico.

Un sistema de cifrado sim�trico utiliza la misma clave para cifrar y
descrifrar el mensaje. En nuestro ejemplo anterior utilizamos la clave
\texttt{2} para cifrar nuestro mensaje utilizando el m�todo de C�sar y
utilizamos el mismo valor \texttt{2} (el cual restamos al mensaje
cifrado) para recuperar el mensaje original.

Los sistemas sim�tricos son en los que la mayor�a de nosotros pensamos
cuando nos hablan de c�digos y cifrados. Probablemente por que estamos
acostumbrados de las pel�culas donde tienen que conseguir una clave para
acceder a ficheros, ordenadores o incluso habitaciones. Sin embargo
existe otro tipo de cifrados que son lo que realmente han revolucionado
las comunicaciones en los �ltimos a�os.

\begin{entradilla}
{\em Los {\color{introcolor}sistemas asim�tricos utilizan claves distintas} para cifrar y descifrar}
\end{entradilla}


Los sistemas de cifrado asim�trico tienen dos claves y lo que nos
permiten es que si ciframos un mensaje con una de las claves, solo
podemos recuperarlo con la otra\ldots. y viceversa. Esto abre un sin fin
de posibilidades para los fikis criptogr�ficos. Las claves suelen
recibir el nombre de clave p�blica y clave privada, ya que, normalmente,
una de ellas la conoce todo el mundo, mientras que la otra es importante
que no la conozca nadie.


Veamos como funciona este sistema cuando queremos cifrar un mensaje. Con
un sistema sim�trico, debemos compartir la clave que usemos con la
persona a la que le queremos enviar el mensaje\ldots{} pero con un
sistema asim�trico la cosa es mucho m�s guay.

Como la clave p�blica de cada persona es\ldots{} bueno\ldots{} p�blica.
Vamos que la conocemos. Entonces solo tenemos que cifrar nuestro mensaje
con la clave p�blica de la persona a la que se lo queremos enviar. Como
el sistema es asim�trico, la �nica forma de descifrar el mensaje es
utilizando la clave privada\ldots{} que solo esa persona conoce\ldots.
Voil�. Veis la diferencia, con el sistema sim�trico en el que tenemos
que hacer llegar la clave a la otra persona de alguna forma. Si esa
clave es interceptada mientras la intercambiamos, entonces el mensaje
estar� comprometido.

Hay formas seguras de intercambiar claves\ldots{} pero de eso hablaremos
m�s tarde.

El algoritmo de cifrado asim�trico m�s famoso es RSA que toma su nombre
de sus inventores: Rivest, Shamir y Adleman) y se basa, as� a lo bestia,
en la factorizaci�n de n�meros primos muy grandes.


\hypertarget{fundamentos-criptogruxe1ficos-3.-hashes}{%
\sectiontext{white}{black}{HASHES}\label{fundamentos-criptogruxe1ficos-3.-hashes}}

El otro concepto criptogr�fico que debemos introducir para entender
muchos de los servicios seguros que existen hoy en d�a es el de las
funciones hash.

Una funci�n \emph{HASH} es una funci�n que es capaz de convertir una
gran cantidad de datos que le pasamos como entrada en un n�mero de un
tama�o fijo. Pero este n�mero tiene una propiedad muy especial. Si
modificamos solo uno de los datos que le pasamos a la funci�n como
entrada el resultado ser� totalmente diferente y cambiar� de una forma
que no se puede predecir.

\begin{entradilla}
{\em Una {\color{introcolor}funci�n hash} genera un valor de tama�o fijo a partir de un n�mero arbitrario de datos}
\end{entradilla}


Las funciones \emph{hash} son la base de los sistemas de firma digital y
de algunos m�todos de autenticaci�n modernos. Las funciones \emph{hash}
normalmente utilizan la representaci�n interna en binario de los
n�meros, ya que necesitan generar una salida de tama�o fijo. Para ver
como funciona esto vamos a usar una funci�n hash muy sencilla y mala,
pero para entender el concepto nos va a valer.

Tomemos la codificaci�n de n�meros y letras que usamos para nuestro
c�digo C�sar y vamos a crear una funci�n HASH que genere un \emph{HASH}
de 2 caracteres. Para ello cogeremos los valores num�ricos de 4 en 4 y
los sumaremos y luego sumaremos los d�gitos hasta obtener un n�mero
entre 0 y 9. Como nuestro mensaje tiene 9 caracteres, completaremos el
�ltimo bloque con 0s.

\begin{verbatim}
H  O  L  A  
7  15 11 0  
M  U  N  D  
12 21 13 3  
O  
15   0 0 0
------------
34 36 24 3
8         <= 3+3+2  (Sumamos el 1er d�gito)
   17     <= 4+6+4+3(Sumamos el 2do d�gito)
8   8
\end{verbatim}

As� que ese ser�a nuestro Hash. Si ahora cambiamos nuestro mensaje a
\emph{Mola Mundo}, es decir, hemos cambiado una sola letra del mensaje.

\begin{verbatim}
M  O  L  A
12 15 11 0
M  U  N  D
12 21 13 3
O
15   0 0 0
------------
39 36 24 3
11        <=  3 + 3 + 2 + 0
   22     <=  9 + 6 + 4 + 3  
2  4
\end{verbatim}

Como pod�is ver, al cambiar la \texttt{H} por una \texttt{M}, el n�mero
que obtenemos es \texttt{24} en lugar de \texttt{88}. Esta funci�n HASH
es muy mala, o como dicen los expertos, no es criptogr�ficamente segura,
pero nos vale para entender la idea\ldots{} no?

Unos de los algoritmos de hash m�s populares durante mucho tiempo fue
MD5, pero hoy en d�a ya no se considera seguro y la familia de
algoritmos SHA se usa en su lugar. SHA-2 en su versi�n de 256 o 512 bits es el m�s popular
a fecha de hoy. Que una funci�n hash no se considere segura significa
que dado un resultado hash, es posible modificar el fichero original de
tal forma que la funci�n hash produzca el mismo resultado.

\hypertarget{firma-digital}{%
\sectiontext{white}{black}{FIRMA DIGITAL}\label{firma-digital}}

Con todos estos fundamentos de criptograf�a que ya dominamos estamos en
condiciones de entender como funciona una firma digital. Una firma
digital nos permite verificar que un determinado mensaje, documento o
informaci�n procede de una determinada persona y que adem�s ese mensaje
no ha sido modificado por una tercera parte.

El concepto es muy sencillo:

\begin{itemize}
\item
  Generamos un hash de los datos originales. En caso de que el documento
  cambie, el hash cambiar�a por completo
\item
  Firmamos el hash con nuestra clave privada. Esto permitir� a cualquier
  que conozca nuestra clave publica desencriptar el hash y recalcularlo
  para verificar que el documento no ha sido cambiado. Si alguien
  intentara modificar el documento, podr�a generar un nuevo hash, pero
  no podr�a encriptarlo ya que desconoce la clave privada del usuario original.
\end{itemize}

Como pod�is ver, la idea es muy sencilla, pero muy potente.

\begin{entradilla}
{\em Los algoritmos de {\color{introcolor}firma digital} combinan el cifrado asim�trico con funciones hash}
\end{entradilla}


As� que si. Firmar digitalmente un documento no tiene nada que ver con
a�adirle una imagen con el escaneo de nuestra firma.

\columnbreak

\hypertarget{cifrado-punto-a-punto}{%
\sectiontext{white}{black}{CIFRADO PUNTO A PUNTO}\label{cifrado-punto-a-punto}}

El cifrado punto a punto (o \emph{End to End}, como dicen los
americanos), se ha convertido en un eslogan de venta para ciertos
productos, y es importante entender como funciona todo esto, para que no
demos por supuesto cosas que no son. Lo que \emph{End to End}
simplemente significa que los datos se cifran en un extremo y se
descifran en el otro, de forma que da igual por cuantos puntos
intermedios pase o si alguien consigue capturarlos mientras viajan por
la red. Nadie podr� leer el contenido de nuestro mensaje.

La forma de asegurar la privacidad entre dos puntos es utilizando un
sistema de cifrado asim�trico. Cada parte cifra su mensaje con la clave
p�blica de la persona que lo va a recibir, de forma que solamente esa
persona pueda decodificar el mensaje.

Hasta aqu� todo claro, sin embargo hay algunas peque�as cuestiones
t�cnicas a tener en cuenta.

\hypertarget{cifrado-asimuxe9trico-lento}{%
\sectiontext{white}{black}{EL CIFRADO ASIM�TRICO ES LENTO}\label{cifrado-asimuxe9trico-lento}}

La primera es que los algoritmos de cifrado asim�trico son muy lentos y
pesados y a d�a de hoy no son pr�cticos para cifrar grandes cantidades
de datos, como documentos o ficheros multimedia. En estos casos, lo que
se hace es generar una clave aleatoria para un sistema de cifrado
sim�trico, cifrar el documento usando esa clave y el sistema sim�trico
(por ejemplo AES-256) y a�adir al mensaje esa clave, pero cifrada con la
clave p�blica del destinatario. De esa forma, solo el destinatario puede
recuperar la clave sim�trica y acceder al documento.

\begin{entradilla}
{\em Los algoritmos de {\color{introcolor}cifrado asim�trico son muy lentos} y no se pueden utilizar para cifrar grandes cantidades de datos}
\end{entradilla}


Resumiendo cuando cifras un mensaje utilizando la clave p�blica de
alguien utilizando, por ejemplo RSA, no estas cifrando todo el mensaje
con el algoritmo RSA, solo una clave aleatoria que ha generado tu
ordenador. Para transferir ficheros por, por ejemplo, mail, esto no es
un gran problema, pero en comunicaciones m�s interactivas el uso de un
algoritmo asim�trico para el cifrado de todos los datos es totalmente
inviable.

\hypertarget{quien-genera-las-claves}{%
\sectiontext{white}{black}{�QUI�N GENERA LAS CLAVES?}\label{quien-genera-las-claves}}

El sistema que hemos descrito es realmente muy seguro. Al menos es tan
seguro como segura es tu clave privada. El gran problema con muchos
servicios disponibles hoy en d�a es que realmente no sabemos como
gestionan todo esto. Cuando inicias un chat en WhatsApp, un bonito
mensaje diciendo que tus comunicaciones son seguras aparece en la
pantalla. Pero que significa eso?

Bueno, puesto que las claves no las podemos generar nosotros por
nuestros propios medios (como por ejemplo en un ordenador aislado no
conectado a ninguna red), y que el c�digo de la aplicaci�n es cerrado,
no existe ninguna garant�a de que la clave privada se haya generado en
nuestro dispositivo y no en un servidor remoto y, aunque hubiera sido
generada en nuestro dispositivo, tampoco tenemos ninguna garant�a de que
esa clave no se haya enviado a alg�n servidor remoto.

Y que significa esto?\ldots{} Bueno, que cualquiera que tenga las claves
privadas de alguien puede ver todos sus mensajes. Si este que tiene las
claves es el servidor de WhatsApp, pues, aunque los mensajes y las fotos
comprometedoras est�n cifradas \emph{End to End}, WhatsApp va a verlo
todo\ldots{}

\hypertarget{grupos}{%
\sectiontext{white}{black}{CHATS DE GRUPOS}\label{grupos}}

Y que pasa con los grupos? Bueno, la verdad que no se como se implementa
la encriptaci�n en estos casos en cada una de las aplicaciones de chat
ah� fuera. Si bien hay varias opciones, una de las mas sencillas (y
aparentemente la usada por, por ejemplo WhatsAPP) es utilizar una clave
sim�trica para cifrar todos los mensajes y distribuir esa clave
utilizando la clave privada de cada miembro cuando este se une al grupo.

Hay esquemas m�s complejos como el que utiliza Signal, pero en general,
la seguridad del grupo depende de mantener la clave utilizada para
cifrar los mensajes del grupo segura, es decir, que solo salga del
dispositivo que la creo, cifrada con la clave publica de los usuarios
que deben usarla. Una vez m�s, si la aplicaci�n que utilizamos es de
c�digo cerrado, no podemos tener la certeza de que esa clave no acabe
siendo filtrada al servidor de la compa��a que gestiona el servicio de
alguna forma.

\hypertarget{cifrado-punto-a-punto.-el-caso-del-correo}{%
\sectiontext{white}{black}{CIFRADO PUNTO A PUNTO.\\EL CASO DEL CORREO}\label{cifrado-punto-a-punto.-el-caso-del-correo}}

El cifrado punto a punto para el caso del correo electr�nico es mucho
m�s complicado e, independientemente de lo que se pueda decir por ah�,
la �nica soluci�n totalmente privada es utilizar GPG para cifrar tus
correos en tu propio dispositivo.

El problema con el correo es que no hay forma de asegurar que el sistema
de correo que utilizan otras personas soporte encriptaci�n. As� que si
enviamos un mensaje a alguien que utiliza correo simple, el mensaje no
puede estar cifrado, o de lo contrario esa persona no podr� leerlo. Esto
tiene una serie de consecuencias:

\begin{itemize}
\item
  La primera es que solo es posible enviar mensajes encriptados punto a
  punto entre usuarios de una determinada plataforma
\item
  Cualquier mensaje desde un usuario utilizando una plataforma que
  permite encriptaci�n E2E a otro usuario de otra plataforma que no lo
  soporte ser� visible por ambas plataformas de correo.
\end{itemize}

Esto �ltimo incluye respuestas y mensajes re-enviados. La plataforma
proporcionando soporte para el cifrado podr� almacenar el mensaje
cifrado, y lo enviar� cifrado a su usuario, pero esa plataforma
recibir�/enviar� el mensaje en texto plano (sin cifrar).

Esta es una limitaci�n dif�cil de solucionar cuando un determinado
servicio soporta distintos tipos de usuarios, unos que usan encriptaci�n
y otros que no.

\hypertarget{cifrado-punto-a-punto.-sitios-web}{%
\sectiontext{white}{black}{CIFRADO PUNTO A PUINTO.\\SITIOS WEB}\label{cifrado-punto-a-punto.-sitios-web}}

El acceso a los sitios web est� cifrado, siempre que utilicemos el
protocolo HTTPS, lo cual hoy en d�a es lo habitual. La mayor�a de los
exploradores utilizan ese protocolo por defecto y si no est� disponible
muestran un aviso al usuario indicando que el sitio web al que intenta
acceder es inseguro. Si bien, dependiendo del servicio que estemos
intentando acceder, puede ser una exageraci�n.

Pero que significa que el acceso a un sitio web sea seguro?. Eso
solamente significa que el tr�fico desde nuestro explorador hasta el
servidor est� cifrado. En otras palabras cualquier entidad que vea el
tr�fico que va desde nuestra casa al servidor, no podr� ver que
contiene. Eso incluye al ISP que nos da acceso a internet y todos los
nodos intermedios de la red por lo que tengamos que pasar para alcanzar el
servidor web.

Si lo que estamos accediendo es un servicio de correo web, que ese
enlace este cifrado es muy importante, ya que de lo contrario, cualquier
que pueda ver el tr�fico ver�a nuestros correos electr�nicos, lo cual es
algo que es considerado informaci�n privada.

Sin embargo, si consideramos el acceso a un blog, o la p�gina web de una
distro linux que queremos bajar, en la que no vamos a intercambiar ning�n
tipo de informaci�n privada (no tenemos ni una cuenta en el sitio), pues
nos da exactamente igual que el sitio web sea seguro o no. Lo que va
a ver alguien que capture el tr�fico es algo que puede ver sin m�s que
acceder a la misma p�gina web.

Pensad en la web de nuestra revista. No necesit�is un usuario para
registraros, ni enviar un correo electr�nico o una clave. Simplemente
busc�is un art�culo y lo le�is. Poder ver vuestro tr�fico al leer un
art�culo le va a permitir al atacante ver el art�culo\ldots{} lo cual ya
puede hacer. Lo �nico que va a ver el curioso es que art�culos est�s
leyendo, lo cual rara vez es algo importante, pero bueno, es
informaci�n.

\hypertarget{como-funciona-https}{%
\sectiontext{white}{black}{�C�MO FUNCIONA HTTPS?}\label{como-funciona-https}}

Para entender mejor las consecuencias de seguridad de usar o no HTTPS,
veamos primero como funciona. HTTPS es una extensi�n del protocolo HTTP,
el protocolo original desarrollado para la web, la diferencia es que la
informaci�n que se env�a y recibe del servidor est� cifrada utilizando
una capa intermedia que puede ser \texttt{SSL} (\emph{Secure Socket
Layer} o Capa de Conexiones Seguras) o \texttt{TLS} (\emph{Transport
Layer Security} Seguridad de Capa de Transporte).

Estas capas intermedias permiten establecer un di�logo inicial con el
que intercambiar informaci�n sobre como establecer una comunicaci�n
segura.

Lo primero que hacen estas capas es negociar una clave de sesi�n
utilizando las claves p�blicas y privadas del cliente y el servidor.
Est�s claves est�n almacenadas en tu browser y no tienes ning�n control
sobre ellas, al menos no de forma directa. Es m�s, el browser podr�a
generar uno de estos pares de claves p�blica y privada por cada
conexi�n, ya que, en HTTPS, el extremo del explorador no necesita ser
autenticado.

Como ya hemos visto antes, esa clave de sesi�n se utiliza para cifrar el
tr�fico utilizando un algoritmo de cifrado sim�trico, el cual tambi�n es
negociado al principio de la sesi�n dependiendo de los algoritmos
soportados por cada parte (ambos necesitan usar el mismo algoritmo
adem�s de la misma clave).

\hypertarget{certificados}{%
\sectiontext{white}{black}{CERTIFICADOS}\label{certificados}}

Hemos dicho en la secci�n anterior que el browser no necesita
autenticarse, pero en general, nosotros si que queremos que el servidor
se autentique. Es decir, queremos estar seguros de que si vamos a
acceder a nuestra cuenta bancaria a trav�s de una p�gina web, la p�gina
a la que estamos accediendo es la de nuestro banco, y no la de alg�n
estafador que intenta robarnos los cuartos. Esta t�cnica se conoce como
phising, y muy pronto ver�s porque funciona incluso con el sistema de
autenticaci�n que existe en la actualidad.

\begin{entradilla}
{\em Los {\color{introcolor}certificados} nos permiten verificar que algui�n es quien dice ser.}
\end{entradilla}

Bien, la forma de verificar que un servidor web es quien dice ser es
utilizando una firma digital. Nuestro banco nos va a enviar un fichero
que dice: ``He soy tu banco. Cr�eme'' y eso lo va a firmar digitalmente
utilizando su clave privada\ldots{}. El
problema con este esquema es que si la clave publica compartida por el
servidor viene del propio servidor\ldots{} bueno, la firma no vale nada
en este caso. Sabemos que el mensaje proviene de ese servidor y que no
ha sido modificado\ldots{} pero eso no nos dice nada sobre la veracidad
de lo que el mensaje contiene.

En general, un fichero con informaci�n sobre la identidad de \emph{algo}
en internet que ha sido firmado digitalmente se conoce como certificado.
Cuando el que emite el certificado es el mismo que lo firma, se dice que
el certificado ha sido autofirmado, y eso es una posible brecha de
seguridad. Esta situaci�n es com�n durante desarrollo, ya que te da
igual tener un certificado real o no cuando est�s programando tu
aplicaci�n.

Sin embargo, como ya os hemos dicho, cualquiera puede crear un
certificado de esta forma, pero eso no nos asegura que el que emite el
certificado sea realmente quien dice ser. As� que lo que necesitamos es
una tercera entidad que de fe de quien es quien. Esta autoridad se
conoce como \emph{Certification Authority} (Autoridad de Certificaci�n,
el nombre es super sorprendente eh?), y normalmente nos referiremos a
ella como \texttt{CA},

As� que cuando quieres un certificado para tu sitio web, le pides a una
de estas CAs que genere un certificado para ti, el cual estar� cifrado
por esta CA, la cual es conocida por todos los exploradores. Ahora,
cuando intentamos acceder a este sitio web, el certificado que
recibimos, lo podemos verificar utilizando la \texttt{CA} la cual nos
asegurar� que el certificado es v�lido. Un certificado puede estar
incluso firmado por m�s de una \texttt{CA}.

\hypertarget{entonces-como-funciona-el-phising}{%
\sectiontext{white}{black}{ENTONCES, �C�MO PUEDE FUNCIONAR EL PHISING?}\label{entonces-como-funciona-el-phising}}

Bueno, aqu� esta el truqui. Si cualquier adquiere un dominio con un
nombre muy parecido a otro, pero diferente\ldots{} bueno, esa persona
puede aplicar por un certificado para ese dominio y se le conceder� ya
que todo est� correcto. El dominio, las direcciones de correo,
etc\ldots{} son todas distintas a las del dominio al que se parece, as�
que no hay ning�n problema en emitir un certificado.

El problema con el phising es que nos enga�an utilizando un nombre muy
parecido a otro real\ldots. Y que, seamos realista, comprobar el
contenido del certificado no nos permite verificar de forma sencilla si
el sitio es real o no. A ver es real, pero si es el que nosotros
pensamos que es.



Si puls�is en el candado que aparece junto al nombre del sitio web al
que quer�is acceder, podr�is, tras unos poco clicks, ver el certificado.
La mitad de la informaci�n son n�meros en hexadecimal que te dicen poco
o nada sobre la identidad real del due�o del certificado. El resto
parece informaci�n legal/empresarial. En otras palabras, los
certificados de sitios web, para los meros mortales como nosotros, no
son ninguna garant�a de que estemos accediendo al sitio que queremos
acceder.

El sistema funciona, pero realmente no es �til. El certificado nos dice
que estamos accediendo al sitio correcto. Y el sitio es correcto, lo que
pasa es que han usado un nombre muy parecido al de otro sitio, y adem�s
han hecho que las p�ginas parezcan exactamente iguales, con lo que, a
efectos pr�cticos asumiremos que estamos accediendo al sitio correcto,
ya que la mayor�a de la gente no va a ser capaz de extraer ninguna
informaci�n �til del certificado.

\begin{entradilla}
{\em Los {\color{introcolor}ataques de PHISING} funcionan porque tener un certificado no significa ser quien intentas parecer ser.}
\end{entradilla}


Es decir, aunque chequees el certificado te vas a quedar igual que
estabas. La mejor forma de defenderte contra el phising es no pulsar en
ning�n enlace recibido por correo o mensaje. Si tienes que hacer algo
con tu banco, vete a la web de tu banco y haz lo que tengas que
hacer\ldots. no uses enlaces que te lleguen por distintos medios.
Incluso si parecen venir de una fuente fiable pueden ser falsos.


\hypertarget{conclusiones}{%
\sectiontext{white}{black}{CONCLUSIONES}\label{conclusiones}}

En esta primera parte de esta breve introducci�n a la criptorgraf�a para meros mortales hemos introducido los conceptos m�s generales relacionados con esta disciplina y algunos de los casos cotidianos en los que la utilizamos aunque no seamos conscientes de ello. Hemos intentado explicaros como se usa esta tecnolog�a en aplicaciones y servicios que usamos todos los d�as para que as� se�is m�s conscientes de los peligros de un uso incorrecto, o mejor a�n, como usar los servicios de forma correcta y as� permanecer seguros.

En la pr�xima entrega os contaremos las formas en las que se pueden atacar este tipo de servicios y que hacer para que ese tipos de ataques no sean efectivos o al menos dificultarlos lo m�ximo posible.

%\raggedcolumns
%\clearpage
\end{multicols}

\rput(7.9,-3.7){\resizebox{!}{9cm}{{\epsfbox{images/promo/promo03.eps}}}}

\hypertarget{fpn:bucle}{}\label{fpn:bucle}
\pagestyle{fpn}

\rput(7.9,-0.5){\resizebox{!}{14cm}{{\epsfbox{images/header-loop.eps}}}}

\psset{fillstyle=solid}
\psframe[fillcolor=black,opacity=0.7](2,-4.5)(17,0)



% -------------------------------------------------
% Cabecera
\begin{flushright}


{\color{pink}\mtitle{14cm}{Bucle Infinito en C. Modo esot�rico}}

\msubtitle{8cm}{Fen�menos Paranormales}

{\sf\color{white}{ por Carolyn Lightrun}}

%{{\psset{linecolor=black,linestyle=dotted}\psline(-12,0)}}
\end{flushright}


\vspace{2mm}
% -------------------------------------------------
\lstset{language=C,frame=tb,framesep=5pt,basicstyle=\scriptsize}

\intro{introcolor}{L}{a maravilla del lenguaje C es que en ocasiones pasan cosas
inexplicables que hablen la puerta al misterio y la superstici�n. Esa es
la raz�n por la que es bastante com�n que los programadores C muestren
sus programas en noches de luna llena tras sacrificar al menos una
gallina.}

\begin{multicols}{2}



Si, C es as�, a veces ocurren fen�menos inexplicables y mientras sean
inexplicables pensaremos que esos hechos est�n controlados por entes
superiores, Dioses del cach� o semidioses de la laguna oscura de la
pila. Pero una vez que encuentras la explicaci�n\ldots{} simplemente se
convierten en bugs.

\hypertarget{bucles-infinitos}{%
\sectiontext{white}{black}{BUCLES INFINITOS}\label{bucles-infinitos}}

Antes de mostraros nuestro bucle infinito modo esot�rico vamos a ver las
distintas formas en las que podemos hacer un bucle infinito en C.

Por cierto, un bucle infinito es una parte de un programa que se ejecuta
una y otra vez sin fin\ldots{} bueno, si, es exactamente lo que su
nombre indica. Normalmente, dentro del bucle suele haber una condici�n
de parada, pero en ocasiones la �nica forma de pararlo es matando el
proceso.

\hypertarget{bucle-infinito-while}{%
\sectiontext{white}{black}{\texorpdfstring{BUCLE INFINITO
\texttt{while}}{Bucle infinito while}}\label{bucle-infinito-while}}

La primera forma de crear un bucle infinito y, probablemente la m�s
com�n y evidente es utilizar un bucle\verb!while!de
la siguiente forma:

\begin{lstlisting}[language=C]
...
while (1) {
    puts ("Me m�s que el chorizo");
}
\end{lstlisting}

Puesto que la condici�n en el\verb!while!es siempre
cierta, todo lo que pongamos en su interior se repetir� hasta que
interrumpamos el programa.

Podemos poner el n�mero que m�s rabia nos de en el
{\verb!while!}, mientras sea distinto de
{\verb!0!}, la condici�n se considerar� verdadera y el
bucle no parar� de ejecutarse.

\hypertarget{bucle-infinito-for}{%
\sectiontext{white}{black}{\texorpdfstring{BUCLE INFINITO
\texttt{for}}{Bucle infinito for}}\label{bucle-infinito-for}}

Otra forma, menos habitual de implementar bucles infinitos es utilizando
{\verb!for!} de la siguiente forma:

\begin{lstlisting}[language=C]
  for (;;) {
    puts ("Me repito m�s que el chorizo");
  }
\end{lstlisting}

En este caso, al no pasar ninguna condici�n se toma por defecto el valor
true. La expresi�n anterior es equivalente a:

\begin{lstlisting}[language=C]
  for (;1;) {
    puts ("Me repito m�s que el chorizo");
  }
\end{lstlisting}

\hypertarget{bucle-infinito-goto}{%
\sectiontext{white}{black}{\texorpdfstring{BUCLE INFINITO
\texttt{goto}}{Bucle infinito goto}}\label{bucle-infinito-goto}}

Este si que no lo vais a ver en el mundo real, ya que todo el mundo
considera los {\verb!goto!}s peligrosos, si bien, tanto
los bucles\verb!while!como los
{\verb!for!} se implementan en ensamblador con un
{\verb!goto!} o salto. Podemos escribir un bucle
infinito con\verb!goto!de la siguiente forma:

\begin{lstlisting}[language=C]
bucle:
    puts ("Me repito m�s que el chorizo");
    goto bucle;
\end{lstlisting}

\hypertarget{bucle-infinito-esotuxe9rico}{%
\sectiontext{white}{black}{BUCLE INFINITO ESOT�RICO}\label{bucle-infinito-esotuxe9rico}}

Para terminar os dejamos una forma de implementar un bucle infinito un
tanto esot�rica. Ve�moslo y luego lo discutimos.

\begin{lstlisting}[language=C]
  for (unsigned char i = 0; i < 256; i++)
    puts ("Me repito m�s que el chorizo");
\end{lstlisting}

Puesto que nosotros queremos hacer un bucle infinito, nuestro programa
es perfectamente v�lido, sin embargo cuando vemos algo como esto es
habitual que la intenci�n del programador no fuera programar un bucle
infinito.

\begin{entradilla}
{\em Un peque�o despieste y un {\color{introcolor}peque�o bucle} se convierte en un {\color{introcolor}bucle INFINITO}}
\end{entradilla}

De hecho esto es algo que me pas� hace muchos a�os cuando empezaba a
programar en C. Aqu� esta la historia:

\sectiontext{white}{black}{UNA HISTORIA DE MIEDO}

Cuando empiezas a programar en C, ves todos esos tipos con distintos
tama�os e intentas ser guay y usar el tipo que mejor se adecue a lo que
tienes que hacer. As� que si vas a hacer un bucle de unas pocas
interacciones, lo suyo es utilizar un byte como contador para no
desperdiciar memoria. Luego aprendes, muchos a�os despu�s que realmente
no estas ahorrando nada\ldots{} pero eso es otra historia.

Pues resulta que estaba programando alg�n tipo de algoritmo de procesado
de se�al. No recuerdo exactamente lo que era, pero en alg�n lugar hab�a
un peque�o bucle de 32 o 64 iteraciones. Todo funcionaba bien, hasta que
aument� el valor a 128 y de repente mi programa dej� de funcionar.

\begin{entradilla}
{\em Este {\color{introcolor}tipo de errores} suelen ser dif�ciles de ver una vez se cometen.}
\end{entradilla}


Recuerdo que me llev� alg�n tiempo reaccionar. Pensaba: ``el programa
funcionaba hace 2 minutos y no he cambiado el algoritmo�. Us� el
depurador y me llevo un rato entender porque cuando llegaba a la �ltima
interacci�n mi contador se reseteaba y el bucle comenzaba a empezar.



En aquel momento estaba usando un contador
{\verb!char!} que en los 90 era con signo por defecto
(esto ha ido cambiando varias veces en las �ltimas d�cadas) y por lo
tanto, al llegar a 128 se produc�a un overflow y el siguiente valor
pasaba a ser -127 en lugar de 129.

\hypertarget{el-bucle-infinito-explicado}{%
\sectiontext{white}{black}{EL BUCLE INFINITO EXPLICADO}\label{el-bucle-infinito-explicado}}

En el c�digo de arriba pasa exactamente lo mismo.
{\verb!i!} es un
{\verb!unsigned char!} que toma valores entre
{\verb!0!} y {\verb!255!}. En la
�ltima iteraci�n, cuando\verb!i!vale
{\verb!255!}\verb!for!hace la
comparaci�n\verb!i < 256!la cual es falsa, ya que
{\verb!i!} todav�a es {\verb!255!},
as� que ejecuta el bucle una vez m�s.

\begin{entradilla}
{\em Los {\color{introcolor}Overflows en tipos enteros} son {\color{introcolor}bugs} habituales}
\end{entradilla}


Al terminar el bucle,\verb!for!ejecuta la �ltima
instrucci�n {\verb!i++!}. Como
{\verb!i!} es\verb!255!al sumarle
uno pasa a ser\verb!256!pero ese valor no se puede
almacenar en un byte ({\verb!unsigned char!}) con lo
que solo los 8 bits menos significativos se asignan a la variable
{\verb!i!} y por lo tanto toma el valor
{\verb!0!} el cual es menor que
{\verb!256!} y por lo tanto el bucle continua
infinitamente.

\hypertarget{conclusiuxf3n}{%
\sectiontext{white}{black}{CONCLUSIONES}\label{conclusiuxf3n}}

Bueno, la conclusi�n m�s importante es que debemos ser consciente del
tama�o de nuestras variables y evitar que se produzcan desbordamientos
que pueden provocar\ldots{} fen�menos extra�os\ldots.. uuuuuuhhhhhhh



\rput(5.0,-12.5){\resizebox{!}{26cm}{{\epsfbox{images/promo/promo02.eps}}}}
\raggedcolumns
\clearpage
\end{multicols}



%\rput(2.5,-1.5){\resizebox{!}{4.5cm}{{\epsfbox{images/general/trucos1.eps}}}}

\hypertarget{divulgacion:compresion1}{}\label{divulgacion:compresion}
\pagestyle{divulgacion}

\rput(7.9,-2.5){\resizebox{!}{11cm}{{\epsfbox{images/header-compresion.eps}}}}

\psset{fillstyle=solid}
\psframe[fillcolor=black,opacity=0.6](2,-4.5)(17,0)


% -------------------------------------------------
% Cabecera
\begin{flushright}


{\color{purple}\mtitle{12cm}{Compresi�n de Datos sin P�rdidas}}

\msubtitle{10cm}{Repasamos los principales algoritmos}

{\sf\color{white}{ por Signul Floid}}

%{\psset{linecolor=black,linestyle=dotted}\psline(-12,0)}
\end{flushright}

\vspace{2mm}
% -------------------------------------------------



\lstset{language=C,frame=tb,framesep=5pt,basicstyle=\scriptsize}

\intro{introcolor}{A}{lguna vez te has preguntado como es posible hacer los ficheros de tu
ordenador m�s peque�os?. Como es posible que \texttt{gzip} coja esos
ficheros de texto y los deje en nada?. O como se quedan las fuentes del
kernel en unos pocos megabytes?. Pues, ag�rrate los machos que todas
esas preguntas van a ser respondidas a continuaci�n\ldots. Si nunca te
has echo alguna de esas preguntas, creo que est�s leyendo la revista
equivocada.}


\begin{multicols}{2}

La compresi�n de datos es un proceso por el cual el \emph{espacio}
requerido para almacenar ciertos datos se reduce con el fin de\ldots.
bueno\ldots{} de ahorrar \emph{espacio}. Si bien hay t�cnicas de
procesado de se�al que se podr�an considerar como algoritmos de
compresi�n sobre se�ales anal�gicas, en este art�culo nos vamos a
centrar en la compresi�n de datos digitales\ldots. vamos, en como coger
una secuencia de n�meros y convertirla en otra secuencia de n�meros, m�s
corta, la cual, de alguna forma, podemos volver a expandir en la
secuencia original.


Vamos a empezar por lo m�s sencillo, e iremos introduciendo distintos
tipos de algoritmos de forma progresiva.

\hypertarget{run-lenght-encoding}{%
\sectiontext{white}{black}{RUN-LENGHT ENCODING}\label{run-lenght-encoding}}

La verdad que no se como traducir este tipo de compresi�n\ldots.
\emph{Codificacion de longitud de secuencia}?. La compresi�n
\href{https://en.wikipedia.org/wiki/Run-length_encoding}{RLE}, como es
normalmente conocida, es uno de los esquemas de compresi�n m�s
sencillos. Popular en los 90 para comprimir im�genes con paleta de
colores como \href{https://en.wikipedia.org/wiki/BMP_file_format}{BMP} o
\href{https://en.wikipedia.org/wiki/PCX}{PCX}, ha ca�do en desuso en la
actualidad en favor de algoritmos m�s eficientes. Sin embargo, la
sencillez de este m�todo hace que merezca la pena tenerlo en mente ya
que nos puede resultar �til en ciertos casos.

\begin{entradilla}
{\em {\color{introcolor}La codificaci�n RLE} es una de las formas m�s sencillas e intuitivas de comprimir datos}
\end{entradilla}


Vamos a ver como funciona esto con un sencillo ejemplo. Imaginemos que
queremos codificar la siguiente imagen monocroma, utilizando este
m�todo, donde el caracter \texttt{.} representa el color blanco y el
caracter \texttt{\#}, representa el color negro:
\columnbreak
\begin{verbatim}
................
..##########....
..#.............
..#.............
..#.............
..#####.........
..#.............
..#.............
..#.............
..##########....
................
\end{verbatim}

Suponiendo que almacenamos el ancho y el alto de forma separada (como
por ejemplo en una imagen
\href{https://en.wikipedia.org/wiki/Netpbm}{\texttt{ppm}}, podemos
almacenar todos los datos en una sola l�nea. Ahora, si os fij�is, que
pasar�a si en lugar de almacenar los 18 \texttt{.} del principio,
almacenamos \texttt{18.}\ldots{} Estar�amos pasando de 18 caracteres a
3\ldots. Voil�\ldots. Compresi�n.

El caracter \texttt{E} de m�s arriba est� inicialmente representado por
176 bytes. Utilizando el m�todo RLE, lo podr�amos representar como

\begin{verbatim}
18.8#6.1#16.1#16.1#16.5#11.1#16.1#16.8#18.
\end{verbatim}

Esto son 43 caracteres, lo que significa que hemos comprimido nuestra
imagen original unas 4 veces, o si lo prefer�s, un 75\%\ldots. Que no
est� nada mal.

\emph{En la realidad los algoritmos de compresi�n trabajan con bits, se
trata de aprovecha al m�ximo el espacio, pero para este primer ejemplo,
hablar de bytes y caracteres hace m�s sencillo entender la idea
general}.

\hypertarget{rle-en-pruxe1ctica}{%
\sectiontext{white}{black}{RLE EN LA PR�CTICA}\label{rle-en-pruxe1ctica}}

En la pr�ctica la implementaci�n de este m�todo de compresi�n es un poco
diferente. Como ya hemos dicho, este m�todo funciona especialmente bien
con datos en los que tenemos largas secuencias de un mismo valor.
Pensemos por ejemplo, en im�genes de dibujos animados, al menos en los
del siglo pasado, en la que se utilizaban colores planos .

Sin embargo, siempre hay partes en las que los n�meros cambian mucho. Si
pensamos en im�genes, esas partes corresponden con los detalles de la
imagen (las altas frecuencias, la parte de la imagen que cambia m�s).
Esos valores, muchas veces vienen en r�fagas, es decir, cuando los
valores cambian, no cambia solo uno. En este caso, RLE estar�a generando
secuencias m�s largas, es decir, en lugar de comprimir, estar�a haciendo
la secuencia m�s grande. Por ejmplo:

\begin{verbatim}
ABCDEF
1A1B1C1D1E1F
\end{verbatim}

Como podemos ver a simple vista, la secuencia generada por el algoritmo
es mucho m�s larga. Por esta raz�n se suele utilizar una codificaci�n un
poco m�s compleja. Vamos a contaros dos formas de mejorar esto. La
primera m�s general y la segunda utilizada en algunos arcanos formatos
de imagen :).

Una forma de evitar este problema es que si la longitud de la secuencia
es m�s de 3, duplicamos el caracter. Por ejemplo:

\begin{verbatim}
AAAAAABCDEFFFFGHHJ
AA6BCDEFF4GHHJ
^^     ^^ 
\end{verbatim}

Otra forma de aplicar este m�todo al trabajar con datos binarios es
reservar el bit de mayor peso para indicar si lo que sigue a
continuaci�n es una secuencia de un pixel repetido, o una secuencia de
valores diferentes. De esta forma, en el caso de utilizar valores de 8
bits, la secuencia m�s larga de valores repetidos que podemos
representar es 128, al igual que la secuencia m�s larga de valores no
repetidos.

Para ilustrar este m�todo, vamos a utilizar valores hexadecimales. El
bit m�s significativo (el que tiene un valor asociado m�s grande) cuando
est� activo genera el valor 0x80. Para que nos resulte f�cil ver la
codificaci�n vamos a representar secuencias de menos de 16 bytes, de
forma que los valores los podamos acomodar en el nibble inferior de
nuestro byte. Por cierto, este byte se suele referir como byte de
repetici�n. Con todo esto:

\begin{verbatim}
0x8C 0xFF 0x04 0x01 0x02 0x03 0x04 0x85 0xFF
\end{verbatim}

Esta secuencia se expandir�a en:


\begin{verbatim}
8C FF                   04      85 FF
FFFFFFFFFFFFFFFFFFFFFFFF01020304FFFFFFFFFF
\end{verbatim}


Este tipo de codificaci�n era com�n, con ciertas variaciones en formatos
de imagen como PCX o
\href{https://en.wikipedia.org/wiki/Truevision_TGA}{TGA}.

\hypertarget{lempel-ziv-y}{%
\sectiontext{white}{black}{LEMPEL, ZIV Y y \ldots{}}\label{lempel-ziv-y}}

Los algoritmos LZ descritos en el legendario paper de Abraham Lempel y
Jacob Ziv en 1977 (ahora ya sab�is de donde viene el nombre LZ),
abrieron la puerta a toda una familia de algoritmos de compresi�n que
todav�a son utilizados masivamente hoy en d�a. El paper de 1977
describ�a dos algoritmos conocidos como LZ1 y LZ2 o
\href{https://en.wikipedia.org/wiki/LZ77_and_LZ78}{LZ77 y LZ78}.

Estos algoritmos sentaron las bases del
\href{https://en.wikipedia.org/wiki/Lempel\%E2\%80\%93Ziv\%E2\%80\%93Welch}{LZW}
que se populariz� con el formato
\href{https://en.wikipedia.org/wiki/GIF}{GIF} en 1984 gracias a la
contribuci�n de Terry Welch (la W en el nombre del algoritmo).
Posteriores versiones como
\href{https://en.wikipedia.org/wiki/Lempel\%E2\%80\%93Ziv\%E2\%80\%93Markov_chain_algorithm}{LZMA},
utilizado por ejemplo por el compresor \texttt{7z} y \texttt{xz}.

Para no aburriros, y no hacer esta secci�n supercomplicada, vamos a
explicar el algoritmo original y dar algunas indicaciones de como fue
modificado para dar lugar a otros memorables miembros de la familia LZ.

LZ77 fu� el algoritmo descrito en el paper original de 1977: \emph{A
Universal Algorithm for Sequential Data Compression}, publicado en el
m�tico \emph{IEEE Transaction on Information Theory}.

LZ77 es un algoritmo muy sencillo, pero es mucho m�s f�cil de entender
sobre un ejemplo, que explicado con palabras\ldots. pero bueno,
b�sicamente lo que el algoritmo hace es buscar ocurrencias del texto en
el propio texto.

\begin{entradilla}
{\em La familia de {\color{introcolor}algoritmos LZ marcaron un antes y un despu�s} en el mundo de la compresi�n de datos.}
\end{entradilla}


El algoritmo funciona desplazando una ventana imaginaria sobre los datos
a comprimir. En este ejemplo, vamos a considerar que el tama�o de la
ventana es muy grande o en otras palabras, que vamos a buscar tan atr�s
como queramos. En las implementaciones reales este tama�o de ventana
suele variar y usualmente se limita a unos 32Kb (valores de 2Kb y 4 Kb
son tambi�n comunes). LZ77 tiene como salida una serie de tuplas
(conjuntos de varios valores de distinta naturaleza)\ldots. un conjunto
de tres valores que vamos a representar como
\texttt{\textless{}d,l,c\textgreater{}}, donde:

\begin{itemize}
\item
  \texttt{d} es el desplazamiento desde la posici�n actual (el caracter
  que estamos comprimiendo) a la ocurrencia anterior m�s cercana. En
  caso de que no haya una ocurrencia anterior, este valor ser� 0.
\item
  \texttt{l} es la longitud de la ocurrencia anterior que se corresponde
  con la cadena actual.
\item
  \texttt{c} es el caracter que tenemos que a�adir a la secuencia que
  hemos encontrado para obtener el texto actual.
\end{itemize}

\hypertarget{comprimiendo-un-texto-con-lz77}{%
\sectiontext{white}{black}{COMPRIMIENDO UN TEXTO CON LZ77}\label{comprimiendo-un-texto-con-lz77}}

Con todo lo que hemos dicho vamos a comprimir el texto que mostramos a
continuacin. El caracter \texttt{\^{}} indica la posici�n actual del
algoritmo en el texto. El algoritmo buscar� ocurrencias de la cadena que
se encuentra a partir de ese s�mbolo, en la cadena que se encuentra
antes que �l. Comenzaremos en el primer caracter

\begin{verbatim}
abcbcbcbabbc
^
\end{verbatim}

Como nuestra secuencia comprimida est� vac�a, no podemos encontrar el
caracter \texttt{a} as� que lo a�adimos y nos movemos al siguiente
car�cter

\begin{verbatim}
abcbcbcbabbc
 ^
LZ77: (0,0,a)
\end{verbatim}

Lo mismo ocurrir� con la \texttt{b} y la \texttt{c} siguientes. As� que
tras esas dos iteraciones extra no encontraremos con el siguiente
estado:

\begin{verbatim}
abcbcbcbabbc
   ^
LZ77: (0,0,a), (0,0,b), (0,0,c)
\end{verbatim}

Si buscamos hacia atr�s, veremos que si nos desplazamos dos caracteres
hacia atr�s encontraremos una \texttt{b} que es nuestro caracter actual.
Ahora podemos intentar ver si hay m�s caracteres que coinciden. Vemos
que la \texttt{c} tambi�n coincide. As� que si nos movemos 2 posiciones
hac�a atr�s (valor \texttt{d} de la tupla), encontramos una secuencia de
2 caracteres (valor \texttt{l}) de la tupla. Una vez a�adidos esos dos
caracteres, necesitamos incluir el caracter \texttt{b} para completar la
secuencia (valor \texttt{c} de la tupla).

\begin{verbatim}
 VV 
abcbcbcbabbc
   ^^
LZ77: (0,0,a), (0,0,b), (0,0,c), (2,2,b)
\end{verbatim}

Llegados a este punto, nuestra ventana de b�squeda se ha terminado
(hemos alcanzado el caracter actual). En este caso se considera que los
caracteres que hemos a�adido ya pasan a formar parte del buffer de
b�squeda, de forma que el \texttt{bc} que ya hemos a�adido se convierte
en un nuevo match. De esta forma hemos sido capaces de encontrar 5
caracteres anteriores que coinciden con los que queremos comprimir,
seguidos del caracter \texttt{a}.

\begin{verbatim}
abcbcbcbabbc
         ^
LZ77: (0,0,a), (0,0,b), (0,0,c), (2,5,a)
      a        b         c        bcbcba
      
\end{verbatim}

Lo siguiente que encontramos es \texttt{bb}, que no existe en nuestra
secuencia de b�squeda, as� que emitiremos un nuevo \texttt{(0,0,b)}

\begin{verbatim}
abcbcbcbabbc
          ^
LZ77: (0,0,a), (0,0,b), (0,0,c), (2,5,a) 
      a        b         c        bcbcba  
      (0,0,b)
      b
\end{verbatim}

Y finalmente el �ltimo \texttt{bc} se convertir�a en:

\begin{verbatim}
abcbcbcbabbc
            ^
LZ77: (0,0,a), (0,0,b), (0,0,c), (2,5,a) 
      a        b         c        bcbcba  
      (0,0,b) (1,1,c)
      b       bc
\end{verbatim}


\hypertarget{codificando-la-salida}{%
\sectiontext{white}{black}{CODIFICANDO LA SALIDA}\label{codificando-la-salida}}

A primera vista, no parece claro que hayamos comprimido mucho nuestro
texto de entrada. Por cada secuencia que encontramos generamos 3
valores\ldots{} Bueno, en la pr�ctica no se suelen almacenar los tres
valores. La mayor�a de implementaciones definen dos tipos de salida. Un
par longitud, posici�n, y lo que se conoce como un literal o los valores
reales que estamos codificando.

As�, en nuestro ejemplo anterior la secuencia se convertir�a en:

\begin{verbatim}
<a><b><c><2,5><a><b><1,1><b>
\end{verbatim}

As� contado a lo bestia la secuencia inicial tiene 12 caracteres,
mientras que la comprimida tiene solo 10. La forma de diferencia entre
un tipo de salida y otro, as� como la codificaci�n final de los valores
depende de la implementaci�n. Por ejemplo, es com�n utilizar c�digo de
longitud variable para los valores de longitud.

Una de los formatos m�s comunes de LZ77 es
\href{https://en.wikipedia.org/wiki/Deflate}{DEFLATE}, desarrollado para
la versi�n 2 de la utilidad PKZIP (que permit�a comprimir ejecutables).
Esta implementaci�n estaba patentada, pero tras mucho trabajo una nueva
implementaci�n que no infling�a la patente fu� desarrollada y publicada
en el \href{https://www.ietf.org/rfc/rfc1951.txt}{RFC-1951}. Esta
implementaci�n libre de patentes permiti� el desarrollo de \texttt{gzip}
o el formato \href{https://en.wikipedia.org/wiki/PNG}{\texttt{PNG}}.

\begin{entradilla}
{\em El {\color{introcolor}formato DEFLATE es uno de los m�s utilizados en Intenet} y los soportan la mayor�a de servidores Web}
\end{entradilla}

El formato DEFLATE combina la salida del algoritmo LZ77 con la
codificaci�n Huffmann, de la que hablaremos en un rato, para definir un
formato binario comprimido.

\hypertarget{el-advenimiento-de-gif}{%
\sectiontext{white}{black}{EL ADVENIMIENTO DE GIF}\label{el-advenimiento-de-gif}}

Ahora que sabemos como funciona el algoritmo LZ77, vamos a ver una cosa
curiosa. Imaginad que queremos comprimir la siguiente secuencia:

\begin{verbatim}
aaaaaaaaaabbbbbbaaaaaaaaaa
\end{verbatim}

Aplicando el algoritmo obtenemos:

\begin{verbatim}
<a><1,9><b><1,5><a><1,7>
\end{verbatim}

Y aplicando RLE obtendr�amos:

\begin{verbatim}
aa10 bb6 aa9
\end{verbatim}

Dejando a un lado la forma final de codificar los datos en los fichero o
en la red, podemos ver que LZ77 funciona como un algoritmo RLE para
datos en los que tenemos largas secuencia de datos iguales. O dicho de
otra forma, el algoritmo se comporta de forma similar a RLE para los
casos m�s favorables para RLE, pero adem�s funciona bien en otros muchos
casos.

Como dijimos m�s arriba, en 1984 la variante LZW fu� publicada, y
algunos a�os despu�s se creo el formato GIF que hac�a uso de este tipo
de compresi�n. Dejando a un lado otras capacidades de GIF como la
posibilidad de hacer animaciones, o la carga progresiva que fueron clave
en los inicios de la Web en Internet, este formato obten�a mejores
resultados que los formatos basados en RLE y r�pidamente desbanc� a
todos ellos.

\begin{entradilla}
{\em En los albores de Internet, {\color{introcolor}el formato GIF era el m�s utilizado}. A�n hoy en d�a es el formato m�s com�n para peque�as animaci�nes.}
\end{entradilla}

La raz�n de que GIF ofreciera mejores resultados\ldots{} bueno, la
acabamos de ver, digamos que como poco LZ77 es tan bueno como RLE, y LZW
es mejor que LZ77.. Lamentablemente la historia de LZW est� int�mamente
ligada al mundo de las patentes y el formato GIF desarrollado por
Compuserve estuvo siempre afectado por la controversia entre los
acuerdos de Compuserrve y Unisys (poseedor de la patente de LZW).
Finalmente, como os comentamos m�s arriba, el desarroll� PNG como
alternativa a GIF y soluci�n al conflicto. En 2004 todas las patentes
relacionadas con GIF expiraron y ahora puede utilizarse libremente, si
bien, hoy en d�a su uso est� mucho m�s limitado que en los primeros
tiempos de Internet.

\hypertarget{cuxf3digos-de-longitud-variable}{%
\sectiontext{white}{black}{C�DIGOS DE LONGITUD VARIABLE}\label{cuxf3digos-de-longitud-variable}}

Normalmente, cuando se habla de compresi�n se suele empezar con esta
parte que es la m�s matem�tica, pero creo que haber empezado por los
algoritmos m�s visuales hace m�s sencillo introducir esta parte\ldots{}
Adem�s, si ahora te pierdes o te cansas\ldots{} bueno ya has le�do la
mitad del art�culo y con suerte aprendido un par de cosas :).

La idea detr�s de los c�digos de longitud variable es muy sencilla:
\emph{�Por qu� no codificamos los caracteres que aparecen m�s a menudo
con menos bits?}. Sencillo no?

En general, la forma m�s sencilla de codificar un mensaje es asignar
c�digos de longitud fija. Imaginemos que queremos codificar 4 s�mbolos
para un sistema de medici�n de precipitaciones, algo como:

\begin{verbatim}
A : Caen Chuzos de Punta
B : Llueve
C : Chispea
D : M�s seco que la mohama. 
\end{verbatim}

La forma habitual de codificar esto es utilizar un c�digo de longitud
fija. Como tenemos 4 s�mbolos necesitamos 2 bits, ya que
\texttt{2\^{}2\ =\ 4} o si lo prefer�s, \texttt{log2\ 4\ =\ 2} y
har�amos algo como esto:

\begin{verbatim}
A -> 00
B -> 10
C -> 01
D -> 11
\end{verbatim}

Perfecto no?\ldots{} Pues s�. Este c�digo funciona perfectamente en
cualquier caso, pero hay casos en los que estaremos desperdiciando bits
y por lo tanto ancho de banda, en el caso de transmitir estos datos de
nuestros sensores a un servidor, o espacio de almacenamiento en el caso
que guardemos los datos en alg�n tipo de memoria. Vamos a introducirnos
en el mundo de la
\href{https://en.wikipedia.org/wiki/Information_theory}{Teor�a de la
Informaci�n} para intentar entender mejor como podemos aprovechar mejor
nuestro ancho de banda o nuestro espacio de almacenamiento.

\hypertarget{la-teoruxeda-de-la-informaciuxf3n}{%
\sectiontext{white}{black}{LA TEOR�A DE LA INFORMACI�N}\label{la-teoruxeda-de-la-informaciuxf3n}}

Como su propio nombre indica, el elemento principal de la teor�a de la
informaci�n es la informaci�n. De esto vamos a hablar en un momentito,
pero antes deb�is saber que todo el fundamento matem�tico tras la teor�a
de la informaci�n, comenz� all� por los locos a�os 20. Primero con
Nyquist (el del l�mite que lleva su mismo nombre) y Hartley, (que dio
nombre a la unidad de medida de informaci�n en base 10\ldots{} como el
bit en base dos, pero en base 10), pero la parte de la que vamos a
hablar aqu� fue introducida unos a�os m�s tarde por
\href{https://en.wikipedia.org/wiki/Claude_Shannon}{Claude Shannon},
quien en un paper de unas 50 p�ginas titulado
\href{https://people.math.harvard.edu/~ctm/home/text/others/shannon/entropy/entropy.pdf}{\emph{A
mathematical Theory of Communications}} sent� las bases de la teor�a de
la informaci�n moderna.

En lo que resta de art�culo vamos a introduciros una peque�a parte de
los conceptos que Shannon desarrolla en ese m�tico art�culo, el cual os
recomendamos que le�is si os interesa el tema, ya que desvela los
fundamentos de muchas tecnolog�as actuales.

Ya esta bien de historia, la parte que nos interesa a nosotros, desde el
punto de vista de la compresi�n de datos, es en la que Shannon define lo
que es la informaci�n y nos da una forma de medirla. Y el concepto es
sorprendentemente sencillo.

Imaginemos que nuestro sistema de sensores de lluvia, cuyos mensajes
hemos codificado con 2 bits, lo instalamos en Amsterdam. Por si no lo
sab�is, los Paises Bajos tienen un tiempo de mi*\ldots{} muy malo y
llueve muy a menudo. En invierno, en verano,\ldots{} Si vives en
Amsterdam y alguien te dice que va a llover, pensar�as algo as� como:
\emph{Vaya, lo t�pico, menuda novedad}. Pero si te dicen que ma�ana va a
hacer sol\ldots. ah, amigo, la cosa cambia mucho. La gente se pedir�
incluso d�as libres en el trabajo, por que\ldots{} bueno, los pocos d�as
de Sol que hay, hay que aprovecharlos.

Sabiendo esto, de los mensajes que definimos anteriormente para nuestro
sistema de alertas de lluvia, cual cre�is que es el mensaje que nos da
m�s informaci�n?\ldots Efectivamente, el D. Pero que pasar�a si
instalamos el sistema en el desierto de Atacama?, el lugar m�s seco de
la tierra?\ldots{} All� que llueva es una gran noticia no?.

Pues con todo esto que os acabamos de contar, parece l�gico pensar que
la cantidad de informaci�n que proporciona un mensaje es inversamente
proporcional a la probabilidad de que ese mensaje ocurra, o si lo
prefer�s que el evento asociado a ese mensaje ocurra\ldots{} que en
nuestro caso es lo mismo. Cuando menor sea la probabilidad de que el
mensaje ocurra, m�s informaci�n nos estar� proporcionando ese mensaje.
As� que podemos concluir que:


\begin{center}
$ I(x) = log (\frac{1}{p(x)}) = - log (p(x))$
\end{center}

Donde \texttt{I(x)} es la cantidad de informaci�n del mensaje \texttt{x}
y \texttt{p(x)} es la probabilidad de que ocurra ese mensaje (o evento
si lo queremos ver desde un punto de vista m�s estad�stico)

\begin{entradilla}
{\em La {\color{introcolor}cantidad de informaci�n} que nos proporciona un mensaje es inversamente proporcional a su probabilidad}
\end{entradilla}

En caso de que no est�is acostumbrados a trabajar con logaritmos, que
sep�is que esa funci�n, el logaritmo, cumple la siguiente propiedad:


\begin{center}
$ log(\frac{1}{a}) = log (a ^ {-1}) = -1 \times log(a)$
\end{center}

En otras palabras, el logaritmo de un n�mero elevado a un determinado
exponente, es igual al exponente, multiplicado por el logaritmo del
n�mero. Como el inverso de cualquier n�mero es ese n�mero elevado a -1,
se cumple la expresi�n de arriba, y por ello \texttt{log(1/p)=-log(p)}

Os preguntar�is, por qu� el logaritmo?. La verdad que es una cuesti�n de
conveniencia, ya que las probabilidades, por definici�n, toman valores
entre 0 y 1. En otras palabras, se trata de valores muy peque�os. La
funci�n logaritmo tiene la propiedad de expandir los valores entre 0 y
1, de forma que resulta m�s sencillo trabajar con ellos. Por ejemplo,
imaginad el siguiente caso:

\begin{align*}
p(x) =0.0001 \\
\frac{1}{p(x)} = 10000 \\
log (\frac{1}{p(x)}) = -4 
\end{align*}


El logaritmo nos permite evitar n�meros muy peque�os o muy grandes y nos
ofrece valores m�s f�ciles de interpretar y por lo tanto m�s f�ciles de
manejar.

Bueno, es hora de continuar, ahora que podemos medir la cantidad de
informaci�n en un mensaje vamos a ver como podemos cuantificar muchos de
estos mensajes juntos.

\hypertarget{entropuxeda}{%
\sectiontext{white}{black}{ENTROP�A}\label{entropuxeda}}

La palabra \href{https://en.wikipedia.org/wiki/Entropy}{entrop�a} se
utiliza como un forma de medir el desorden de las cosas. En
termodin�mica nos indica como de revolucionados est�n los �tomos dentro
de un determinado sistema (y sistema termodin�micamente hablando es un
concepto muy amplio), pero
\href{https://en.wikipedia.org/wiki/Entropy_(information_theory)}{Entrop�a},
en teor�a de la informaci�n, se utiliza para representar lo aleatoria
que es una determinada secuencia de datos.

Dicho de otra forma, la entrop�a nos da una medida de la redundancia (o
de lo aleatorio si lo prefer�s) que hay en un mensaje y por lo tanto nos
indica cuanto se puede comprimir un determinado mensaje. Pensad por un
segundo en RLE o LZ77\ldots{} Ambos algoritmos se basan en buscar
valores repetidos y re-escribiros de una forma m�s eficiente. Es decir,
esos algoritmos intentan reducir la redundancia en el mensaje. Si no hay
redundancia\ldots{} bueno, los algoritmos simplemente no funcionan y en
general, lejos de comprimir los datos van a aumentar su tama�o\ldots{}
Que pasar�a si a LZ77 le pasamos un mensaje formado por una �nica
repetici�n de cada s�mbolo?\ldots{} No tendr�a nada que comprimir
verdad\ldots.

La entrop�a en teor�a de la informaci�n se calcula con la siguiente
f�rmula. Os adelanto que parece complicada pero en seguida ver�is que es
mucho m�s sencilla de lo que parece.

\begin{center}
$ H = \sum\limits_{i=o}^N p(x)log(\frac{1}{p(x)}) = - \sum\limits_{i=o}^N p(x)log(p(x))$
\end{center}

Vista as� da un poco de miedo la formulita no?. Pero y si la escribimos
as�:


\begin{center}
$ H = \sum\limits_{x} p(x)I(x)$
\end{center}

Donde, \texttt{p(x)} es la probabilidad de que ocurra el mensaje
\texttt{x}, e \texttt{I(x)} es la cantidad de informaci�n de ese
mensaje. Dicho de otra forma, la entrop�a no es m�s que la media
ponderada de la cantidad de informaci�n de todos los mensajes (o m�s
correctamente s�mbolos)q ue utilizamos en un determinado sistema de
comunicaci�n.



\begin{center}
{\colorbox{excolor}{
\begin{minipage}{.8\linewidth}
{\em
Normalmente cuando calcul�is una media supon�is que los valores tiene
todos la misma probabilidad, es decir, la distribuci�n es uniforme. En
ese caso la probabilidad de cada valor es 1/N, siendo N el n�mero de
elementos, As�, cuando hac�is una media, sum�is los valores y los
divid�s por N. En el caso general, en lugar de dividir por N hay que
multiplicar cada valor por su probabilidad asociada y as� es como
obtenemos la f�rmula anterior.}
\end{minipage}
}
}
\end{center}

\hypertarget{breve-comentario}{%
\sectiontext{white}{black}{BREVE COMENTARIO}\label{breve-comentario}}

Hasta ahora hemos estado utilizando las palabras mensaje y s�mbolo de
una forma un poco ambigua, ya que la diferencia no era realmente
importante, pero llegados a este punto es importante que conozcamos la
diferencia entre cada una de estas palabras.

Un s�mbolo es cada uno de los elementos m�nimos que podemos transmitir o
almacenar. El conjunto de todos los s�mbolos utilizados en un
determinado sistema de comunicaciones se denomina alfabeto. Los
s�mbolos, pueden ser cualquier cosa: N�meros, rayas y puntos, valores de
voltaje\ldots. Pero en general para nosotros cada s�mbolo va a ser un
valor normalmente codificado en binario.

En este contexto, un mensaje no es m�s que una secuencia de s�mbolos.
Espero no haberos liado mucho. Si pensamos en este articulo:

\begin{itemize}

\item
  Las letras son los s�mbolos utilizados para codificar el mensaje
\item
  El alfabeto\ldots{} bueno, es el alfabeto, la lista de las letras que
  podemos usar :).
\item
  Y el mensaje es el articulos que combina los s�mbolos del alfabeto de
  una forma concreta para transmitir cierta informaci�n.
\end{itemize}

En nuestros ejemplos anteriores, los s�mbolos que defin�amos, eran por
si mismo un mensaje. O si lo prefer�s, el n�mero de mensajes posible
estaba limitado y cada mensaje contenia solo un s�mbolo. Por ejemplo, si
nuestro mensaje es que esta lloviendo, utilizamos el s�mbolo D, pero no
tiene sentido a�adir ninguno de los otros. El mensaje y el s�mbolo
coinciden en ese caso.

Sin entrar en profundidad en otros aspectos del art�culo de Shannon,
espero que la idea general est� clara y resulte m�s sencillo comprender
el resto del art�culo.

\hypertarget{entropuxeda-y-redundancia}{%
\sectiontext{white}{black}{ENTROP�A Y REDUNDANCIA}\label{entropuxeda-y-redundancia}}

Tras esta peque�a disgresi�n, continuemos con el tema principal.

La entrop�a es capaz de medir la cantidad de redundancia que hay en una
determinada secuencia de s�mbolos o mensaje, y por lo tanto nos dar�,
entre otras cosas, una medida de cuanto podemos comprimir esos mensajes.
Volvamos por un momento a nuestro ejemplo de la lluvia y vamos a
inventarnos unos valores de ejemplo para tener algo con lo que trabajar.
Imaginemos que en nuestra ciudad imaginaria Razorburgo, llueve a lo
bestia unos 80 d�as al a�o, hace sol unos 181 d�as al a�o, 52 d�as al
a�o est� nublado y llueve de forma normal y el resto del a�o, unos 53
d�as, esta nublado y chispea. Esto nos dar�a las siguientes
probabilidades.

{\scriptsize
\begin{verbatim}
A : Caen Chuzos de Punta     | p(A) = 80/365  = 0.219
B : Llueve                   | p(B) = 52/365  = 0.142
C : Chispea                  | p(C) = 53/365  = 0.145
D : M�s seco que la mohama.  | p(D) = 181/365 = 0.496
\end{verbatim}
}

En este ejemplo, los mensajes que nos dan m�s informaci�n son los que
nos dice si llueve. Especialmente si llueve normal o chispea. Ahora
podemos calcular la entrop�a para este conjunto de s�mbolos.

\begin{verbatim}
H = - (  0.219 * log2(0.219)  +
         0.142 * log2 (0.142) +
         0.145 * log2 (0.145) +
         0.496 * log2 (0.496) = 1.785
\end{verbatim}

Comparemos esto con la que obtendr�amos en un lugar en el que llueve y
hace sol por igual. Es decir, la probabilidad de cualquier s�mbolo es
1/4:

\begin{center}
$ H = 4 \Big(\frac{1}{4} log_2(4)\Big) = 4 \frac{2}{4} = 2$
\end{center}

Sip, la entrop�a no solo nos dice la cantidad de redundancia que hay en
un mensaje, sino que tambi�n nos da el l�mite inferior de la cantidad de
bits que necesitamos para codificar cada s�mbolo de un mensaje. Bueno,
si hemos usado logaritmos en base dos. Si us�is logaritmos neperianos
(base e), obtendr�is NATs y si us�is logaritmos naturales (base 10)
obtendr�is HARTLEYs.

\begin{entradilla}
{\em {\color{introcolor}La entrop�a nos dice la cantidad de redundancia} que hay en un determinado mensaje}
\end{entradilla}

Observad que si definimos un conjunto de mensajes equi-probables (todos
con la misma posibilidad, la formula se reduce a:

\begin{center}
$H = N \Big( \frac{1}{N} log_2(\frac{1}{N})\Big) = - log_2 (N)$
\end{center}
Y eso no es otra cosa que el n�mero de bits necesarios para codificar un
s�mbolos de cualquier alfabeto con s�mbolos equiprobables\ldots. Mind
blowing!!!!

\hypertarget{teorema-de-codificaciuxf3n-de-fuente}{%
\sectiontext{white}{black}{TEOREMA DE CODIFICACI�N DE FUENTE}\label{teorema-de-codificaciuxf3n-de-fuente}}

Esto que os acabamos de contar es la base de lo que se conoce como el
teorema de codificaci�n de fuente de Shannon, y puede ser demostrado
matem�ticamente. El teorema en su formulaci�n original dice lo
siguiente.

\emph{Dadas N variables aleatorias independientes e id�nticamente
distribuidas, cada una de ellas con entrop�a H(X), estas N variables
pueden ser comprimidas en N * H(X) bits con un riesgo despreciable de
p�rdida de informaci�n cuando N tiende a infinito. Por el contrario, si
esas variables aleatorias son comprimidas en menos de N * H(x) bits es
virtualmente seguro que se perder� informaci�n.}

Es decir, la entrop�a no solo nos permite determinar lo aleatorio que es
un mensaje sino que nos indica cual es el l�mite te�rico de compresi�n
de una determinada fuente de datos.

Como os contamos en la secci�n anterior, la entrop�a nos indica el
n�mero m�nimo de bits que necesitamos por s�mbolo mientras que el
teorema de codificaci�n de fuente habla de mensajes de N s�mbolos. Por
eso el tama�o total del que habla est� multiplicado por \texttt{N}.

Ahora que sabemos cuanto podemos comprimir una determinada fuente de
datos, solo necesitamos saber como codificar esos datos para acercarnos
lo m�s posible a el valor te�rico que nos proporciona el teorema.
Observad que el teorema asegura que podemos conseguir ese nivel de
compresi�n cuando el mensaje tiene una longitud infinita, o si lo
prefer�s cuando es m� grande. Lo que esto significa es que, en la
pr�ctica, la tasa de compresi�n que obtendremos ser� un poco superior a
la te�rica\ldots{} Pero para explorar todo esto necesitamos ser capaces
de codificar nuestros datos de forma que ocupen menos. Vamos all�.

\hypertarget{cuxf3digos-de-longitud-variable-1}{%
\sectiontext{white}{black}{C�DIGOS DE LONGITUD VARIABLE\\OTRA VEZ}\label{cuxf3digos-de-longitud-variable-1}}

Ahora que ya hemos introducido los principales elementos que necesitamos
para entender como comprimir mensajes desde el punto de vista de la
teor�a de la informaci�n, vamos a ver como usar todo esto en la
pr�ctica. Lo que vamos a hacer es algo bastante obvio: \emph{Usar menos
bits para codificar los mensajes que aparezcan m�s a menudo}. O lo que
es lo mismo, utilizar c�digos de longitud variable en los que los
s�mbolos no tienen todos el mismo n�mero de bits. Que es exactamente
otra forma de decir lo mismo :).

\begin{entradilla}
{\em El {\color{introcolor}c�digo Morse} es quiz�s el c�digo de longitud variable m�s conocido.}
\end{entradilla}

Un ejemplo de esto es el conocido
\href{https://en.wikipedia.org/wiki/Morse_code}{c�digo Morse}, el cual
asigna secuencias de puntos y rayas m�s cortas a las letras m�s
comunes\ldots{} al menos en el idioma ingl�s.

Para ilustrar todo esto, usaremos un ejemplo muy sencillo con 4
s�mbolos: A, B, C, D.

Empecemos por el caso m�s sencillo en el que todos los s�mbolos tienen
la misma posibilidad de aparecer:

\begin{verbatim}
A | p(A) = 1/4 | log(p(A)) = -2 | 00
B | p(B) = 1/4 | log(p(B)) = -2 | 01
C | p(C) = 1/4 | log(p(C)) = -2 | 10
D | p(D) = 1/4 | log(p(D)) = -2 | 11
\end{verbatim}

La entrop�a para este alfabeto es de 2 (la media de muchos valores
iguales es precisamente el valor de cada uno de ellos\ldots{} no
necesitamos calcularlo). Y tiene sentido que necesitemos 2 bits para
codificar nuestros 4 mensajes. Si ninguno de los mensajes es m�s
probable que los dem�s, no hay ninguna informaci�n que podamos utilizar
as� que la mejor soluci�n es utilizar nuestro cl�sico c�digo de longitud
fija.

Ahora imaginemos que las probabilidades de los s�mbolos son diferentes.
Pensad por ejemplo en las letras en un texto, algunas letras aparecen
m�s a menudo que otras o, dicho de otra forma, algunas letras tienen una
probabilidad m�s alta de aparecer en un mensaje.

\begin{verbatim}
A | p(A) = 1/2 | log(p(A)) = -1 
B | p(B) = 1/4 | log(p(B)) = -2 
C | p(C) = 1/8 | log(p(C)) = -3 
D | p(D) = 1/8 | log(p(D)) = -3 
\end{verbatim}

La entrop�a para este alfabeto es: 0.4375. Vamos que estamos
desperdiciando un mont�n de espacio utilizando una n�mero fijo de bits
para codificar esta informaci�n.

As� que vamos a intentar asignar menos bits a los mensajes m�s probables
(los que aparecen m�s a menudo), y dar m�s bits a los menos probables.
Comenzaremos d�ndole el c�digo \texttt{0} al mensaje \texttt{A}. El
mensaje \texttt{B} es el siguiente, pero ya no le podemos dar solo un
bit, porque sino, no podr�amos decodificar el mensaje, as� que le vamos
a dar el c�digo \texttt{10} (veremos por qu� en un momento).

Finalmente, los s�mbolos \texttt{C} y \texttt{D} son los menos probables
y por lo tanto deben tener el c�digo m�s largo, as� que les asignaremos
los valores \texttt{110} y \texttt{111} respectivamente. Nuestro c�digo
resultante es:

\begin{verbatim}
A | 0
B | 10
C | 110
D | 111
\end{verbatim}

\emph{Fijaros que el n�mero de bits que hemos asignado a cada s�mbolo se
corresponde con la cantidad de informaci�n asociada a cada uno de ellos;
1, 2, 3 y 3.}

Este c�digo no lo hemos creado de forma arbitraria, es un tipo especial
de c�digo conocido como
\href{https://en.wikipedia.org/wiki/Prefix_code}{c�digo prefijo}. Un
c�digo prefijo es un c�digo de longitud variable en el que cada palabra
c�digo no es prefijo de ninguna otra. Lo que esto nos permite es poder
decodificar los mensajes cuando recibimos los bits todos juntos sin
ambiguedad posible. Por ejemplo, dada la secuencia de
bits\texttt{0001101111011101110}, obtendremos el resultado:

\begin{verbatim}
0 0 0 110 111 10 111 0 111 0
A A A C   D   B  D   A D   A
\end{verbatim}

Como pod�is ver, esa es la �nica forma de decodificar la secuencia de
bits inicial.

Utilizando el teorema de codificaci�n de fuente de Shannon, sabemos que,
para este conjunto de s�mbolos, es posible codificar un mensaje
suficientemente largo con \texttt{N\ *\ H(X)}. M�s arriba vimos que la
entrop�a de este alfabeto era de 0.4375. Puesto que tenemos 4 s�mbolos
el n�mero m�nimo de bits por s�mbolo te�rico que podemos alcanzar es:

\begin{verbatim}
N * H(X) = 4 * 0.4375 = 1.75
\end{verbatim}

Ahora calculemos la longitud media del c�digo de longitud variable que
hemos generado:

\begin{verbatim}
E{L} = 1*p(A) + 2*p(B) + 3*p(C) + 3*p(D) =
     = 1*1/2  + 2*1/4  + 3*1/8  + 3*1/8  =  1.75
     
\end{verbatim}

Es decir, el c�digo que hemos calculado es el c�digo �ptimo para este
conjunto de s�mbolos y tiene un nombre especial\ldots{}

\hypertarget{cuxf3digos-de-huffman}{%
\sectiontext{white}{black}{C�DIGOS DE HUFFMAN}\label{cuxf3digos-de-huffman}}

El c�digo que hemos construido en la secci�n anterior es un tipo
especial de c�digo prefijo conocido como
\href{https://en.wikipedia.org/wiki/Huffman_coding}{c�digo de Huffman}.
El c�digo de Huffman es �ptimo para un determinado alfabeto y se puede
generar muy f�cilmente utilizando un sencillo algoritmo (b�sicamente
tenemos que construir un �rbol binario usando las probabilidades):

Comenzamos colocando todos los s�mbolos en una tabla ordenada por sus
probabilidades:

\begin{verbatim}
S�mbolo |    A   |    B   |    C  |    D  |
--------+--------+--------+-------+-------+
Prob    |   .5   |   .25  | 0.125 | 0.125 |
\end{verbatim}

A partir de aqu�, construiremos el c�digo de forma iteractiva, a�adiendo
un bit cada vez. Comenzamos eligiendo los dos s�mbolos con la
probabilidad m�s baja. En este caso \texttt{C} y \texttt{D} y asignamos
un bit a cada uno de ellos.

\begin{verbatim}
S�mbolo |    A   |    B   |    C  |    D  |
--------+--------+--------+-------+-------+
Prob    |   .5   |   .25  | 0.125 | 0.125 |
        |        |        |   0   |   1   |
\end{verbatim}

A continuaci�n, asignamos al conjunto formado por esos dos s�mbolos la
suma de las probabilidades de los s�mbolos originales,

\begin{verbatim}
S�mbolo |    A   |    B   |    C  |    D  |
--------+--------+--------+-------+-------+
Prob    |   .5   |   .25  | 0.125 | 0.125 |
        |        |        |   0   |   1   |
        |   .5   |   .25  | 0.25  |       |
\end{verbatim}

Y repetimos el proceso a�adiendo cada vez bits a la izquierda.

\begin{verbatim}
S�mbolo |  A  |  B  |    C  |    D  | PASO |
--------+-----+-----+-------+-------+------+
Prob    | .5  | .25 | 0.125 | 0.125 |  0   |
        |     |     |   0   |   1   |  1   |
        | .5  | .25 | 0.25  |       |  2   |
        |     |  0  |  10   |   11  |  3   |
        | .5  | .5  |       |       |  4   |
        |  0  | 10  |  110  |  111  |  5   |
\end{verbatim}

Realmente hemos construido un �rbol binario, pero empezando desde abajo.

\begin{verbatim}
   A       B      C      D
 .5|    .25|  .125|  .125|    0
   |       |   '0'|   '1'|    1
   |       |      +--+---+
 .5|    .25|       .5|        2
   |    '0'|      '1'|        3
   |       +----+----+
 .5|          .5|             4
'0'|         '1'|             5
   +------+-----+
          | ROOT
\end{verbatim}

Ahora, si recorremos el �rbol desde la ra�z a cada uno de los s�mbolos,
obtenemos los c�digos que obtuvimos anteriormente. El hecho de que se
trate de un �rbol binario, nos asegura que se trata de un c�digo
prefijo\ldots{} simplemente por la forma de construirlo.

Como pod�is observar, dependiendo de como decidamos ir asignando los
distintos bits, podemos obtener diferentes c�digos. No solo eso, en
casos m�s complejos en los que varios s�mbolos tienen la misma
probabilidad, podremos agruparlos de diferentes formas y obtener c�digo
completamente distintos. Eso hace que un mismo alfabeto pueda producir
diferentes c�digos Huffman, dependiendo de como se asignen los bits.

Veamos ahora como quedar�a codificado un mensaje de ejemplo, utilizando
nuestro c�digo Huffman u un c�digo de longitud fija:

{\scriptsize
\begin{verbatim}
A  A  A  B  A  A  A  B  A  B  A  C   A  A  B  D   A  : 17 S�mbolos
00 00 00 01 00 00 00 01 00 01 00 10  00 00 01 11  00 : 34 bits
0  0  0  10 0  0  0  10 0  10 0  110 0  0  10 111 0  : 25 bits
\end{verbatim}
}
Pues parece que nos hemos ahorrado unos cuantos bits. En concreto
necesitamos 25 bits para codificar 17 s�mbolos, lo que significa que
hemos utilizado 1.47 bits por s�mbolo y hemos comprimido el mensaje algo
m�s de un 25\%.

\begin{entradilla}
{\em Los c�digos Huffman se utilizan junto con otros algoritmos como {\color{introcolor}LZW, JPEG o el formato Deflate}}
\end{entradilla}

Los c�digos Huffman son m�s comunes de lo que podr�ais imaginas, si
bien, no se utilizan como sistema de compresi�n, sino como una forma
eficiente de codificar informaci�n normalmente como parte de un
algoritmo de compresi�n m�s general. Por ejemplo, muchas
implementaciones de LZ77 utilizan c�digos Huffman para codificar las
longitudes de las secuencias. JPEG, utiliza c�digos Huffman para
codificar los coeficientes de la DCT (realmente no es un c�digo Huffman,
sino c�digo prefijo pre-calculado, pero se suelen referir a ellos de
esta forma).

\hypertarget{extensiuxf3n-de-fuente}{%
\sectiontext{white}{black}{EXTENSI�N DE FUENTE}\label{extensiuxf3n-de-fuente}}

Una forma de comprimir todav�a m�s un mensaje es utilizando una t�cnica
que se conoce como extensi�n de fuente. El concepto es muy sencillo, en
lugar de tomar s�mbolos aislados, los tomamos en grupos. Veamos como
funciona esto con un sencillo ejemplo utilizando solo 2 s�mbolos.
Imaginemos que nuestro sistema de comunicaci�n utiliza los dos s�mbolos
siguientes:

\begin{verbatim}
A -> p(A) = 3/4
B -> p(B) = 1/4
\end{verbatim}

La entrop�a para este conjunto de s�mbolos es \texttt{0.81}, es decir,
te�ricamente podr�amos comprimir nuestros mensajes para que, en media,
cada s�mbolo utilice \texttt{0.81} bits\ldots{} Pero con solo dos
s�mbolos, los �nicos c�digos que podemos generar son

\begin{verbatim}
     C�DIGO 1 | C�DIGO 2
A ->     0    |    1
B ->     1    |    0
\end{verbatim}

Entonces, como podemos acercarnos al l�mite te�rico de Shannon indicado
por la entrop�a. Bueno, una forma es extendiendo la fuente. En lugar de
considerar los s�mbolos \texttt{A} y \texttt{B}, consideremos todas sus
combinaciones:

\begin{verbatim}
AA -> p(A) * p(A) = 9/16
AB -> p(A) * p(B) = 3/16
BA -> p(B) * p(B) = 3/16
BB -> p(B) * p(B) = 1/16
\end{verbatim}

Ahora podemos generar un nuevo c�digo Huffman con estos nuevos valores
de probabilidad:

\begin{verbatim}
AA -> p(AA) = 9/16  -> 0
AB -> p(AB) = 3/16  -> 10
BA -> p(BB) = 3/16  -> 110
BB -> p(BB) = 1/16  -> 111
\end{verbatim}

En toda esta teor�a siempre se asumen que los valores son independientes
e id�nticamente distribuidos. Esto simplemente significa que todos los
valores siguen la misma distribuci�n de probabilidades y su valor no
depende de ning�n otro valor. Esta es la raz�n por la que podemos
calcular la probabilidad de dos s�mbolos consecutivos como el producto
de las probabilidades de cada s�mbolo separado\ldots{} Cuando los
valores (la variable aleatoria como le llaman en estad�stica) no cumplen
esta propiedad\ldots{} bueno, las cosas se complican un poco m�s.

Si ahora calculamos la entrop�a y longitud media del c�digo obtendremos:

\begin{verbatim}
H = 1.62
L = 9/16*1 + 3/16*2 + 3/16*3 + 1/16*4 = 1.68
\end{verbatim}

Como pod�is ver, ahora estamos mucho m�s cerca del valor te�rico. De
hecho, el cociente entre la entrop�a y la longitud media del c�digo mide
la eficiencia del c�digo. Para el primer caso ten�amos:

\begin{verbatim}
0.81/1 = 0.81 -> 81%
\end{verbatim}

Mientras que ahora nuestro c�digo tiene una eficiencia de:

\begin{verbatim}
1.62/1.68 = 0.96 -> 96%
\end{verbatim}

Observad que ahora, cada uno de nuestros s�mbolos incluye dos de los
s�mbolos originales. As� que, respecto en media, estamos utilizando
\texttt{1.68/2\ =\ 0.84} bits por cada uno de nuestos s�mbolos
originales (\texttt{A} y \texttt{B}) y como pod�is ver, este es un valor
m�s cercano al valor de la entrop�a original.

\begin{entradilla}
{\em El {\color{introcolor}proceso de extension de fuente} nos permite acercarnos al {\color{introcolor}l�mite te�rico de compresi�n} que marca el teorema de Shannon}
\end{entradilla}

Este proceso se puede extender varias veces (probad a generar el c�digo
para la extensi�n a 3 simbolos), si bien, la ganancia en compresi�n es
cada vez menor con cada extensi�n. Si lo pensamos un poco, el algoritmo
LZ del que hablamos al principio, en cierto modo, saca beneficio de este
resultado ya que agrupa s�mbolos asignandole c�digos m�s cortos cuanto
m�s frecuentes son (cuanto m�s aparezcan en el mensaje a comprimir).

\hypertarget{conclusiuxf3n}{%
\sectiontext{white}{black}{CONCLUSIONES}\label{conclusiuxf3n}}

En este art�culo hemos repasado algunos de los m�todos m�s comunes para
comprimir datos sin p�rdidas. Si bien existen muchos m�s tipos de
algoritmos de compresi�n, estos tres que hemos discutido cubren un
amplio espectro de los existentes. Los c�digos Huffman basados en
entrop�a, los c�digos LZ basados en diccionario de palabras clave y
otros algoritmos m�s intuitivos como RLE que no caen en una familia
concreta.

Nos hemos introducido t�midamente en el apasionante mundo de la teor�a
de la informaci�n y hemos visto como conceptos como la entrop�a y la
longitud de los c�digo se relacionan estad�sticamente con los datos que
queremos comprimir

\end{multicols}


\hypertarget{geek:bf1}{}\label{geek:bf}
\pagestyle{geek}

\rput(6.0,-1.7){\resizebox{!}{11.5cm}{{\epsfbox{images/header-geek.eps}}}}

\psset{fillstyle=solid}
\psframe[fillcolor=black,opacity=0.4](2,-4.5)(17,0)


% -------------------------------------------------
% Cabecera
\begin{flushright}


{\color{violet}\mtitle{12cm}{Interprete BrainFuck. Parte I}}

\msubtitle{6cm}{En ensamblador RISC-V}

{\sf\color{white}{ por Don Bit0}}

%{\psset{linecolor=black,linestyle=dotted}\psline(-12,0)}
\end{flushright}

\vspace{2mm}
% -------------------------------------------------



%\lstset{language=C,frame=tb,framesep=5pt,basicstyle=\footnotesize}
\lstset{language=C,frame=tb,framesep=5pt,basicstyle=\scriptsize}




\intro{introcolor}{I}{niciamos una nueva secci�n en la que vamos a hacer cosas muy geeks y
empezamos con un art�culo que creemos reflejar� muy bien la filosof�a de
esta secci�n. Vamos a programar un int�rprete del lenguaje esot�rico
\emph{BrainFuck} en ensamblador RISC-V\ldots{} por que?\ldots{} pues
porque podemos :)}

\begin{multicols}{2}
Para los que no conozc�is brainfuck, bueno enseguida entender�is de
donde viene el nombre. Es quiz�s el lenguaje esot�rico m�s conocido y
emblem�tico, adem�s de raro de narices. Para los que no conozc�is
RISC-V, pues se trata de una arquitectura de procesadores RISC que tiene
la peculiaridad de estar libre de licencias, lo que permite a cualquier
desarrollar libremente su propio procesador basado en esa arquitectura.

Arquitecturas como ARM por ejemplo requiere que los desarrolladores de
procesadores basados en ARM paguen una licencia. Los usuarios finales no
nos enteramos, pero para los desarrolladores de HW es un coste a tener
en cuenta.

Sin m�s pre�mbulos vamos a contaros que es BrainFuck, las peculiaridades
de RISC-V y ponernos al tajo. Vamos a ir construyendo el programa con
pasos reproducibles que pod�is seguir y que os ayuden a depurar el
programa en caso de que teng�is cualquier problema\ldots{} Y como
siempre, no dud�is en poneros en contacto con nosotros ;)

\hypertarget{brainfuck}{%
\sectiontext{white}{black}{BRAINFUCK}\label{brainfuck}}

BrainFuck es uno de los lenguajes esot�ricos m�s emblem�ticos.
Probablemente el m�s conocido. Si quer�is conocer m�s sobre su historia,
la p�gina de la \href{https://en.wikipedia.org/wiki/Brainfuck}{wikipedia}
contiene toda la informaci�n
que pod�is buscar. En este art�culo nosotros solo vamos a explicaros la
sintaxis del lenguaje de forma que el art�culo sea auto-contenido.

\begin{entradilla}
{\em {\color{introcolor}BrainFuck} es uno de los {\color{introcolor}lenguajes esot�ricos} m�s conocidos}
\end{entradilla}

El lenguaje est� compuesto por 8 caracteres que funcionan sobre una,
digamos, m�quina virtual muy sencilla. Esta m�quina define una zona de
memoria en la que se almacena el programa (una secuencia de los 8
caracteres que describiremos a continuaci�n), y otra zona de memoria en
la que se almacenan datos.

La m�quina dispone de dos registros. Un contador de programa, que apunta
a la instrucci�n del lenguaje que estamos ejecutando (en la primera zona
de memoria) y un puntero de datos que apunta a alguna de las posiciones
disponibles en el segundo bloque de datos.

Veamos las instrucciones y un programa de ejemplo:

\begin{lstlisting}
> Incrementa el puntero de datos en 1
< Decrementa el puntero de datos en 1
+ Incrementa el valor apuntado por el puntero de datos
- Decrementa el valor apuntado por el puntero de datos
. Muestra en pantalla el valor apuntado por el puntero de
  datos
, Lee un valor del teclado y lo almacena en la posici�n
  indicada por el puntero de datos
[ Si el valor apuntado por el puntero de datos es 0 salta al ]
  correspondiente. Sino continua la ejecuci�n normalmente
] Si el valor apuntado por el puntero de datos no es cero,
  salta al [ correspondiente, sino continua la ejecuci�n
  normalmente.
\end{lstlisting}

En contra de lo que pueda parecer el lenguaje es mucho m�s sencillo de
lo que parece, si bien, hacer cosas �tiles es complicado. Veamos un
sencillo programa.

\hypertarget{ejecutando-un-programa-brainfuck}{%
\sectiontext{white}{black}{EJECUTANDO UN PROGRAMA BF PASO A PASO}\label{ejecutando-un-programa-brainfuck}}

Este va a ser nuestro programa de ejemplo. Como pod�is ver utiliza todos
los comandos, as� que nos va a permitir ver como funciona el lenguaje en
su totalidad.

\begin{lstlisting}
++>,.<[>+.<-]
\end{lstlisting}

Vamos a ir ejecutando el programa caracter a caracter para que entend�is
como funciona cada uno de los comandos. Los caracteres
{\verb!\^!} indican la posici�n a la que apunta el
puntero de programa (PC) y el puntero de datos (PD):

\begin{lstlisting}
PRG   | ++>,.<[>+.<-]
PC    | ^ 
DATOS | 0 0 0 0 0 0 0 ....
PD    | ^
\end{lstlisting}

Al comenzar el programa encontramos dos caracteres
{\verb!+!}. Estos caracteres incrementan el valor al
que apunta el puntero de datos. As� que tras ejecutar las dos primeras
instrucciones la m�quina virtual tendr� el siguiente estado:

\begin{lstlisting}
PRG   | ++>,.<[>+.<-]
PC    |   ^ 
DATOS | 2 0 0 0 0 0 0 ....
PD    | ^
\end{lstlisting}

Hemos almacenado el valor dos en la primera posici�n de memoria. Este va
a ser nuestro contador para el bucle que implementamos un poco m�s
adelante. Ahora nos encontramos un caracter {\verb!>!},
el cual nos mueve a la siguiente posici�n de memoria:

\begin{lstlisting}
PRG   | ++>,.<[>+.<-]
PC    |    ^ 
DATOS | 2 0 0 0 0 0 0 ....
PD    |   ^
\end{lstlisting}

En esta nueva posici�n vamos a leer un caracter de la consola e
imprimirlo inmediatamente. Eso es lo que hacen los comandos
{\verb!,!} y {\verb!.!}. Supongamos
que hemos introducido el caracter {\verb!A!}. La
m�quina virtual mostrar� el siguiente estado:

\begin{lstlisting}
PRG   | ++>,.<[>+.<-]
PC    |      ^ 
DATOS | 2 64 0 0 0 0 0 ....
PD    |    ^
  
\end{lstlisting}

Estamos a punto de comenzar el bucle, as� que primero nos movemos a la
posici�n de memoria que contiene nuestro contador de bucle. Esto lo
hacemos con el comando {\verb!<!} con el que nos
moveremos a la posici�n inicial.

\begin{lstlisting}
PRG   | ++>,.<[>+.<-]
PC    |       ^ 
DATOS | 2 64 0 0 0 0 0 ....
PD    | ^
\end{lstlisting}

Ahora estamos en condiciones de ejecutar nuestro bucle. Puesto que el
valor al que apunta el puntero de datos no es cero, continuamos la
ejecuci�n normalmente entrando en el bucle. Dentro del bucle haremos dos
cosas, incrementar el valor que hemos le�do por consola y decrementar el
contador del bucle:

\begin{lstlisting}
PRG   | ++>,.<[>+.<-]
PC    |        ^ 
DATOS | 2 64 0 0 0 0 0 ....
PD    |   ^
\end{lstlisting}

Ahora que apuntamos a la {\verb!A!}, sumamos uno e
imprimimos.

\begin{lstlisting}
PRG   | ++>,.<[>+.<-]
PC    |          ^ 
DATOS | 2 65 0 0 0 0 0 ....
PD    |   ^
\end{lstlisting}

Ahora apuntamos a nuestro contador y lo de crementamos

\begin{lstlisting}
PRG   | ++>,.<[>+.<-]
PC    |           ^ 
DATOS | 2 65 0 0 0 0 0 ....
PD    | ^

PRG   | ++>,.<[>+.<-]
PC    |             ^ 
DATOS | 1 65 0 0 0 0 0 ....
PD    | ^
\end{lstlisting}

Llegados a este punto nos encontramos el caracter
{\verb!]!}, el cual, si el valor actual de los datos no
es cero, saltar� a la posici�n del caracter {\verb![!}
correspondiente. En este caso el valor al que estamos apuntando es
{\verb!1!}, as� que volvemos hacia atr�s.

\begin{lstlisting}
PRG   | ++>,.<[>+.<-]
PC    |       ^       
DATOS | 1 65 0 0 0 0 0 ....
PD    | ^
\end{lstlisting}

Y procedemos a ejecutar una vez m�s el cuerpo del bucle.

\begin{lstlisting}
PRG   | ++>,.<[>+.<-]
PC    |             ^       
DATOS | 0 66 0 0 0 0 0 ....
PD    | ^
\end{lstlisting}

Tras lo cual volveremos a comprobar nuestro contador, el cual ahora vale
{\verb!0!} y por lo tanto el bucle termina.

Este sencillo programa imprime el caracter que introducimos por teclado
y los dos siguientes.

Como pod�is ver, lo que hacen los programas es super-sencillo, pero
programar cualquir cosa �til es muy complicado.

En nuestras pruebas, vamos a utilizar este programa que acabamos de
mostrar, adem�s de el \emph{Hola Mundo} para la prueba final. Aqu� os lo
dejamos para que os divirt�is un poco:

\begin{lstlisting}
++++++++[>++++[>++>+++>+++>+<<<<-]>+>+>->>+[<]<-]>>.>---.
+++++++..+++.>>.<-.<.+++.------.--------.>>+.>++.
\end{lstlisting}

\hypertarget{risc-v}{%
\sectiontext{white}{black}{RISC-V}\label{risc-v}}

Hablemos ahora brevemente del la arquitectura RISC-V, en concreto la
versi�n de 64 bits. La arquitectura RISC-V es, como su propio nombre
indica una arquitectura {\verb!RISC!} (\emph{Reduced
Instruction Set Computer} u \emph{Ordenadores con Conjunto de
Instrucciones Reducido}). Este tipo de procesadores normalmente ofrecen
muy pocas instrucciones, pero es capaz de ejecutarlas muy r�pido. Otros
procesadores como los Intel tienen una arquitectura
{\verb!CISC!}, donde la {\verb!C!}
significa complejo y ofrecen varios cientos de instrucciones con tiempos
de ejecuci�n muy variados.

\begin{entradilla}
{\em {\color{introcolor}La plataforma RISC-V} guarda muchas similitudes con los procesadores MIPS.}
\end{entradilla}

Tanto los nombres de los registros como los de las instrucciones se
corresponden casi uno a uno con los procesadores MIPS, sin embargo hay
algunas diferencias peque�as diferencias entre ellos. Vamos, que si
sab�is ensamblador MIPS tambi�n sab�is ensamblador RISC-V.

Si no sab�is ninguno, no os preocup�is, iremos explicando cada una de
las instrucciones que usemos en nuestro programa.

\hypertarget{risc-v-registros}{%
\sectiontext{white}{black}{RISC-V REGISTROS}\label{risc-v-registros}}

La arquitecura RISC-V define 32 registros de los cuales los 5 primeros
tienen funciones especiales, y el resto se agrupan de acuerdo a la forma
en la que se usan normalmente, si bien, esto es solo una convenci�n. Los
registros se nombran por su �ndice o por un nombre simb�lico, al igual
que sucede con los procesadores MIPS.

\begin{lstlisting}
x0      zero    Siempre devuelve 0
x1      ra      Direcci�n de Retorno
x2      sp      Puntero de pila
x3      gp      Puntero Global
x4      tp      Puntero de Hebra
x6      t0      Temporal/Direcci�n Retorno Alternativa
x6-7    t1-2    Registros temporales
x8      s0/fp   Registros Guardados/Frame Pointer
x9      s1      Registro Guardado
x10-11  a0-a1   Parametros Funciones/Valores Retorno
x12-18  a2-7    Par�metros Funciones
x18-27  s2-11   Registros Guardados
x28-31  t3-t6   Temporales
\end{lstlisting}

En nuestro programa nosotros usaremos los registros
{\verb!a0-a7!} para hacer las llamadas al sistema y los
registros {\verb!t0-t6!} para almacenar nuestros
valores temporales.

En nuestro caso no vamos a utilizar funciones, as� que, en realidad,
podemos utilizar los registros que nos de la gana, pero en un programa
normal, tendremos varias funciones y unas llamar�n a otras. Los
registros que aparecen en la lista como
{\verb!Registros Guardados!} deben ser guardados por la
funci�n que es llamada. Es decir nuestra funci�n no debe modificar esos
registros, puesto que el c�digo que llame a nuestra funci�n, va a
esperar que esos registros se preserven. En otras palabras, si queremos
usar esos registros en un programa normal, debemos guardar sus valores
originales y restaurarlos antes de volver de la funci�n.

Como os adelantamos, en nuestra implementaci�n no vamos a tener ese
problema, pero en caso de que alguien se preguntara esto\ldots{} bueno,
pues ah� lo ten�is.

\hypertarget{bf01.-leyendo-el-programa-en-memoria}{%
\sectiontext{white}{black}{LEYENDO PROGRAMA EN MEMORIA}\label{bf01.-leyendo-el-programa-en-memoria}}

Comenzaremos haciendo un programa muy sencillo que lea la entrada del
usuario, y almacene esos datos en memoria. El c�digo ser�a algo as�:

\begin{lstlisting}
    .text
    .global _start
    .equ MEM_SIZE, 4096
_start:
    # Lee datos desde entrada standard
    la   a1, prog  # Buffer de lectura
    move t1, zero  # Contador longitud del programa
bucle01:    
    li   a7, 63   # SYS_read
    li   a0, 0    # stdin
    li   a2, 1    # leemos caracter a caracter
    # a0 = read (a7=SYS_read, a0=stdin, a1=prog, a2=1);
    ecall       
    blez a0, cont01 # Si a0 <= 0 dejamos de leer
    
    add  t1, t1, 1 # Contamos un caracter
    add  a1, a1, 1 # Actualizamos el pumtero
    j    bucle01     # Repetimos
cont01:
    # Muestra la cadena leida por pantalla
    li   a7, 64     # SYS_write
    li   a0, 1      # a0 = stdout
    la   a1, prog   # a1 = prog
    move a2, t1   # a2 = size
    # write (a7=SYS_write, a0=stdout, a1=prog, a2=size);
    ecall         
exit:
    li   a7, 93   # Exit  
    li   a0, 0
    ecall
prog:
    .fill MEM_SIZE  
\end{lstlisting}

Para los que no sep�is mucho de programaci�n en ensamblador en
GNU/Linux, vamos a daros un curso intensivo con este programa. Para los
que ya sep�is\ldots{} bueno, os pod�is saltar la siguiente secci�n ;).

\hypertarget{curso-intensivo-de-ensamblador-gnulinux}{%
\sectiontext{white}{black}{CURSO INTENSIVO DE ENSAMBLADOR GNU/Linux}\label{curso-intensivo-de-ensamblador-gnulinux}}

Para este cursillo intensivo, simplemente vamos a ir comentando las
distintas partes del programa que os acabamos de mostrar. Vamos a ello:

\begin{lstlisting}
    .text
    .global _start
_start:
\end{lstlisting}

As� suelen empezar todos los programas en ensamblador,
independientemente del procesador. La primera l�nea
{\verb!.text!} indica que lo que sigue es c�digo (no
datos). La segunda declara un s�mbolo global llamado
{\verb!_start!}. Este es el s�mbolo que va a utilizar
el linker como nuestro punto de entrada\ldots{} o dicho de otra forma,
es donde el programa va a empezar a ejecutarse.

A continuaci�n tenemos la etiqueta {\verb!_start:!}
que indica la direcci�n en la que empezar� nuestro programa.

\hypertarget{llamadas-al-sistema-en-risc-v}{%
\sectiontext{white}{black}{LLAMADAS AL SYSTEMA CON RISC-V}\label{llamadas-al-sistema-en-risc-v}}

Para poder hacer cosas como leer datos del usuario o imprimir mensajes
en pantalla debemos hacer llamadas al sistema operativo o
{\verb!syscalls!}. Para cada sistema operativo esto se
hace de una determinada forma y, dentro de cada sistema operativo, cada
arquitectura usa un m�todo distinto.

\begin{entradilla}
{\em Podemos ejecutar llamadas al sistema en Linux siguiendo {\color{introcolor}el ABI System V}}
\end{entradilla}

Para UNIX se suele utilizar lo que se conoce como System V ABI. ABI es
el acr�nimo de \emph{Application Binary Interface}\ldots{} es como un
API, pero en vez de decirte que funciones puedes utilizar, te dice como
llamar a esas funciones (entre otras cosas). Linux utiliza esta ABI y en
general funciona de la siguiente forma:

\begin{itemize}

\item
  Almacenamos en un registro el n�mero de syscall que queremos invocar.
\item
  Almacenamos en otros registros los par�metros que queremos pasar a la
  syscall
\item
  Utilizamos un comando especial para pasar a modo kernel
\item
  Recibimos en un registro (o dos) el resultado de la operaci�n.
  Normalmente se usa el mismo registro que el utilizado para indicar la
  syscall o el primer par�metro.
\end{itemize}

Para el caso de RISC-V:

La forma de ejecutar estar esta {\verb!syscall!} ser�a:

\begin{itemize}

\item
  Almacenamos el n�mero de syscall en registro
  {\verb!a7!}
\item
  Pasamos los par�metros en los registros
  {\verb!a1-a6!}
\item
  Ejecutamos la instrucci�n {\verb!ecall!}
\item
  Procesamos el resultado devuelto en {\verb!a0!}
\end{itemize}

Poniendo como ejemplo la llamada al sistema
{\verb!SYS\_READ!} (que tiene valor
{\verb!63!} para RISC-V) que tiene el siguiente
prototipo C:

\begin{lstlisting}[language=C]
 ssize_t read(int fd, void buf[.count], size_t count);
\end{lstlisting}

El siguiente c�digo leer�a un caracter de la entrada est�ndar
(descriptor {\verb!0!}):

\begin{lstlisting}
  li   a7, 63     # SYS_read
  move a0, zero   # stdin
  la   a1, buffer # Puntero a nuestro buffer de lectura
  li   a2, 1      # leemos un caracter
  ecall           # SYS_read (a0=0, a1=buffer, a2=1)
\end{lstlisting}

Expliquemos brevemente las instrucciones que hemos utilizado:

\begin{itemize}

\item
  {\verb!li!} (\emph{Load Immediate}). Permite cargar
  un valor en registro.
\item
  {\verb!la!} (\emph{Load Address}). Permite cargar una
  direcci�n en un registro. Esta es una pseudo instrucci�n
\item
  {\verb!move!}. Permite mover valores entre registros.
  Est� tambi�n es una pseudo instrucci�n.
\end{itemize}

\hypertarget{pseudo-instrucciones}{%
\sectiontext{white}{black}{PSEUDO-INSTRUCCIONES}\label{pseudo-instrucciones}}

Los procesadores RISC tienen muy pocas instrucciones, como su propio
nombre indica y sucede a menudo que, para realizar operaciones
relativamente comunes es necesario utiliza dos instrucciones. En esos
casos se definen pseudo instrucciones, es decir, instrucciones que no
existen en el procesador, pero que se han especificado y los
ensambladores entiende. ARM, MIPS y tambi�n RISC-V definen varias de
estas pseudo instrucciones.

\begin{entradilla}
{\em Los procesadores RISC suelen ofrecer {\color{introcolor}pseudo-instruciones} para operaciones comunes que requiren m�s de una instrucci�n}
\end{entradilla}

Para nuestro ejemplo anterior, veremos como se implementan las dos
pseudo instrucciones de las que os hemos hablado. Vamos a comenzar con
{\verb!move!} que es m�s sencilla:

\begin{lstlisting}
   move t1, zero   <=> addi t1, zero, 0
   move a0, a1     <=> addi a0, a1, 0
\end{lstlisting}

La instrucci�n {\verb!addi!} \emph{Add Immediate} nos
permite sumar un valor num�rico a un registro y almacenarlo en otro. Si
le sumamos {\verb!0!}\ldots{} bueno, estamos copiando
el valor de un registro en otro\ldots{} V�is por donde vamos no?

La instrucci�n {\verb!li!} tambi�n se implementa usando
{\verb!addi!} en algunos casos:

\begin{lstlisting}
   li a1, 10  <=> addi a1, zero, 10
\end{lstlisting}

La instrucci�n {\verb!la!} es un poco m�s complicada.
Se expande de la siguiente forma:

\begin{lstlisting}
    la a1, buffer  <=> auipc a1, XX
                       addi  a1, a1, YY
\end{lstlisting}

La instrucci�n {\verb!auipc r, 0xAABB!} realiza la
siguiente operaci�n:

\begin{lstlisting}
r = 0xAABB0000 + PC
\end{lstlisting}

Es decir, toma el par�metro y lo carga en la parte alta del registro, y
luego le suma el valor del contador del programa. Es una forma de
direccionamiento relativo. El {\verb!addi!} que sigue
ajusta la parte baja de la direcci�n si es necesario. Afortunadamente no
tenemos que preocuparnos de calcular todas estas cosas, el ensamblador
lo har� por nosotros.

\begin{quote}
Pod�is estudiar vosotros que instrucciones son pseudo instrucciones y en
que se traducen utilizando el siguiente comando sobre vuestros binarios:
riscv64-linux-gnu-objdump -d -M no-aliases PROG
\end{quote}

\hypertarget{leyendo-nuestro-programa-brainfck-en-un-buffer}{%
\sectiontext{white}{black}{LEYENDO PROGRAMA BF EN UN BUFFER}\label{leyendo-nuestro-programa-brainfck-en-un-buffer}}

Con todo lo que acabamos de ver el c�digo para leer la entrada del
usuario (el programa BrainF*ck) en un buffer ser�a tal que as�:

\begin{lstlisting}
    # Lee datos desde entrada standard
    la   a1, prog  # Buffer de lectura
    move t1, zero  # Contador longitud del programa
bucle01:    
    li   a7, 63   # SYS_read
    li   a0, 0    # stdin
    li   a2, 1    # leemos caracter a caracter
    # a0 = read (a7=SYS_read, a0=stdin, a1=prog, a2=1);
    ecall       
    blez a0, cont01 # Si a0 <= 0 dejamos de leer
    
    add  t1, t1, 1 # Contamos un caracter
    add  a1, a1, 1 # Actualizamos el pumtero
    j    bucle01     # Repetimos
cont01:
.data
prog:
    .fill MEM_SIZE  
\end{lstlisting}

Deber�ais de identificar f�cilmente la llamada al sistema
{\verb!SYS\_read!}. Tambi�n pod�is ver como declaramos
el buffer en la secci�n {\verb!.data!}. Esto le dice al
linker que, lo que pongamos a partir de ah� se cargar� en un bloque de
memoria con permisos de lectura y escritura, pero no de ejecuci�n.

Veamos como funciona el bucle de lectura. Una vez que ejecutamos la
llamada al sistema con {\verb!ecall!} en el registro
{\verb!a0!} tendremos el resultado de la operaci�n.
Para {\verb!read!} un valor negativo indica un error,
un valor 0 indica fin de fichero ({\verb!EOF!}) y un
valor {\verb!1!} indica que hemos le�do nuestro
caracter correctamente.

As� que usamos la instrucci�n {\verb!blez a0, cont01!}.
Esta instrucci�n salta a {\verb!cont01!} si el valor de
{\verb!a0!} es menor o igual que cero. Y si. Esta es
tambi�n una pseudo instrucci�n. RISC-V solo define un par de
instrucciones de salto condicional y el resto de posibles saltos se
generan como pseudo instrucciones. En este caso:

\begin{lstlisting}
blez r, offset   <=> bge zero, r, offset
\end{lstlisting}

Que significa \_Salta si {\verb!0!} es mayor o igual el
registro {\verb!r!}. Que es lo mismo que salta si
{\verb!r!} es menor o igual que cero.

Esa es nuestra condici�n de final de buffer. En caso de tener que hacer
otra iteraci�n simplemente incrementamos el puntero dentro de nuestro
buffer y una variable que usamos para contar los caracteres le�dos y
saber el tama�o de nuestro programa.

\hypertarget{comprobando-que-vamos-bien}{%
\sectiontext{white}{black}{COMPROBANDO PROGRESO}\label{comprobando-que-vamos-bien}}

Para comprobar que todo ha ido bien, pod�is a�adir una llamada a
{\verb!SYS\_write!} para imprimir el buffer.

\begin{lstlisting}
    # Muestra la cadena leida por pantalla
    li   a7, 64     # SYS_write
    li   a0, 1      # a0 = stdout
    la   a1, prog   # a1 = prog
    move a2, t1   # a2 = size
    # write (a7=SYS_write, a0=stdout, a1=prog, a2=size);
    ecall         
\end{lstlisting}

Finalmente terminamos el programa con una llamada a
{\verb!SYS\_EXIT!} para que todo termine de forma
ordenada.

\begin{lstlisting}
exit:
    li a7, 93   # Exit  
    li a0, 0
    ecall
prog:
    .fill MEM_SIZE  
\end{lstlisting}

El c�digo que hemos escrito hasta el momento deber�a funcionar sin
problemas. Leer� todo lo que escribamos en la consola hasta que pulsemos
{\verb!CTRL+D!} (esta es la forma de mandar un
{\verb!EOF!} desde la consola) y luego lo imprimir� en
pantalla.

\hypertarget{ejecutando-nuestro-primer-programa}{%
\sectiontext{white}{black}{EJECUTANDO NUESTRO PRIMER PROGRAMA}\label{ejecutando-nuestro-primer-programa}}

La forma m�s sencilla de que pod�is instalar todas las herramientas que
necesit�is sin problemas de configuraci�n especial de vuestra m�quina es
utilizando un contenedor docker. Eso s�, ten�is que tener docker
configurado correctamente y eso no lo vamos a tratar en este art�culo.
As� que vamos a ello:

\begin{entradilla}
{\em La forma m�s sencilla de compilar y ejecutar nuestro programa es {\color{introcolor}desplegando un contendor docker}}
\end{entradilla}

Supongamos que tenemos nuestro programa en ensamblador en un
subdirectorio llamado {\verb!bf!} najo nuestro
directorio actual. As� que iniciamos un contenedor Debian montando ese
directorio:

{\scriptsize
\begin{verbatim}
$ docker run -it --rm -v $PWD/bf:/opt/src debian:12
root@5bca9d382ba1:/# apt update
(...)
root@5bca9d382ba1:/# apt upgrade
(...)
root@5bca9d382ba1:/# apt upgrade
root@5bca9d382ba1:/# apt install gcc-riscv64-linux-gnu qemu-user-static
root@5bca9d382ba1:/# cd /opt/src/
root@5bca9d382ba1:/# riscv64-linux-gnu-as -o bf-riscv.o bf-riscv.s
root@5bca9d382ba1:/tmp# riscv64-linux-gnu-ld -o bf-riscv bf-riscv.o 
root@5bca9d382ba1:/tmp# qemu-riscv64-static ./bf-riscv
\end{verbatim}
}

O si ten�is alguna m�quina RISC-V como mi Milk-V Duo, pod�is subir el
binario y ejecutarlo all�.


\sectiontext{white}{black}{CONCLUSIONES}\label{conclusiones}

En este art�culo hemos aprendido a programar en BrainFuck y unas cuantas cosas de como programar ensamblador para RISC-V, Hemos programa el esqueleto de nuestro int�rprete y el c�digo para leer programas desde la consola y procesarlos. En el pr�ximo n�mero a�adiremos los comandos b�sicos para poder empezar a ejecutar nuestros propios programas.

\end{multicols}


\rput(7.9,-3){\resizebox{!}{6.0cm}{{\epsfbox{images/promo/promo04.eps}}}}

\rput(2.5,-1.5){\resizebox{!}{4.5cm}{{\epsfbox{images/general/trucos1.eps}}}}
\hypertarget{trucos1}{}\label{trucos}

\pagestyle{trucos}


% -------------------------------------------------
% Cabecera
\begin{flushright}


{\color{introcolor}\mtitle{8cm}{Con un par... de l�neas}}

\msubtitle{5cm}{Los Mejores Trucos}

{\sf por Tamariz el de la Perdiz}

{\psset{linecolor=black,linestyle=dotted}\psline(-12,0)}
\end{flushright}

\vspace{2mm}
% -------------------------------------------------

\begin{multicols}{2}

%\lstset{language=C,frame=tb,framesep=5pt,basicstyle=\footnotesize}
\lstset{language=C,frame=tb,framesep=5pt,basicstyle=\scriptsize}



\hypertarget{documentos-here-y-luxedneas-here}{%
\sectiontext{white}{black}{DOCUMENTOS HERE Y L�NEAS HERE}\label{documentos-here-y-luxedneas-here}}

Probablemente conozcas los documentos \texttt{HERE}. Una forma comoda de
hacer que nuestros scripts escriban varias lineas en pantalla de una
sola vez.

\begin{verbatim}
echo << EOM
Esto son varias lineas que podemos imprimir
Todo lo que escribamos hasta la marca
aparecera en la consola
EOM
\end{verbatim}

Tambien los podemos utilizar para generar ficheros cuando no tenemos un
editor de textos

\begin{verbatim}
cat << EOM  > hola.c
> #include <stdio.h>
> int main(void) { puts("Hola Mundo");}
> EOM
\end{verbatim}

Pero en ocasiones solamente necesitamos una l�nea. Los siguientes
comandos son equivalentes en bash

\begin{verbatim}
$ base64 <<< "Codifica esto en base64"
$ echo "Codifica esto en base64" | base64
\end{verbatim}

Finalmente, la sustituci�n de procesos nos permite utilizar la entrada o
la salida de un comando como entrada o salida de otro. La sustituci�n de
procesos crea un \texttt{pipe} con el que redirigir la entrada o salida
de procesos. Veamos un par de ejemplos

Un ejemplo cl�sico del uso de este operador es la comparaci�n de la
salida de dos procesos. En general, para poder hacer eso, deberiamos
redireccionar la salida de cada proceso a un fichero y luego comparar
esos ficheros. Pero en su lugar podemos hacer lo siguiente.

\begin{verbatim}
$ diff <(cmd1) <(cmd2)
\end{verbatim}

Observad que no hay espacio entre \texttt{\textless{}} y \texttt{(}. De
la misma forma podemos redireccionar la salida de un comando

{\footnotesize
\begin{verbatim}
$ echo "Hola mundo" | tee >(md5sum) >(sha256sum)
\end{verbatim}
}

En este caso el problema esta en que \texttt{tee} espera un fichero como
par�metro, no un comando y por lo tanto no podriamos enciar nuestra
cadena de texto a los dos programas que calculan el hash.

\hypertarget{lista-errores-del-sistema}{%
\sectiontext{white}{black}{LISTA DE ERRORES DEL SISTEMA}\label{lista-errores-del-sistema}}

Durante el desarrollo de programas en C es normal obtener errores como
valores numericos en la variable global \texttt{errno}. Podemos utilizar
la funci�n \texttt{perror} para traducirlos en una cadena comprensible
para nosotros, o podemos utilizar el comando \texttt{errno} para
traducirlos\ldots{} en caso de que, por ejemplo, el programador de
alguna utilidad simplemente nos devuelva un cr�ptico \texttt{error\ 10}

{\small
\begin{verbatim}
$ errno -l  # Muestra todos los errores del sistema
$ errno 10  # Muestra descripcion del error 10
ECHILD 10 No child processes
$ errno EAGAIN # Muestra descripcion del error EAGAIN
EAGAIN 11 Resource temporarily unavailable
\end{verbatim}
}

Si el programa no est� instalado, podr�is encontrarlo en el paquete
\texttt{moreutils} al menos en la distribucion Debian.

\hypertarget{ejecutar-comandos-desde-tu-editor-de-textos-sobre-documento}{%
\sectiontext{white}{black}{EJECUTAR COMANDOS DIRECTAMENTE SOBRE TU TEXTO}\label{ejecutar-comandos-desde-tu-editor-de-textos-sobre-documento}}

Sabias que puedes ejecutar comandos shell sobre bloques de texto que
estas editando en vim. Por ejemplo, puedes ordenar alfabeticamente una
lista que hayas escrito desordenada en tu documento utilizando el
comando:

\begin{verbatim}
:%!cmd
\end{verbatim}

El \texttt{\%} significa aplicar el comando a todo el fichero, pero
tambien podeis usarlo en modo visual.

EN emacs podies hacer lo mismo utilizando el siguiente comando.

\begin{verbatim}
CTRL + U  ALT + | cmd
\end{verbatim}

\hypertarget{servidor-ftp-temporal}{%
\sectiontext{white}{black}{SERVIDOR FTP TEMPORAL}\label{servidor-ftp-temporal}}

Hay muchas formas de transferir ficheros r�pidamente entre ordenadores
de una red. Aqu� os traemos una opci�n m�s, utilizando el modulo
\texttt{pyftplib}. Poner en marcha un servidor web es tan sencillo como
esto.

\begin{verbatim}
$ python -m pyftpdlib --directory=FTP \\
         --port=2121 --write
\end{verbatim}

Para acceder al servidor utilizad el nombre de usuario \texttt{anon}

Pod�is instalar \texttt{pyftplib} con el comando.

\begin{verbatim}
$ pip install --user pyftpdlib
\end{verbatim}

\end{multicols}




%\include{trucos}


\pagestyle{empty}

\rput(8,-8.0){\resizebox{21.5cm}{!}{{\epsfbox{images/general/contraportada.eps}}}}



\end{document}
% Convert files
%  iconv -c -t ISO_8859-1 paper3.tex -o paper4.tex
% -c elimina caracteres invalidos
%  file -ib paper1.tex

% Emacs substitutions
%  \\passthrough{\\lstinline!\(.+\)!} -> \\verb!\1!
% 

% Colors
% black, blue, brown, cyan, darkgray, gray, green, lightgray, lime, magenta, olive, orange, pink, purple, red, teal, violet, white, yellow.
